% Generated mechanically from the frozen obsolete.tex target.
% Do not edit this generated file; edit the bound locale source instead.

% OK, start here.
%

\setcounter{chapter}{114}
\chapter{已废弃内容}


\phantomsection
\label{obsolete-section-phantom}





\section{引言}
\label{obsolete-section-introduction}

\noindent
本章收录一些已经变得“已废弃”的引理（见 \cite{Miller}）。


\section{预备知识}
\label{obsolete-section-preliminaries}

\begin{remark}
\label{obsolete-remark-composition-of-adjoints-isomorphic-to-identity}
本注记原先包含的信息，现在已由下列两个引理的结合涵盖：
Categories，引理 \ref{categories-lemma-adjoint-fully-faithful} 与
\ref{categories-lemma-left-adjoint-composed-fully-faithful}。
\end{remark}


\section{同调代数}
\label{obsolete-section-homological-algebra}

\begin{remark}
\label{obsolete-remark-weak-serre-subcategory}
下列注记已经废弃，因为它们已由
Homology，引理 \ref{homology-lemma-biregular-ss-converges} 与
\ref{homology-lemma-first-quadrant-ss} 涵盖。
设 $\mathcal{A}$ 是阿贝尔范畴，且
$\mathcal{C} \subset \mathcal{A}$ 是弱 Serre 子范畴（见
Homology，定义 \ref{homology-definition-serre-subcategory}）。
设 $K^{\bullet, \bullet}$ 是适用
Homology，引理 \ref{homology-lemma-first-quadrant-ss} 的双复形，并且对某个
$r \geq 0$，所有对象 ${}'E_r^{p, q}$ 都属于 $\mathcal{C}$。
则所有上同调群 $H^n(sK^\bullet)$ 都属于 $\mathcal{C}$。事实上，这些假设意味着
${}'d_r^{p, q}$ 的核与像都在 $\mathcal{C}$ 中，于是每个
${}'E_{r + 1}^{p, q}$ 也在 $\mathcal{C}$ 中。归纳可知每个
${}'E_\infty^{p, q}$ 都在 $\mathcal{C}$ 中。因此每个
$H^n(sK^\bullet)$ 都有一个有限过滤，其子商属于 $\mathcal{C}$。
利用 $\mathcal{C}$ 对扩张封闭，可得 $H^n(sK^\bullet)$ 属于
$\mathcal{C}$，如所述。
对 $K^{\bullet, \bullet}$ 关联的第二个谱序列也有同样结论。类似地，若
$(K^\bullet, F)$ 是适用
Homology，引理 \ref{homology-lemma-biregular-ss-converges} 的过滤复形，并且对某个
$r \geq 0$，所有对象 $E_r^{p, q}$ 都属于 $\mathcal{C}$，则每个
$H^n(K^\bullet)$ 都是 $\mathcal{C}$ 中的对象。
\end{remark}




\section{已废弃的代数引理}
\label{obsolete-section-algebra}

\begin{lemma}
\label{obsolete-lemma-finite-presentation-module-independent}
设 $M$ 是有限表示的 $R$-模。对任意满射
$\alpha : R^{\oplus n} \to M$，$\alpha$ 的核是有限生成的 $R$-模。
\end{lemma}

\begin{proof}
这是 Algebra，引理 \ref{algebra-lemma-extension} 的特殊情形。
\end{proof}

\begin{lemma}
\label{obsolete-lemma-p-ring-map}
设 $\varphi : R \to S$ 是环映射。若
\begin{enumerate}
\item 对任意 $x \in S$，存在 $n > 0$ 使得
$x^n$ 属于 $\varphi$ 的像；并且
\item 对任意 $x \in \Ker(\varphi)$，存在 $n > 0$
使得 $x^n = 0$；
\end{enumerate}
则 $\varphi$ 在谱空间上诱导同胚。给定素数 $p$，若
\begin{enumerate}
\item[(a)] $S$ 作为 $R$-代数由元素 $x$ 生成，并且对每个这样的 x，
存在 $n > 0$ 使 $x^{p^n} \in \varphi(R)$ 且
$p^nx \in \varphi(R)$；并且
\item[(b)] $\varphi$ 的核由幂零元素生成，
\end{enumerate}
则 (1) 与 (2) 成立；对任意环映射 $R \to R'$，
环映射 $R' \to R' \otimes_R S$ 也满足 (a)、(b)、(1) 和 (2)，
尤其在谱空间上诱导同胚。
\end{lemma}

\begin{proof}
这是 Algebra，引理 \ref{algebra-lemma-powers} 与
\ref{algebra-lemma-p-ring-map} 的组合。
\end{proof}

\noindent
下面的技术性引理说明：对于在适当意义下本质有限型的环映射，
可以把纤维中的任意关系序列提升到整个空间。

\begin{lemma}
\label{obsolete-lemma-lift-elements-ideal}
设 $R \to S$ 是环映射。
设 $\mathfrak p \subset R$ 是素理想，
$\mathfrak q \subset S$ 是位于 $\mathfrak p$ 上方的素理想。
假设 $S_{\mathfrak q}$ 在 $R_\mathfrak p$ 上本质有限型。
给定
\begin{enumerate}
\item 整数 $n \geq 0$；
\item 素理想 $\mathfrak a \subset \kappa(\mathfrak p)[x_1, \ldots, x_n]$；
\item 满的 $\kappa(\mathfrak p)$-同态
$$
\psi : (\kappa(\mathfrak p)[x_1, \ldots, x_n])_{\mathfrak a}
\longrightarrow
S_{\mathfrak q}/\mathfrak p S_{\mathfrak q},
$$
以及
\item $\Ker(\psi)$ 中的元素 $\overline{f}_1, \ldots, \overline{f}_e$。
\end{enumerate}
则存在
\begin{enumerate}
\item 整数 $m \geq 0$；
\item 元素 $g \in S$，$g \not\in \mathfrak q$；
\item 映射
$$
\Psi :
R[x_1, \ldots, x_n, x_{n + 1}, \ldots, x_{n + m}]
\longrightarrow
S_g,
$$
以及
\item $\Ker(\Psi)$ 中的元素
$f_1, \ldots, f_e, f_{e + 1}, \ldots, f_{e + m}$
\end{enumerate}
使得
\begin{enumerate}
\item 下图交换：
$$
\xymatrix{
R[x_1, \ldots, x_{n + m}] \ar[d]_\Psi
\ar[rr]_-{x_{n + j} \mapsto 0} & &
(\kappa(\mathfrak p)[x_1, \ldots, x_n])_{\mathfrak a} \ar[d]^\psi \\
S_g \ar[rr] & &
S_{\mathfrak q}/\mathfrak p S_{\mathfrak q}
},
$$
\item 元素 $f_i$（$i \leq n$）在局部环中映到
$\overline{f}_i$ 的某个单位倍；
$$
(\kappa(\mathfrak p)[x_1, \ldots, x_{n + m}])_{
(\mathfrak a, x_{n + 1}, \ldots, x_{n + m})},
$$
\item 元素 $f_{e + j}$ 在同一局部环中映到
$x_{n + j}$ 的某个单位倍；并且
\item 诱导映射 $R[x_1, \ldots, x_{n + m}]_{\mathfrak b}
\to S_{\mathfrak q}$ 是满射，其中
$\mathfrak b = \Psi^{-1}(\mathfrak qS_g)$。
\end{enumerate}
\end{lemma}

\begin{proof}
我们断言，只需在 $R$ 与 $S$ 都是局部环，
且极大理想分别为 $\mathfrak p$ 与 $\mathfrak q$ 的情形证明引理。
事实上，假设已经构造出
$$
\Psi' : R_{\mathfrak p}[x_1, \ldots, x_{n + m}]
\longrightarrow
S_{\mathfrak q}
$$
以及 $f_1', \ldots, f_{e + m}' \in R_{\mathfrak p}[x_1, \ldots, x_{n + m}]$，
并满足所有要求。则存在元素
$f \in R$，$f \not \in \mathfrak p$，使每个
$ff_k'$ 都来自某个 $f_k \in R[x_1, \ldots, x_{n + m}]$。
此外，对某个适当的 $g \in S$，$g \not \in \mathfrak q$，
元素 $\Psi'(x_i)$ 是某些元素 $y_i \in S_g$ 的像。
令 $\Psi$ 为由规则 $\Psi(x_i) = y_i$ 定义的 $R$-代数映射。
由于 $\Psi(f_i)$ 在局部化 $S_{\mathfrak q}$ 中为零，
必要时替换 $g$ 后可假设 $\Psi(f_i) = 0$。断言得证。

\medskip\noindent
因此可假设 $R$ 与 $S$ 都是局部环，
其极大理想分别为 $\mathfrak p$ 与 $\mathfrak q$。
取 $y_1, \ldots, y_n \in S$，使得
$y_i \bmod \mathfrak pS = \psi(x_i)$。
取元素 $y_{n + 1}, \ldots, y_{n + m} \in S$，使它们生成某个
$R$-子代数，而 $S$ 是该子代数的局部化。
这样的元素存在，因为假设 $S$ 在 $R$ 上本质有限型。
由于 $\psi$ 是满射，可对某个
$h_j \in \kappa(\mathfrak p)[x_1, \ldots, x_n]_{\mathfrak a}$ 写成
$y_{n + j} \bmod \mathfrak pS = \psi(h_j)$。
写 $h_j = g_j/d$，其中 $g_j \in \kappa(\mathfrak p)[x_1, \ldots, x_n]$，
公共分母 $d \in \kappa(\mathfrak p)[x_1, \ldots, x_n]$ 满足
$d \not \in \mathfrak a$。取 $g_j$ 与 $d$ 的提升
$G_j, D \in R[x_1, \ldots, x_n]$。置
$y_{n + j}' = D(y_1, \ldots, y_n) y_{n + j} - G_j(y_1, \ldots, y_n)$。
由构造，$y_{n + j}' \in \mathfrak p S$。
显然，$y_1, \ldots, y_n, y_n', \ldots, y_{n + m}'$
生成 $S$ 的一个 $R$-子代数，其局部化为 $S$。
定义
$$
\Psi : R[x_1, \ldots, x_{n + m}] \to S
$$
为把 $x_i$（$i = 1, \ldots, n$）映到 $y_i$，
并把 $x_{n + j}$（$j = 1, \ldots, m$）映到 $y'_{n + j}$ 的映射。
由构造，性质 (1) 与 (4) 显然成立。此外，理想
$\mathfrak b$ 满射到多项式环 $\kappa(\mathfrak p)[x_1, \ldots, x_{n + m}]$ 中的理想
$(\mathfrak a, x_{n + 1}, \ldots, x_{n + m})$。

\medskip\noindent
记 $J = \Ker(\Psi)$。有短正合列
$$
0 \to J_{\mathfrak b}
\to R[x_1, \ldots, x_{n + m}]_{\mathfrak b}
\to S_{\mathfrak q}
\to 0.
$$
满射性来自上面对
$y_1, \ldots, y_n, y_n', \ldots, y_{n + m}'$ 的选择。
这意味着
$$
J_{\mathfrak b}/ \mathfrak pJ_{\mathfrak b}
\to \kappa(\mathfrak p)[x_1, \ldots, x_{n + m}]_{
(\mathfrak a, x_{n + 1}, \ldots, x_{n + m})}
\to S_{\mathfrak q}/\mathfrak pS_{\mathfrak q}
\to 0
$$
正合。由构造，在最后一个映射下 $x_i$ 映到 $\psi(x_i)$，
$x_{n + j}$ 映到零。
因此容易按引理的 (2) 与 (3) 选择 $f_i$。
\end{proof}

\begin{lemma}
\label{obsolete-lemma-smooth-independent-presentation}
设 $R \to S$ 是有限表示的环映射。
若对 $S$ 在 $R$ 上的某个表示 $\alpha$，其
朴素余切复形 $\NL(\alpha)$ 拟同构于置于次数 $0$ 的
有限投射 $S$-模，则
对任意表示都成立。
\end{lemma}

\begin{proof}
立即由 Algebra，引理 \ref{algebra-lemma-NL-homotopy} 得出。
\end{proof}

\begin{remark}[投射分解]
\label{obsolete-remark-projective-resolution}
设 $R$ 是环。
对任意集合 $S$，记 $F(S)$ 为以 $S$ 为基的自由 $R$-模。
则任意左 $R$-模都有如下两步分解
$$
F(M \times M) \oplus F(R \times M) \to F(M) \to M \to 0.
$$
第一个映射由规则
$$
[m_1, m_2] \oplus [r, m] \mapsto [m_1 + m_2] - [m_1] - [m_2] + [rm] - r[m].
$$
\end{remark}

\begin{lemma}
\label{obsolete-lemma-spec-localization-first}
设 $S$ 是 $A$ 的乘法闭集。则映射
$$
f: \Spec(S^{-1}A)\longrightarrow \Spec(A)
$$
由典范环映射
$A \to S^{-1}A$ 诱导的映射到其像上是同胚，且
$\Im(f) = \{ \mathfrak p \in \Spec(A) : \mathfrak p\cap S = \emptyset \}$。
\end{lemma}

\begin{proof}
这是 Algebra，引理 \ref{algebra-lemma-spec-localization} 的重复。
\end{proof}

\begin{lemma}
\label{obsolete-lemma-finite-type-flat-over-integral-algebra}
设 $A \to B$ 是有限型平坦环映射，且 $A$ 是整
环。则 $B$ 是有限表示的 $A$-代数。
\end{lemma}

\begin{proof}
这是 More on Flatness，命题
\ref{flat-proposition-flat-finite-type-finite-presentation-domain} 的特殊情形。
\end{proof}

\begin{lemma}
\label{obsolete-lemma-helper-finite-type-flat-finite-presentation}
设 $R$ 是分式域为 $K$ 的整环。
设 $S = R[x_1, \ldots, x_n]$ 是 $R$ 上的多项式环。
设 $M$ 是有限 $S$-模，并假设 $M$ 平坦于 $R$。
若对每个子环 $R \subset R' \subset K$ 且 $R \not = R'$，
$M \otimes_R R'$ 在
$S \otimes_R R'$ 上有限表示，则 $M$ 在 $S$ 上有限表示。
\end{lemma}

\begin{proof}
该引理成立，是因为即使不作引理陈述中对每个 $R'$ 的
$M \otimes_R R'$ 有限表示这一假设，$M$ 仍是有限表示的。
这由 More on Flatness，命题
\ref{flat-proposition-flat-finite-type-finite-presentation-domain} 得出。
这个引理原先有一个错误的证明（感谢 Ofer Gabber
发现其中的缺口），并曾用于所引命题的另一种证明。要恢复这个引理，
在 $R$ 为局部正规整环时，我们需要只使用交换代数各章结果的
正确论证；若有论证，请发送邮件至
\href{mailto:stacks.project@gmail.com}{stacks.project@gmail.com}。
\end{proof}

\begin{lemma}
\label{obsolete-lemma-relative-effective-cartier-algebra}
设 $A \to B$ 是环映射，且 $f \in B$。假设
\begin{enumerate}
\item $A \to B$ 平坦；
\item $f$ 是非零因子；并且
\item $A \to B/fB$ 平坦。
\end{enumerate}
于是对任意理想 $I \subset A$，映射
$f : B/IB \to B/IB$ 是单射。
\end{lemma}

\begin{proof}
注意，$IB = I \otimes_A B$ 以及 $I(B/fB) = I \otimes_A B/fB$。
这是因为 $B$ 与 $B/fB$ 都平坦于 $A$。
特别地，$IB/fIB \cong I \otimes_A B/fB$ 单射地映入
$B/fB$。因此将蛇形引理应用于下列图表即可得到结论。
$$
\xymatrix{
0 \ar[r] &
I \otimes_A B \ar[r] \ar[d]^f &
B \ar[r] \ar[d]^f &
B/IB \ar[r] \ar[d]^f &
0 \\
0 \ar[r] &
I \otimes_A B \ar[r] &
B \ar[r] &
B/IB \ar[r] &
0
}
$$
其中各行都是正合的。
\end{proof}

\begin{lemma}
\label{obsolete-lemma-faithfully-flat-injective}
若 $R \to S$ 是忠实平坦的环映射，则对每个 $R$-模
$M$，映射 $M \to S \otimes_R M$，$x \mapsto 1 \otimes x$ 是单射。
\end{lemma}

\begin{proof}
这是 Algebra，引理
\ref{algebra-lemma-faithfully-flat-universally-injective} 的重复。
\end{proof}

\begin{remark}
\label{obsolete-remark-section-colimits}
这个参考/标签原先指向 Smoothing Ring Maps 一章中的某个小节，
但相关材料后来已并入 Algebra 的小节
\ref{algebra-section-colimits-flat}。
\end{remark}

\begin{lemma}
\label{obsolete-lemma-not-domain}
设 $(R, \mathfrak m)$ 是维数为 $1$ 的约化诺特局部环，
且 $x \in \mathfrak m$ 是非零因子。设
$\mathfrak q_1, \ldots, \mathfrak q_r$ 是 $R$ 的极小素理想。则
$$
\text{length}_R(R/(x)) = \sum\nolimits_i \text{ord}_{R/\mathfrak q_i}(x)
$$
\end{lemma}

\begin{proof}
这是 Chow Homology，引理
\ref{chow-lemma-additivity-divisors-restricted} 的特殊（非常容易）情形。
\end{proof}

\begin{lemma}
\label{obsolete-lemma-bound-primes}
设 $A$ 是维数为 $2$ 的诺特局部正规整环。
对非零的 $f \in \mathfrak m$，记
$\text{div}(f) = \sum n_i (\mathfrak p_i)$
为 $f$ 在 $A$ 的去心谱上的因子。
置 $|f| = \sum n_i$。存在整数 $N$ 与 $M$，
使得对所有 $g \in \mathfrak m^N$ 都有 $|f + g| \leq M$。
\end{lemma}

\begin{proof}
取 $h \in \mathfrak m$ 使得 $f, h$ 是 $A$ 中的正则序列
（这由 Algebra，引理 \ref{algebra-lemma-criterion-normal} 与
\ref{algebra-lemma-depth-drops-by-one} 推出）。
取 $M = \text{length}_A(A/(f, h))$，并取满足
$\mathfrak m^N \subset (f, h)$ 的任意整数 $N$ 来证明引理。这样的
$N$ 存在，因为 $\sqrt{(f, h)} = \mathfrak m$。注意，对所有 $g \in \mathfrak m^N$，
都有 $M = \text{length}_A(A/(f + g, h))$，因为
$(f, h) = (f + g, h)$。这也说明 $f + g, h$
同样是 $A$ 中的正则序列，见
Algebra，引理 \ref{algebra-lemma-reformulate-CM}。
现在设 $\text{div}(f + g ) = \sum m_j (\mathfrak q_j)$。
考虑映射
$$
c : A/(f + g) \longrightarrow \prod A/\mathfrak q_j^{(m_j)}
$$
其中 $\mathfrak q_j^{(m_j)}$ 是符号幂，见
Algebra，小节 \ref{algebra-section-symbolic-power}。
由于 $A$ 正规，可知 $A_{\mathfrak q_i}$ 是
离散赋值环，因此
$$
A_{\mathfrak q_i}/(f + g) =
A_{\mathfrak q_i}/\mathfrak q_i^{m_i} A_{\mathfrak q_i} =
(A/\mathfrak q_i^{(m_i)})_{\mathfrak q_i}
$$
由于 $V(f + g, h) = \{\mathfrak m\}$，这意味着 $c$ 在令 $h$ 可逆后
成为同构（略去一个细节）。由于 $h$ 是 $A/(f + g)$ 上的非零因子，
可知 $A/(f + g, h)$ 的长度等于 Chow Homology 小节
\ref{chow-section-periodic-complexes} 中定义的 Herbrand 商
$e_A(A/(f + g), 0, h)$。
类似地，$A/(h, \mathfrak q_j^{(m_j)})$ 的长度等于
$e_A(A/\mathfrak q_j^{(m_j)}, 0, h)$。于是
\begin{align*}
M & = \text{length}_A(A/(f + g, h) \\
& =
e_A(A/(f + g), 0, h) \\
& =
\sum\nolimits_i e_A(A/\mathfrak q_j^{(m_j)}, 0, h) \\
& =
\sum\nolimits_i \sum\nolimits_{m = 0, \ldots, m_j - 1}
e_A(\mathfrak q_j^{(m)}/\mathfrak q_j^{(m + 1)}, 0, h)
\end{align*}
这些等式由 Chow Homology，引理
\ref{chow-lemma-additivity-periodic-length} 与
\ref{chow-lemma-finite-periodic-length} 推出；特别使用了
上面所述 $c$ 的余核具有有限长度这一事实。
容易证明，
$e_A(\mathfrak q^{(m)}/\mathfrak q^{(m + 1)}, 0, h)$
由 Nakayama 引理可知至少为 $1$。引理得证。
\end{proof}

\begin{lemma}
\label{obsolete-lemma-flat-over-gorenstein-gorenstein-fibre}
设 $A \to B$ 是诺特局部环之间的平坦局部同态。
若 $A$ 与 $B/\mathfrak m_A B$ 都是 Gorenstein 环，则 $B$ 是 Gorenstein 环。
\end{lemma}

\begin{proof}
立即由 Dualizing Complexes，引理
\ref{dualizing-lemma-flat-under-gorenstein} 得出。
\end{proof}

\begin{lemma}
\label{obsolete-lemma-kill-local}
设 $(A, \mathfrak m)$ 是诺特局部环。
设 $I \subset A$ 是理想，$M$ 是有限 $A$-模。
设 $s$ 是整数。假设
\begin{enumerate}
\item $A$ 有对偶化复形；
\item 若 $\mathfrak p \not \in V(I)$ 且
$V(\mathfrak p) \cap V(I) \not = \{\mathfrak m\}$，则
$\text{depth}_{A_\mathfrak p}(M_\mathfrak p) + \dim(A/\mathfrak p) > s$。
\end{enumerate}
则存在 $n > 0$ 与理想 $J \subset A$，满足
$V(J) \cap V(I) = \{\mathfrak m\}$，且 $JI^n$ 湮灭
$i \leq s$ 时的 $H^i_\mathfrak m(M)$。
\end{lemma}

\begin{proof}
根据
Local Cohomology，引理 \ref{local-cohomology-lemma-sitting-in-degrees}，
只需对有限 $A$-模
$E^i = \text{Ext}^{-i}_A(M, \omega_A^\bullet)$
（$i \leq s$）证明这一点。$E^0 \oplus \ldots \oplus E^s$ 的支集 $Z$
在 $\Spec(A)$ 中闭合，且不包含满足 (2) 的任何素理想。
因此对某个符合引理陈述的 $J$，它包含于
$V(JI^n)$ 中。
\end{proof}

\begin{lemma}
\label{obsolete-lemma-algebraize-local-cohomology-bis-bis}
设 $(A, \mathfrak m)$ 是诺特局部环。
设 $I \subset A$ 是理想，$M$ 是有限 $A$-模。
设 $s$ 与 $d$ 是整数。假设
\begin{enumerate}
\item[(a)] $A$ 有对偶化复形；
\item[(b)] $\text{cd}(A, I) \leq d$；
\item[(c)] 若 $\mathfrak p \not \in V(I)$，则
$\text{depth}_{A_\mathfrak p}(M_\mathfrak p) > s$，或
$\text{depth}_{A_\mathfrak p}(M_\mathfrak p) + \dim(A/\mathfrak p) > d + s$。
\end{enumerate}
则 Algebraic and Formal Geometry，引理
\ref{algebraization-lemma-algebraize-local-cohomology-bis} 对
$A, I, \mathfrak m, M$ 成立，并且
$H^i_\mathfrak m(M) \to \lim H^i_\mathfrak m(M/I^nM)$
在 $i \leq s$ 时是同构；这些模
被 $I$ 的某个幂湮灭。
\end{lemma}

\begin{proof}
假设由更一般的 Algebraic and Formal Geometry，引理
\ref{algebraization-lemma-algebraize-local-cohomology-bis}，以及
Algebraic and Formal Geometry，引理
\ref{algebraization-lemma-bootstrap-bis-bis} 给出。
然后 Algebraic and Formal Geometry，引理
\ref{algebraization-lemma-algebraize-local-cohomology-bis} 的结论
给出第二个断言。
\end{proof}

\begin{lemma}
\label{obsolete-lemma-combine-one}
在 Algebraic and Formal Geometry，情形
\ref{algebraization-situation-bootstrap} 中，
有 $H^s_\mathfrak a(M) = \lim H^s_\mathfrak a(M/I^nM)$。
\end{lemma}

\begin{proof}
这立即由 Algebraic and Formal Geometry，定理
\ref{algebraization-theorem-final-bootstrap} 得出。
这个引理的原版本带有额外假设，
已被该定理取代。
\end{proof}

\begin{lemma}
\label{obsolete-lemma-fully-faithful-simple-one}
设 $A$ 是诺特环，且 $f \in \mathfrak a$ 属于 $A$ 的某个理想。
设 $U = \Spec(A) \setminus V(\mathfrak a)$。假设
\begin{enumerate}
\item $A$ 有对偶化复形，且关于 $f$ 完备；
\item $A_f$ 满足 $(S_2)$，并且对每个极小素理想 $\mathfrak p \subset A$，
$f \not \in \mathfrak p$，且对
$\mathfrak q \in V(\mathfrak p) \cap V(\mathfrak a)$ 有
$\dim((A/\mathfrak p)_\mathfrak q) \geq 3$。
\end{enumerate}
则完备化函子
$$
\textit{Coh}(\mathcal{O}_U)
\longrightarrow
\textit{Coh}(U, I\mathcal{O}_U),
\quad
\mathcal{F} \longmapsto \mathcal{F}^\wedge
$$
在有限局部自由对象构成的满子范畴上是全忠实的。
\end{lemma}

\begin{proof}
这是 Algebraic and Formal Geometry，引理
\ref{algebraization-lemma-fully-faithful-simple-two} 的特殊情形。
\end{proof}





\section{与 ZMT 相关的引理}
\label{obsolete-section-ZMT}

\noindent
本节的引理原先用于证明（代数版本的）Zariski 主定理，
即 Algebra，定理 \ref{algebra-theorem-main-theorem}。

\begin{lemma}
\label{obsolete-lemma-change-equation-multiply}
设 $R$ 是环，且 $\varphi : R[x] \to S$ 是
环映射。设 $t \in S$。若 $t$ 在
$R[x]$ 上整，则存在 $\ell \geq 0$，使得对每个 $a \in R$，元素
$\varphi(a)^\ell t$ 关于由下式定义的
$\varphi_a : R[y] \to S$ 是整的：
$y \mapsto \varphi(ax)$ 且 $r \mapsto \varphi(r)$，
其中 $r\in R$。
\end{lemma}

\begin{proof}
写成 $t^d + \sum_{i < d} \varphi(f_i)t^i = 0$，其中
$f_i \in R[x]$。令 $\ell$ 为所有
$f_i$ 关于 $x$ 的最大次数。将等式乘以 $\varphi(a)^\ell$，得到
$\varphi(a)^\ell t^d + \sum_{i < d} \varphi(a^\ell f_i)t^i = 0$。
注意，每个 $\varphi(a^\ell f_i)$ 都在
$\varphi_a$ 的像中。结论由
Algebra，引理 \ref{algebra-lemma-make-integral-trivial} 得出。
\end{proof}

\begin{lemma}
\label{obsolete-lemma-make-integral-less-trivial}
设 $\varphi : R \to S$ 是环映射。
设 $t \in S$ 满足关系
$\varphi(a_0) + \varphi(a_1)t + \ldots + \varphi(a_n) t^n = 0$。
置 $u_n = \varphi(a_n)$，$u_{n-1} = u_n t + \varphi(a_{n-1})$，
依此递推至 $u_1 = u_2 t + \varphi(a_1)$。
则 $u_n, u_{n-1}, \ldots, u_1$ 以及
$u_nt, u_{n-1}t, \ldots, u_1t$ 都在 $R$ 上整，
并且 $S$ 中的理想 $(\varphi(a_0), \ldots, \varphi(a_n))$ 与
$(u_n, \ldots, u_1)$ 相等。
\end{lemma}

\begin{proof}
我们对 $n$ 作归纳证明。由于 $u_n = \varphi(a_n)$，由
Algebra，引理 \ref{algebra-lemma-make-integral-trivial}
可得 $u_nt$ 在 $R$ 上整。当然
$u_n = \varphi(a_n)$ 在 $R$ 上整。于是
$u_{n - 1} = u_n t  + \varphi(a_{n - 1})$ 在 $R$ 上整（见
Algebra，引理 \ref{algebra-lemma-integral-closure-is-ring}），
并且有
$$
\varphi(a_0) + \varphi(a_1)t + \ldots + \varphi(a_{n - 1})t^{n - 1} +
u_{n - 1}t^{n - 1} = 0.
$$
因此将归纳假设应用于映射
$S' \to S$（其中 $S'$ 是 $S$ 中 $R$ 的整闭包）以及上面的等式，
可知 $u_{n-1}, \ldots, u_1$ 与 $u_{n-1}t, \ldots, u_1t$
也都属于 $S'$。关于理想的断言直接由
这些元素的形式以及 $u_1t + \varphi(a_0) = 0$ 得出。
\end{proof}

\begin{lemma}
\label{obsolete-lemma-make-integral-not-in-ideal}
设 $\varphi : R \to S$ 是环映射。
设 $t \in S$ 满足
关系 $\varphi(a_0) + \varphi(a_1)t + \ldots + \varphi(a_n) t^n = 0$。
设 $J \subset S$ 是理想，并且至少存在一个 $i$ 使
$\varphi(a_i) \not \in J$。
则存在 $u \in S$，$u \not\in J$，使得
$u$ 与 $ut$ 都在 $R$ 上整。
\end{lemma}

\begin{proof}
这立即由引理 \ref{obsolete-lemma-make-integral-less-trivial} 得出，
因为某个元素 $u_i$ 不属于 $J$。
\end{proof}

\noindent
下面两个引理给出一种描述由一个（非退化）等式
切出的 $\mathbf{P}^1_R$ 闭子概形的方法。

\begin{lemma}
\label{obsolete-lemma-P1}
设 $R$ 是环。
设 $F(X, Y) \in R[X, Y]$ 是次数为
$d$ 的齐次多项式。假设对 $R$ 的每个素理想 $\mathfrak p$，
$F$ 至少有一个系数不属于 $\mathfrak p$。
设 $S = R[X, Y]/(F)$ 为分次环。
则对所有 $n \geq d$，$R$-模 $S_n$
都是秩为 $d$ 的有限局部自由模。
\end{lemma}

\begin{proof}
$R$-模 $S_n$ 有如下表示：
$$
R[X, Y]_{n-d} \to R[X, Y]_n \to S_n \to 0.
$$
因此由 Algebra，引理 \ref{algebra-lemma-cokernel-flat}，
只需证明乘以 $F$ 诱导单射
$\kappa(\mathfrak p)[X, Y]
\to \kappa(\mathfrak p)[X, Y]$
对所有素理想 $\mathfrak p$ 都成立。
这由 $F$ 模 $\mathfrak p$ 不映为零多项式的假设立即可见。
秩的断言也同样显然。
\end{proof}

\begin{lemma}
\label{obsolete-lemma-rel-prime-pols}
设 $k$ 是域，且 $F, G \in k[X, Y]$ 齐次，
次数分别为 $d, e$。假设 $F, G$ 互素。
则乘以 $G$ 在 $S = k[X, Y]/(F)$ 上是单射。
\end{lemma}

\begin{proof}
这是一种定义“互素”的方式。若采用另一种
定义，可以证明它与此定义等价。
\end{proof}

\begin{lemma}
\label{obsolete-lemma-P1-localize}
设 $R$ 是环。设 $F(X, Y) \in R[X, Y]$ 是次数为
$d$ 的齐次多项式，且分次环 $S = R[X, Y]/(F)$。
设 $\mathfrak p \subset R$ 是素理想，且
$F$ 的某个系数不属于 $\mathfrak p$。
则存在 $f \in R$，$f \not\in \mathfrak p$，
整数 $e$ 及 $G \in R[X, Y]_e$，
使得乘以 $G$ 对所有 $n \geq d$ 诱导同构
$(S_n)_f \to (S_{n + e})_f$。
\end{lemma}

\begin{proof}
在证明过程中，我们可以用 $R_f$ 替换 $R$，其中
$f\in R$ 且 $f\not\in \mathfrak p$（有限次进行）。
首先作这样的替换，使得
$F$ 的某个系数在 $R$ 中可逆。
特别地，由引理 \ref{obsolete-lemma-P1}，现在对 $n \geq d$，模 $S_n$
是秩为 $d$ 的局部自由模。取任意 $G \in R[X, Y]_e$，使其在
$G$ 在 $\kappa(\mathfrak p)[X, Y]$ 中的像与 $F(X, Y)$ 的像
互素（对某个 $e$ 这是可能的）。
将 Algebra，引理 \ref{algebra-lemma-cokernel-flat} 应用于由乘以 $G$
从 $S_d \to S_{d + e}$ 诱导的映射。由 $G$ 的选择及
引理 \ref{obsolete-lemma-rel-prime-pols}，可知
$S_d \otimes \kappa(\mathfrak p) \to S_{d + e} \otimes \kappa(\mathfrak p)$
是双射。因此，对适当的 $f$ 用 $R_f$ 替换 $R$ 后，可设
$G : S_d \to S_{d + e}$ 是双射。这又意味着
对 $R$ 的所有素理想 $\mathfrak p'$，$G$ 在
$\kappa(\mathfrak p')[X, Y]$ 中的像与 $F$ 的像
互素。于是再次由 Algebra，引理 \ref{algebra-lemma-cokernel-flat}
可知所有映射
$G : S_d \to S_{d + e}$（$n \geq d$）都是同构。
\end{proof}

\begin{remark}
\label{obsolete-remark-algebra}
设 $R$ 是环。假设有 $F \in R[X, Y]_d$ 与
$G \in R[X, Y]_e$，置 $S = R[X, Y]/(F)$，并满足：(1) 对
所有 $n \geq d$，$S_n$ 是秩为 $d$ 的有限局部自由模；(2) 乘以 $G$
对所有 $n \geq d$ 诱导同构 $S_n \to S_{n + e}$。此时
可以如下定义有限局部自由的 $R$-代数
$A$：
\begin{enumerate}
\item 作为 $R$-模，$A = S_{ed}$；
\item 乘法 $A \times A \to A$ 由下列规则给出：
在 $S_{2ed}$ 中，$H_1 H_2 = H_3$ 当且仅当 $G^d H_3 = H_1 H_2$。
\end{enumerate}
这是合理的，因为乘以 $G^d$
诱导双射 $S_{de} \to S_{2de}$。
容易看出这定义了一个环结构。
需要注意一个容易混淆的事实：元素 $G^d$
定义了环 $A$ 的单位元。
\end{remark}

\begin{lemma}
\label{obsolete-lemma-finite-after-localization}
设 $R$ 是环，且 $f \in R$。
设有 $S$、$S'$ 及实线箭头，
组成下列环的交换图：
$$
\xymatrix{
& S'' \ar@{-->}[rd] \ar@{-->}[dd] &
\\
R \ar[rr] \ar@{-->}[ru] \ar[d] &  & S \ar[d]
\\
R_f \ar[r] & S' \ar[r] & S_f
}
$$
假设 $R_f \to S'$ 是有限映射。则可找到
有限环映射 $R \to S''$ 及图中所示的虚线箭头，
使得 $S' = (S'')_f$。
\end{lemma}

\begin{proof}
具体地，设 $S'$ 由
$x_i$ 在 $R_f$ 上生成，其中 $i = 1, \ldots, w$。令 $P_i(t) \in R_f[t]$
为满足 $P_i(x_i) = 0$ 的首一多项式。设
$P_i$ 的次数为 $d_i > 0$。写成
$P_i(t) = t^{d_i} + \sum_{j < d_i} (a_{ij}/f^n) t^j$，
其中 $n$ 统一。再将
$x_i$ 在 $S_f$ 中的像写成 $g_i / f^n$，
其中 $g_i \in S$ 适当。于是元素
$\xi_i = f^{nd_i} g_i^{d_i} + \sum_{j < d_i} f^{n(d_i - j)} a_{ij} g_i^j$
$S$ 被 $f$ 的某个幂湮灭。
因此将 $n$ 增大为 $n'$（这会把
$g_i$ 替换为 $f^{n' - n}g_i$）后，可设 $\xi_i = 0$。
于是 $S'$ 由元素
$f^n x_i$ 生成；每个元素都是下列首一多项式的根：
$Q_i(t) = t^{d_i} + \sum_{j < d_i} f^{n(d_i - j)} a_{ij} t^j$，
其系数在 $R$ 中。按构造，
$Q_i(f^ng_i) = 0$ 在 $S$ 中。因此得到有限 $R$-代数
$S'' = R[z_1, \ldots, z_w]/(Q_1(z_1), \ldots, Q_w(z_w))$，
它适合于上面的交换图。
映射 $\alpha : S'' \to S$ 将 $z_i$ 映到 $f^ng_i$，
映射 $\beta : S'' \to S'$ 将 $z_i$ 映到 $f^nx_i$。
此时 $\beta$ 可能尚未诱导
同构 $(S'')_f \cong S'$。
目前只知道该映射是满射。问题在于，可能存在
元素 $h/f^n \in (S'')_f$ 映到 $S'$ 中的零但自身不为零。此时 $\beta(h)$
是 $S$ 中的元素，且对某个 $N$ 有 $f^N \beta(h) = 0$。因此 $f^N h$ 属于理想
$J = \{h \in S'' \mid \alpha(h) = 0 \text{ 且 }
\beta(h) = 0\}$ 的 $S''$。容易看出 $S''/J$ 即可满足要求。
\end{proof}




\section{形式光滑的环映射}
\label{obsolete-section-formally-smooth}

\begin{lemma}
\label{obsolete-lemma-formally-smooth-smooth}
设 $R$ 是环，$S$ 是 $R$-代数。
若 $S$ 有限表示且形式光滑于 $R$，
则 $S$ 光滑于 $R$。
\end{lemma}

\begin{proof}
见 Algebra，命题 \ref{algebra-proposition-smooth-formally-smooth}。
\end{proof}

\begin{remark}
\label{obsolete-remark-equation-derivatives}
该标签原先指向 Algebraization of Formal Spaces，命题
\ref{restricted-proposition-approximate} 的证明中的一个等式，
由于材料重新编排而不再使用。
\end{remark}

\begin{remark}
\label{obsolete-remark-equation-ci}
该标签原先指向 Algebraization of Formal Spaces，命题
\ref{restricted-proposition-approximate} 的证明中的一个等式，
由于材料重新编排而不再使用。
\end{remark}

\begin{remark}
\label{obsolete-remark-equation-in-ideal}
该标签原先指向 Algebraization of Formal Spaces，命题
\ref{restricted-proposition-approximate} 的证明中的一个等式，
由于材料重新编排而不再使用。
\end{remark}

\begin{remark}
\label{obsolete-remark-equation-derivatives-analogue}
该标签原先指向 Algebraization of Formal Spaces，命题
\ref{restricted-proposition-approximate} 的证明中的一个等式，
由于材料重新编排而不再使用。
\end{remark}

\begin{remark}
\label{obsolete-remark-equation-go-down}
该标签原先指向 Algebraization of Formal Spaces，引理
\ref{restricted-lemma-lift-approximation} 的证明中的一个等式，
由于材料重新编排而不再使用。
\end{remark}

\begin{lemma}
\label{obsolete-lemma-get-morphism-general}
设 $A$ 是诺特环，且 $I \subset A$ 是理想。
设 $t$ 是 $I$ 的最小生成元个数。
设 $C$ 是关于 $I$-进拓扑完备的诺特 $A$-代数。
存在只依赖于 $I \subset A \to C$ 的整数 $d \geq 0$，
具有如下性质：给定
\begin{enumerate}
\item $c \geq 0$ 及 Algebraization of Formal Spaces，
等式 (\ref{restricted-equation-C-prime}) 中的 $B$，满足对 $a \in I^c$，
乘以 $a$ 在 $\NL_{B/A}^\wedge$ 上于 $D(B)$ 中为零；
\item 整数 $n > 2t\max(c, d)$；
\item $A/I^n$-代数映射 $\psi_n : B/I^nB \to C/I^nC$；
\end{enumerate}
存在 $A$-代数映射 $\varphi : B \to C$，使得
$\psi_n \bmod I^{m - c} = \varphi \bmod I^{m - c}$，
其中 $m = \lfloor \frac{n}{t} \rfloor$。
\end{lemma}

\begin{proof}
这个引理已被更强的
Algebraization of Formal Spaces，引理
\ref{restricted-lemma-get-morphism-general-better} 所废止。
事实上，我们将由该引理推出本引理。

\medskip\noindent
设 $I \subset A \to C$ 如上面的陈述。
记 $d(\text{Gr}_I(C))$ 与 $q(\text{Gr}_I(C))$ 为
Local Cohomology，小节 \ref{local-cohomology-section-uniform} 中得到的整数。
注意，$t$ 是 $IC$ 的最小生成元个数的上界，因此
$d(\text{Gr}_I(C)) + 1 \leq t$；见
Local Cohomology，小节 \ref{local-cohomology-section-uniform} 中的讨论。
可假设 $t \geq 1$，因为否则引理没有内容。
我们断言取
$$
d = q(\text{Gr}_I(C))
$$
具体地，设 $c$、$B$、$n$、$\psi_n$ 如上面的陈述。
则有
$$
n > 2t\max(c, d) \Rightarrow n \geq 2tc + 1 \Rightarrow
n \geq 2(d(\text{Gr}_I(C)) + 1)c + 1
$$
另一方面，有
$$
n > 2t\max(c, d) \Rightarrow n > t(c + d) \Rightarrow
n \geq q(C) + tc \geq q(\text{Gr}_I(C)) + (d(\text{Gr}_I(C)) + 1)c
$$
因此 Algebraization of Formal Spaces，引理
\ref{restricted-lemma-get-morphism-general-better} 的假设
得到满足，从而得到 $A$-代数同态
$\varphi : B \to C$，它与 $\psi_n$ 在
模 $I^{n - (d(\text{Gr}_I(C)) + 1)c}C$ 上同余。
由于
\begin{align*}
n - (d(\text{Gr}_I(C)) + 1)c
& = \frac{n}{t} + \frac{(t - 1)n}{t} - (d(\text{Gr}_I(C)) + 1)c \\
& \geq \frac{n}{t} + \frac{(d(\text{Gr}_I(C))n}{t} - (d(\text{Gr}_I(C)) + 1)c \\
& > \frac{n}{t} + \frac{d(\text{Gr}_I(C))2tc}{t} - (d(\text{Gr}_I(C)) + 1)c \\
& = \frac{n}{t} + 2d(\text{Gr}_I(C))c - (d(\text{Gr}_I(C)) + 1)c \\
& = \frac{n}{t} + d(\text{Gr}_I(C))c - c \\
& \geq m - c
\end{align*}
可知如所需，$\varphi$ 与 $\psi_n$ 在模 $I^{m - c}C$ 上同余。
\end{proof}




\section{网站与层}
\label{obsolete-section-sites}

\begin{remark}[不存在从零延拓到推前的映射]
\label{obsolete-remark-from-shriek-to-star}
设 $U$ 是网站 $\mathcal{C}$ 的对象。对 $\mathcal{C}/U$ 上的任意阿贝尔层
$\mathcal{G}$，可以问是否存在典范映射
$$
c : j_{U!}\mathcal{G} \longrightarrow j_{U*}\mathcal{G}
$$
构造这样的映射等价于构造映射
$j_U^{-1}j_{U!}\mathcal{G} \to \mathcal{G}$。
注意限制与层化交换。
因此可以使用
Modules on Sites，引理 \ref{sites-modules-lemma-extension-by-zero} 中的预层。
所以只需对 $V/U$ 定义映射
$$
\bigoplus\nolimits_{\varphi \in \Mor_\mathcal{C}(V, U)}
\mathcal{G}(V \xrightarrow{\varphi} U)
\longrightarrow
\mathcal{G}(V/U)
$$
使其与限制相容。看起来可以令该映射在除以下求和项外的
所有求和项上为零：当 $\varphi$ 是结构态射 $\varphi_0 : V \to U$ 时，
在该项上取值 $1$。
然而这与限制映射不相容：具体地，若
$\alpha : V' \to V$ 是 $\mathcal{C}$ 中的态射，记
$V'/U$ 为 $\mathcal{C}/U$ 中结构
态射为 $\varphi'_0 = \varphi_0 \circ \alpha$ 的对象。
需要检查下图
$$
\xymatrix{
\bigoplus\nolimits_{\varphi \in \Mor_\mathcal{C}(V, U)}
\mathcal{G}(V \xrightarrow{\varphi} U)
\ar[d] \ar[r] &
\mathcal{G}(V/U) \ar[d] \\
\bigoplus\nolimits_{\varphi' \in \Mor_\mathcal{C}(V', U)}
\mathcal{G}(V' \xrightarrow{\varphi'} U)
\ar[r] &
\mathcal{G}(V'/U)
}
$$
交换。问题在于，可能存在不同于 $\varphi_0$ 的态射
$\varphi : V \to U$，满足
$\varphi \circ \alpha = \varphi'_0$。
因此左侧竖直箭头会把对应于 $\varphi$ 的求和项送到下方
水平箭头取值为 $1$ 的求和项，图表几乎肯定不交换。
\end{remark}




\section{上同调}
\label{obsolete-section-cohomology}

\noindent
关于上同调的已废弃引理。

\begin{lemma}
\label{obsolete-lemma-ML-general}
设 $I$ 是环 $A$ 的理想，$X$ 是位于 $\Spec(A)$ 上的概形。设
$$
\ldots \to \mathcal{F}_3 \to \mathcal{F}_2 \to \mathcal{F}_1
$$
是 $\mathcal{O}_X$-模的逆系统，
满足 $\mathcal{F}_n = \mathcal{F}_{n + 1}/I^n\mathcal{F}_{n + 1}$。
假设
$$
\bigoplus\nolimits_{n \geq 0} H^1(X, I^n\mathcal{F}_{n + 1})
$$
作为分次
$\bigoplus_{n \geq 0} I^n/I^{n + 1}$-模满足升链条件。
则逆系统 $M_n = \Gamma(X, \mathcal{F}_n)$ 满足
Mittag--Leffler 条件。
\end{lemma}

\begin{proof}
这是更一般的
Cohomology，引理 \ref{cohomology-lemma-ML-general} 的特殊情形。
\end{proof}

\begin{lemma}
\label{obsolete-lemma-ML-general-better}
设 $I$ 是环 $A$ 的理想，$X$ 是 $\Spec(A)$ 上的概形。设
$$
\ldots \to \mathcal{F}_3 \to \mathcal{F}_2 \to \mathcal{F}_1
$$
为 $\mathcal{O}_X$-模的逆系统，
满足 $\mathcal{F}_n = \mathcal{F}_{n + 1}/I^n\mathcal{F}_{n + 1}$。
给定 $n$，定义
$$
H^1_n =
\bigcap\nolimits_{m \geq n}
\Im\left(
H^1(X, I^n\mathcal{F}_{m + 1}) \to H^1(X, I^n\mathcal{F}_{n + 1})
\right)
$$
若 $\bigoplus H^1_n$ 作为分次
$\bigoplus_{n \geq 0} I^n/I^{n + 1}$-模满足升链条件，则逆系统
$M_n = \Gamma(X, \mathcal{F}_n)$ 满足 Mittag--Leffler 条件。
\end{lemma}

\begin{proof}
这是更一般的
Cohomology，引理 \ref{cohomology-lemma-ML-general-better} 的特殊情形。
\end{proof}

\begin{lemma}
\label{obsolete-lemma-topology-I-adic-general}
设 $I$ 是环 $A$ 的有限生成理想。
设 $X$ 是位于 $\Spec(A)$ 上的概形。设
$$
\ldots \to \mathcal{F}_3 \to \mathcal{F}_2 \to \mathcal{F}_1
$$
是 $\mathcal{O}_X$-模的逆系统，满足
$\mathcal{F}_n = \mathcal{F}_{n + 1}/I^n\mathcal{F}_{n + 1}$。假设
$$
\bigoplus\nolimits_{n \geq 0} H^0(X, I^n\mathcal{F}_{n + 1})
$$
作为分次
$\bigoplus_{n \geq 0} I^n/I^{n + 1}$-模满足升链条件。
则 $M = \lim \Gamma(X, \mathcal{F}_n)$ 上的极限拓扑
就是 $I$-进拓扑。
\end{lemma}

\begin{proof}
这是更一般的
Cohomology，引理 \ref{cohomology-lemma-topology-I-adic-general} 的特殊情形。
\end{proof}

\begin{lemma}
\label{obsolete-lemma-pullback-K-flat}
设 $(\Sh(\mathcal{C}), \mathcal{O}_\mathcal{C})$ 是环化拓扑斯。
对任意 $\mathcal{O}_\mathcal{C}$-模复形 $\mathcal{G}^\bullet$，
存在拟同构 $\mathcal{K}^\bullet \to \mathcal{G}^\bullet$，使得
对任意环化拓扑斯态射
$f : (\Sh(\mathcal{D}), \mathcal{O}_\mathcal{D}) \to
(\Sh(\mathcal{C}), \mathcal{O}_\mathcal{C})$，$f^*\mathcal{K}^\bullet$ 是
$\mathcal{O}_\mathcal{D}$-模的 K-平坦复形。
\end{lemma}

\begin{proof}
这由 Cohomology on Sites，引理
\ref{sites-cohomology-lemma-K-flat-resolution} 与
\ref{sites-cohomology-lemma-pullback-K-flat} 得出。
\end{proof}

\begin{remark}
\label{obsolete-remark-pullback-K-flat}
本注记原先讨论 K-平坦复形的拉回是否仍为 K-平坦，
但现在已被
Cohomology on Sites，引理 \ref{sites-cohomology-lemma-pullback-K-flat} 废止。
\end{remark}

\noindent
下面的引理在特殊情形下计算
导出极限的上同调层。

\begin{lemma}
\label{obsolete-lemma-Rlim-of-system}
设 $(\mathcal{C}, \mathcal{O})$ 是环化网站。设 $(K_n)$
是 $D(\mathcal{O})$ 中对象的逆系统。
设 $\mathcal{B} \subset \Ob(\mathcal{C})$ 是子集。
设 $d \in \mathbf{N}$。假设
\begin{enumerate}
\item 对所有 $n$，$K_n$ 是 $D^+(\mathcal{O})$ 中的对象；
\item 对 $q \in \mathbf{Z}$，存在
$n(q)$，使得当 $n \geq n(q)$ 时，$H^q(K_{n + 1}) \to H^q(K_n)$ 是同构；
\item $\mathcal{C}$ 的每个对象都有一个覆盖，其成员是
$\mathcal{B}$ 的元素；
\item 对每个 $U \in \mathcal{B}$，有 $H^p(U, H^q(K_n)) = 0$，
对所有 $q$ 及 $p > d$ 均成立。
\end{enumerate}
则对所有 $m \in \mathbf{Z}$，有 $H^m(R\lim K_n) = \lim H^m(K_n)$。
\end{lemma}

\begin{proof}
置 $K = R\lim K_n$。设 $U \in \mathcal{B}$。对每个 $n$ 有谱
序列
$$
H^p(U, H^q(K_n)) \Rightarrow H^{p + q}(U, K_n)
$$
它收敛，因为 $K_n$ 有下界；见
Derived Categories，引理 \ref{derived-lemma-two-ss-complex-functor}。
固定 $m \in \mathbf{Z}$，由假设 (4) 可知，只有
$H^p(U, H^q(K_n))$ 对 $H^m(U, K_n)$ 有贡献，
其中 $0 \leq p \leq d$ 且 $m - d \leq q \leq m$。由假设 (2)，
当 $n \geq \max{n(m), \ldots, n(m - d)}$ 时，
$H^m(U, K_{n + 1}) \to H^m(U, K_n)$ 是同构。函子 $R\Gamma(U, -)$
与导出极限交换，依据
Injectives，引理 \ref{injectives-lemma-RF-commutes-with-Rlim}。
因此
$$
H^m(U, K) = H^m(R\lim R\Gamma(U, K_n))
$$
另一方面，我们刚刚看到复形 $R\Gamma(U, K_n)$ 的上同调群最终恒定。因此由
More on Algebra，注记 \ref{more-algebra-remark-compare-derived-limit}，
可知对某个与 $U \in \mathcal{B}$ 无关的界，所有 $n \gg 0$ 都有
$H^m(U, K)$ 等于 $H^m(U, K_n)$。取这样的 $n$。最后回忆，$H^m(K)$ 是预层
$U \mapsto H^m(U, K)$ 的层化，而 $H^m(K_n)$ 是预层
$U \mapsto H^m(U, K_n)$ 的层化。在 $\mathcal{B}$ 的元素上，这些预层取值相同。
因此假设 (3) 保证层化也相同。
引理得证。
\end{proof}

\begin{lemma}
\label{obsolete-lemma-trivialities-cohomological-descent-abelian}
在 Simplicial Spaces，情形
\ref{spaces-simplicial-situation-simplicial-site} 中，
设 $a_0$ 是指向网站 $\mathcal{D}$ 的增广，如
Simplicial Spaces，注记 \ref{spaces-simplicial-remark-augmentation-site} 所述。
给定严格满弱 Serre 子范畴
$$
\mathcal{A} \subset \textit{Ab}(\mathcal{D}),\quad
\mathcal{A}_n \subset \textit{Ab}(\mathcal{C}_n)
$$
则
\begin{enumerate}
\item[(1)]
$\mathcal{C}_{total}$ 上所有满足其到每个 $\mathcal{C}_n$ 的限制属于 $\mathcal{A}_n$（对所有 $n$）的
阿贝尔层 $\mathcal{F}$ 构成严格满弱 Serre 子范畴
$\mathcal{A}_{total} \subset \textit{Ab}(\mathcal{C}_{total})$。
\end{enumerate}
若对所有 $n$，$a_n^{-1}$ 将 $\mathcal{A}$ 映入 $\mathcal{A}_n$，则
\begin{enumerate}
\item[(2)] $a^{-1}$ 将 $\mathcal{A}$ 映入 $\mathcal{A}_{total}$；
\item[(3)] $a^{-1}$ 将 $D_\mathcal{A}(\mathcal{D})$ 映入
$D_{\mathcal{A}_{total}}(\mathcal{C}_{total})$。
\end{enumerate}
若对所有 $n, q$，$R^qa_{n, *}$ 将 $\mathcal{A}_n$ 映入 $\mathcal{A}$，则
\begin{enumerate}
\item[(4)] 对所有 $q$，$R^qa_*$ 将 $\mathcal{A}_{total}$ 映入 $\mathcal{A}$；
并且
\item[(5)] $Ra_*$ 将 $D_{\mathcal{A}_{total}}^+(\mathcal{C}_{total})$
映入 $D_\mathcal{A}^+(\mathcal{D})$。
\end{enumerate}
\end{lemma}

\begin{proof}
唯一有实质内容的断言是 (4) 与 (5)。
(4) 由下列谱序列推出：
Simplicial Spaces，引理
\ref{spaces-simplicial-lemma-augmentation-spectral-sequence}
以及 Homology，引理 \ref{homology-lemma-biregular-ss-converges}。
然后考虑由截断给出的
$D_{\mathcal{A}_{total}}^+(\mathcal{C}_{total})$ 中对象
$K$ 的典范过滤所关联的谱序列，即得 (5)。
细节略去。
\end{proof}

\begin{remark}
\label{obsolete-remark-cohomology-topics}
该标签原先指向上调章节中列出待处理主题的一个小节。
\end{remark}

\begin{remark}
\label{obsolete-remark-sites-cohomology-topics}
该标签原先指向上调章节中列出待处理主题的一个小节。
\end{remark}

\begin{remark}
\label{obsolete-remark-V-implies-C}
该标签原先指向 Cohomology on Sites，引理
\ref{sites-cohomology-lemma-V-implies-C-general} 在
Cohomology on Sites，引理 \ref{sites-cohomology-lemma-compare-qc-zar}
所述情形下的特殊情形。
\end{remark}

\begin{remark}
\label{obsolete-remark-V-implies-cohomology}
该标签原先指向 Cohomology on Sites，引理
\ref{sites-cohomology-lemma-V-implies-cohomology-general} 在
Cohomology on Sites，引理 \ref{sites-cohomology-lemma-compare-qc-zar}
所述情形下的特殊情形。
\end{remark}

\begin{remark}
\label{obsolete-remark-induction-step-V-C}
该标签原先指向 Cohomology on Sites，引理
\ref{sites-cohomology-lemma-induction-step-V-C-general} 在
Cohomology on Sites，引理 \ref{sites-cohomology-lemma-compare-qc-zar}
所述情形下的特殊情形。
\end{remark}

\begin{remark}
\label{obsolete-remark-V-implies-C-etale-fppf}
该标签原先指向 Cohomology on Sites，引理
\ref{sites-cohomology-lemma-V-implies-C-general} 在
\'Etale Cohomology，引理 \ref{etale-cohomology-lemma-compare-fppf-etale}
所述情形下的特殊情形。
\end{remark}

\begin{remark}
\label{obsolete-remark-V-implies-cohomology-etale-fppf}
该标签原先指向 Cohomology on Sites，引理
\ref{sites-cohomology-lemma-V-implies-cohomology-general} 在
\'Etale Cohomology，引理 \ref{etale-cohomology-lemma-compare-fppf-etale}
所述情形下的特殊情形。
\end{remark}

\begin{remark}
\label{obsolete-remark-induction-step-V-C-etale-fppf}
该标签原先指向 Cohomology on Sites，引理
\ref{sites-cohomology-lemma-induction-step-V-C-general} 在
\'Etale Cohomology，引理 \ref{etale-cohomology-lemma-compare-fppf-etale}
所述情形下的特殊情形。
\end{remark}

\begin{remark}
\label{obsolete-remark-V-implies-C-etale-ph}
该标签原先指向 Cohomology on Sites，引理
\ref{sites-cohomology-lemma-V-implies-C-general} 在
\'Etale Cohomology，引理 \ref{etale-cohomology-lemma-compare-ph-etale}
所述情形下的特殊情形。
\end{remark}

\begin{remark}
\label{obsolete-remark-V-implies-cohomology-etale-ph}
该标签原先指向 Cohomology on Sites，引理
\ref{sites-cohomology-lemma-V-implies-cohomology-general} 在
\'Etale Cohomology，引理 \ref{etale-cohomology-lemma-compare-ph-etale}
所述情形下的特殊情形。
\end{remark}

\begin{remark}
\label{obsolete-remark-V-implies-cohomology-etale-ph-extra}
该标签原先指向 Cohomology on Sites，引理
\ref{sites-cohomology-lemma-V-implies-cohomology-extra-general} 在
\'Etale Cohomology，引理 \ref{etale-cohomology-lemma-compare-ph-etale}
所述情形下的特殊情形。
\end{remark}

\begin{remark}
\label{obsolete-remark-make-class-zero}
该标签原先指向 Cohomology on Sites，引理
\ref{sites-cohomology-lemma-make-class-zero-general} 在
\'Etale Cohomology，引理 \ref{etale-cohomology-lemma-compare-ph-etale}
所述情形下的特殊情形。
\end{remark}

\begin{remark}
\label{obsolete-remark-induction-step-V-C-etale-ph}
该标签原先指向 Cohomology on Sites，引理
\ref{sites-cohomology-lemma-induction-step-V-C-general} 在
\'Etale Cohomology，引理 \ref{etale-cohomology-lemma-compare-ph-etale}
所述情形下的特殊情形。
\end{remark}

\begin{remark}
\label{obsolete-remark-V-implies-C-etale-h}
该标签原先指向 Cohomology on Sites，引理
\ref{sites-cohomology-lemma-V-implies-C-general} 在
\'Etale Cohomology，引理 \ref{etale-cohomology-lemma-compare-h-etale}
所述情形下的特殊情形。
\end{remark}

\begin{remark}
\label{obsolete-remark-V-implies-cohomology-etale-h}
该标签原先指向 Cohomology on Sites，引理
\ref{sites-cohomology-lemma-V-implies-cohomology-general} 在
\'Etale Cohomology，引理 \ref{etale-cohomology-lemma-compare-h-etale}
所述情形下的特殊情形。
\end{remark}

\begin{remark}
\label{obsolete-remark-V-implies-cohomology-etale-h-extra}
该标签原先指向 Cohomology on Sites，引理
\ref{sites-cohomology-lemma-V-implies-cohomology-extra-general} 在
\'Etale Cohomology，引理 \ref{etale-cohomology-lemma-compare-h-etale}
所述情形下的特殊情形。
\end{remark}

\begin{remark}
\label{obsolete-remark-induction-step-V-C-etale-h}
该标签原先指向 Cohomology on Sites，引理
\ref{sites-cohomology-lemma-induction-step-V-C-general} 在
\'Etale Cohomology，引理 \ref{etale-cohomology-lemma-compare-h-etale}
所述情形下的特殊情形。
\end{remark}

\begin{remark}
\label{obsolete-remark-how-used}
该标签原先位于 \'etale 上同调章节，但因重新组织已不适合留在那里。
标签内容如下：
\'Etale Cohomology，引理 \ref{etale-cohomology-lemma-when-ctf}
可用于证明：若 $f : X \to Y$ 是概形的分离有限型态射且 $Y$ 诺特，
则 $Rf_!$ 诱导函子
$D_{ctf}(X_\etale, \Lambda) \to D_{ctf}(Y_\etale, \Lambda)$。
当 $Y$ 是域的谱且 $X$ 是曲线时，这一论证的例子见 The Trace Formula，
命题 \ref{trace-proposition-projective-curve-constructible-cohomology}。
\end{remark}

\begin{lemma}
\label{obsolete-lemma-glue-f-upper-shriek}
设 $f : X \to Y$ 是概形的局部拟有限态射。
存在唯一函子
$f^! : \textit{Ab}(Y_\etale) \to \textit{Ab}(X_\etale)$，满足
\begin{enumerate}
\item 对任意开浸入 $j : U \to X$，若 $f \circ j$ 分离，则存在典范同构
$j^! \circ f^! = (f \circ j)^!$；并且
\item 当 $U \subset U' \subset X$ 时，这些同构与 More \'Etale Cohomology，引理
\ref{more-etale-lemma-upper-shriek-restriction} 中的同构相容。
\end{enumerate}
\end{lemma}

\begin{proof}
立即由 More \'Etale Cohomology，引理
\ref{more-etale-lemma-lqf-f-upper-shriek} 与
\ref{more-etale-lemma-upper-shriek-restriction} 得出。
\end{proof}

\begin{proposition}
\label{obsolete-proposition-lqf-shriek}
设 $f : X \to Y$ 是局部拟有限态射。存在伴随函子
$f_! : \textit{Ab}(X_\etale) \to \textit{Ab}(Y_\etale)$
与 $f^! : \textit{Ab}(Y_\etale) \to \textit{Ab}(X_\etale)$，
具有以下性质
\begin{enumerate}
\item 函子 $f^!$ 是 More \'Etale Cohomology 中由
引理 \ref{more-etale-lemma-lqf-f-upper-shriek} 构造的函子；
\item 对任意开浸入 $j : U \to X$，若 $f \circ j$ 分离，存在典范同构
$f_! \circ j_! = (f \circ j)_!$；并且
\item 当 $U \subset U' \subset X$ 时，这些同构与 More \'Etale Cohomology，
引理 \ref{more-etale-lemma-f-shriek-composition} 中的同构相容。
\end{enumerate}
\end{proposition}

\begin{proof}
见 More \'Etale Cohomology，小节
\ref{more-etale-section-finite-support} 与
\ref{more-etale-section-duality-locally-quasi-finite}。
\end{proof}

\begin{lemma}
\label{obsolete-lemma-lqf-f-upper-shriek-stalk}
设 $f : X \to Y$ 是局部拟有限的概形态射。
对阿贝尔群 $A$ 与几何点
$\overline{y} : \Spec(k) \to Y$，有
$f^!(\overline{y}_*A) = \prod\nolimits_{f(\overline{x}) = \overline{y}}
\overline{x}_*A$.
\end{lemma}

\begin{proof}
由 More \'Etale Cohomology，引理
\ref{more-etale-lemma-lqf-f-upper-shriek} 中对应的陈述推出。
\end{proof}

\begin{lemma}
\label{obsolete-lemma-lqf-f-shriek-composition}
设 $f : X \to Y$ 与 $g : Y \to Z$ 是可复合的局部拟有限概形态射。
则 $g_! \circ f_! = (g \circ f)_!$，并且
$f^! \circ g^! = (g \circ f)^!$。
\end{lemma}

\begin{proof}
这是 More \'Etale Cohomology 中引理
\ref{more-etale-lemma-lqf-shriek-composition} 与
\ref{more-etale-lemma-upper-shriek-restriction} 的组合。
\end{proof}









\section{微分分次代数}
\label{obsolete-section-dga}


\begin{lemma}
\label{obsolete-lemma-P-not-preserved-base-change}
设 $(A, \text{d})$ 与 $(B, \text{d})$ 是微分分次代数。
设 $N$ 是具有性质 (P) 的微分分次 $(A, B)$-双模，设 $M$ 是具有性质 (P)
的微分分次 $A$-模。则 $Q = M \otimes_A N$ 是微分分次 $B$-模，
在 $D(B)$ 中表示 $M \otimes_A^\mathbf{L} N$，并带有滤链
$$
0 = F_{-1}Q \subset F_0Q \subset F_1Q \subset \ldots \subset Q
$$
其微分分次子模满足 $Q = \bigcup F_pQ$；包含映射
$F_iQ \to F_{i + 1}Q$ 是可容许单射；商
$F_{i + 1}Q/F_iQ$ 作为微分分次 $B$-模同构于
$(A \otimes_R B)[k]$ 的直和。
\end{lemma}

\begin{proof}
在 $M$ 与 $N$ 上选取滤链 $F_\bullet$。考虑 $Q = M \otimes_A N$ 上由下式给出的滤链
$$
F_n(Q) = \sum\nolimits_{i + j = n}
F_i(M) \otimes_A F_j(N)
$$
这显然是微分分次 $B$-子模。我们有
$$
F_n(Q)/F_{n - 1}(Q) =
\bigoplus\nolimits_{i + j = n}
F_i(M)/F_{i - 1}(M) \otimes_A F_j(N)/F_{j - 1}(N)
$$
例如，这是因为 $M$ 的滤链在分次 $A$-模范畴中是可分裂的。由于假设右侧的商同构于
$A$ 与 $A \otimes_R B$ 的移位直和，并且
$A \otimes_A (A \otimes_R B) = A \otimes_R B$，
所以左侧作为微分分次 $B$-模是 $A \otimes_R B$ 的移位直和。
（警告：$Q$ 没有 $(A, B)$-双模结构。）
这证明了引理的第一个断言。第二个断言直接由
Differential Graded Algebra，引理 \ref{dga-lemma-derived-bc} 中函子的定义得到。
\end{proof}












\section{单纯方法}
\label{obsolete-section-simplicial}


\begin{lemma}
\label{obsolete-lemma-equiv}
假设与记号同 Simplicial，引理 \ref{simplicial-lemma-section}。
存在态射 $f$ 的截面 $g : U \to V$，且复合 $g \circ f$ 在 $V$ 上与恒等态射同伦等价。
特别地，态射 $f$ 是同伦等价。
\end{lemma}

\begin{proof}
立即由 Simplicial，引理 \ref{simplicial-lemma-section} 与
\ref{simplicial-lemma-trivial-kan-homotopy} 得出。
\end{proof}

\begin{lemma}
\label{obsolete-lemma-cosk-hom-deltak}
设 $\mathcal{C}$ 是具有有限余积与有限极限的范畴。设 $X$ 是 $\mathcal{C}$ 中的对象。
设 $k \geq 0$。规范映射
$$
\Hom(\Delta[k], X)
\longrightarrow
\text{cosk}_1 \text{sk}_1 \Hom(\Delta[k], X)
$$
是同构。
\end{lemma}

\begin{proof}
对任意单纯对象 $V$，有
\begin{eqnarray*}
\Mor(V, \text{cosk}_1 \text{sk}_1 \Hom(\Delta[k], X))
& = &
\Mor(\text{sk}_1 V, \text{sk}_1 \Hom(\Delta[k], X)) \\
& = &
\Mor(i_{1!} \text{sk}_1 V, \Hom(\Delta[k], X)) \\
& = &
\Mor(i_{1!} \text{sk}_1 V \times \Delta[k], X)
\end{eqnarray*}
第一个等式来自 $\text{sk}$ 与 $\text{cosk}$ 的伴随性，第二个等式来自
$i_{1!}$ 与 $\text{sk}_1$ 的伴随性，第三个等式来自
Simplicial，定义 \ref{simplicial-definition-hom-from-simplicial-set}，
其中最后的 $X$ 表示取值为 $X$ 的常值单纯对象。由 Simplicial，引理
\ref{simplicial-lemma-augmentation-howto}，该集合中的元素只依赖于
$i_{1!} \text{sk}_1 V \times \Delta[k]$ 的 $0$ 次与 $1$ 次项。
这些项与 $V \times \Delta[k]$ 的 $0$ 次与 $1$ 次项一致；见
Simplicial，引理 \ref{simplicial-lemma-recovering-U-for-real}。
因此上面的集合等于
$\Mor(V \times \Delta[k], X) = \Mor(V, \Hom(\Delta[k], X))$.
\end{proof}

\begin{lemma}
\label{obsolete-lemma-cosk0-hom-deltak}
设 $\mathcal{C}$ 是范畴。设 $X$ 是 $\mathcal{C}$ 中的对象，并且自积
$X \times \ldots \times X$ 存在。设 $k \geq 0$，令 $C[k]$ 如
Simplicial，例 \ref{simplicial-example-simplex-cosimplicial-set} 中所述。
采用 Simplicial，引理 \ref{simplicial-lemma-morphism-into-product} 的记号，
规范映射
$$
\Hom(C[k], X)_1
\longrightarrow
(\text{cosk}_0 \text{sk}_0 \Hom(C[k], X))_1
$$
可识别为映射
$$
\prod\nolimits_{\alpha : [k] \to [1]} X
\longrightarrow
X \times X
$$
它是投影到 $\alpha$ 为常值映射的那些因子的投影。
\end{lemma}

\begin{proof}
这在 Hypercoverings，引理 \ref{hypercovering-lemma-covering} 的证明中给出。
\end{proof}




\section{关于概形的结果}
\label{obsolete-section-devissage}

\noindent
看似多余的引理。

\begin{lemma}
\label{obsolete-lemma-stein-projective}
设 $(R, \mathfrak m, \kappa)$ 是局部环。设 $X \subset \mathbf{P}^n_R$ 是闭子概形。
假设 $R = \Gamma(X, \mathcal{O}_X)$。则特殊纤维 $X_k$ 几何连通。
\end{lemma}

\begin{proof}
这是 More on Morphisms，定理
\ref{more-morphisms-theorem-stein-factorization-general}.
\end{proof}

\begin{lemma}
\label{obsolete-lemma-property-irreducible-higher-rank}
设 $X$ 是诺特概形。设 $Z_0 \subset X$ 是不可约闭子集，其一般点为 $\xi$。
设 $\mathcal{P}$ 是 $X$ 上凝聚层的性质，满足
\begin{enumerate}
\item 对任意凝聚层的短正合列，若其中三个对象中的两个具有性质 $\mathcal{P}$，
则第三个也具有该性质。
\item 若凝聚层的直和具有 $\mathcal{P}$，则该性质对两个直项均成立。
\item 对每个积分闭子概形 $Z \subset Z_0 \subset X$，若 $Z \not = Z_0$，且每个拟凝聚理想层
$\mathcal{I} \subset \mathcal{O}_Z$，都有 $(Z \to X)_*\mathcal{I}$ 具有 $\mathcal{P}$。
\item 存在 $X$ 上的凝聚层 $\mathcal{G}$，使得
\begin{enumerate}
\item $\text{Supp}(\mathcal{G}) = Z_0$；
\item $\mathcal{G}_\xi$ 被 $\mathfrak m_\xi$ 湮灭；
\item $\mathcal{G}$ 具有性质 $\mathcal{P}$。
\end{enumerate}
\end{enumerate}
则 $X$ 上支集包含于 $Z_0$ 的每个凝聚层 $\mathcal{F}$ 都具有性质 $\mathcal{P}$。
\end{lemma}

\begin{proof}
证明是 Cohomology of Schemes，引理
\ref{coherent-lemma-property-irreducible} 证明的变体。
以完全相同的方式可见，支集严格包含于 $Z_0$ 的任何凝聚层都具有性质 $\mathcal{P}$。

\medskip\noindent
考虑如 (3) 中的凝聚层 $\mathcal{G}$。由 Cohomology of Schemes，引理
\ref{coherent-lemma-prepare-filter-irreducible}，在 $Z_0$ 上存在理想层 $\mathcal{I}$，以及短正合列
$$
0 \to
\left((Z_0 \to X)_*\mathcal{I}\right)^{\oplus r} \to
\mathcal{G} \to
\mathcal{Q} \to 0
$$
其中 $\mathcal{Q}$ 的支集严格包含于 $Z_0$。特别地，由于 $\mathcal{G}$ 的支集等于 $Z$，
有 $r > 0$ 且 $\mathcal{I}$ 非零。由于 $\mathcal{Q}$ 具有性质 $\mathcal{P}$，我们得到
$\left((Z_0 \to X)_*\mathcal{I}\right)^{\oplus r}$ 也具有性质 $\mathcal{P}$。
由 (2) 推出 $(Z_0 \to X)_*\mathcal{I}$ 具有性质 $\mathcal{P}$。
将此放入 Cohomology of Schemes，引理
\ref{coherent-lemma-property-irreducible} 证明中的适当位置，即得引理。
细节略。
\end{proof}

\begin{lemma}
\label{obsolete-lemma-property-higher-rank}
设 $X$ 是诺特概形。设 $\mathcal{P}$ 是 $X$ 上凝聚层的性质，满足
\begin{enumerate}
\item 对任意凝聚层的短正合列，若其中三个对象中的两个具有性质 $\mathcal{P}$，
则第三个也具有该性质。
\item 若凝聚层的直和具有 $\mathcal{P}$，则该性质对两个直项均成立。
\item 对每个一般点为 $\xi$ 的积分闭子概形 $Z \subset X$，存在凝聚层 $\mathcal{G}$，使得
\begin{enumerate}
\item $\text{Supp}(\mathcal{G}) = Z$；
\item $\mathcal{G}_\xi$ 被 $\mathfrak m_\xi$ 湮灭；
\item $\mathcal{G}$ 具有性质 $\mathcal{P}$。
\end{enumerate}
\end{enumerate}
则 $X$ 上每个凝聚层都具有性质 $\mathcal{P}$。
\end{lemma}

\begin{proof}
这由引理 \ref{obsolete-lemma-property-irreducible-higher-rank} 得出，方式与
Cohomology of Schemes，引理 \ref{coherent-lemma-property} 由
Cohomology of Schemes，引理 \ref{coherent-lemma-property-irreducible} 得出的方式完全相同。
\end{proof}

\begin{lemma}
\label{obsolete-lemma-section-maps-back-into}
设 $X$ 是概形。设 $\mathcal{L}$ 是可逆 $\mathcal{O}_X$-模。
设 $s \in \Gamma(X, \mathcal{L})$ 是截面。设 $\mathcal{F}' \subset \mathcal{F}$ 是拟凝聚
$\mathcal{O}_X$-模。假设
\begin{enumerate}
\item $X$ 拟紧；
\item $\mathcal{F}$ 是有限型；
\item $\mathcal{F}'|_{X_s} = \mathcal{F}|_{X_s}$。
\end{enumerate}
则存在 $n \geq 0$，使得 $\mathcal{F}$ 上乘以 $s^n$ 的映射经由 $\mathcal{F}'$ 分解。
\end{lemma}

\begin{proof}
换言之，我们要证明
$s^n\mathcal{F} \subset
\mathcal{F}' \otimes_{\mathcal{O}_X} \mathcal{L}^{\otimes n}$
对某个 $n \geq 0$ 成立。换言之，商映射
$\mathcal{F} \to \mathcal{F}/\mathcal{F}'$ 在乘以 $s$ 的某个幂后变为零。
这由 Properties，引理
\ref{properties-lemma-section-maps-backwards}.
\end{proof}

\begin{lemma}
\label{obsolete-lemma-bound-degree-in-nbhd-generic-point}
设 $f : X \to Y$ 是概形态射。假设
\begin{enumerate}
\item $X$ 与 $Y$ 是积分概形；
\item $f$ 局部有限型且占优；
\item $f$ 拟紧或分离；
\item $f$ 泛有限，即 Morphisms，引理 \ref{morphisms-lemma-finite-degree} 的
(1)--(5) 中至少有一个成立。
\end{enumerate}
则存在非空开集 $V \subset Y$，使得 $f^{-1}(V) \to V$ 是次数为 $\deg(X/Y)$ 的有限局部自由态射。
特别地，$f^{-1}(V) \to V$ 的纤维次数由 $\deg(X/Y)$ 控制。
\end{lemma}

\begin{proof}
可选取 $V$ 使 $f^{-1}(V) \to V$ 有限。再缩小 $V$，由一般平坦性可设
$f^{-1}(V) \to V$ 平坦且有限表示（Morphisms，命题
\ref{morphisms-proposition-generic-flatness}）。由 Morphisms，引理
\ref{morphisms-lemma-finite-flat}，该态射是有限局部自由的。
由于 $V$ 不可约，该态射具有固定次数。最后一个断言由此以及 Morphisms，引理
\ref{morphisms-lemma-finite-locally-free-universally-bounded} 得出。
\end{proof}







\section{簇的导出范畴}
\label{obsolete-section-equiv}

\noindent
以下若干引理原属簇的导出范畴一章，
但现已不再需要。

\begin{lemma}
\label{obsolete-lemma-preserves-Coh}
设 $k$ 是域。设 $X$ 是 $k$ 上有限型分离且正则的概形。设
$F : D_{perf}(\mathcal{O}_X) \to D_{perf}(\mathcal{O}_X)$ 是 $k$-线性正合函子。
假设对每个满足 $\dim(\text{Supp}(\mathcal{F})) = 0$ 的凝聚
$\mathcal{O}_X$-模 $\mathcal{F}$，存在 $k$-向量空间同构
$$
\Hom_X(\mathcal{F}, M) = \Hom_X(\mathcal{F}, F(M))
$$
且该同构在 $D_{perf}(\mathcal{O}_X)$ 中关于 $M$ 函子性。
则存在定义在 $k$ 上的自同构 $f : X \to X$，它在底层拓扑空间上诱导恒等映射\footnote{这通常迫使 $f$ 为恒等映射；见 Varieties，引理 \ref{varieties-lemma-automorphism}。}
，以及可逆 $\mathcal{O}_X$-模 $\mathcal{L}$，使得 $F$ 与
$F'(M) = f^*M \otimes_{\mathcal{O}_X}^\mathbf{L} \mathcal{L}$ 是兄弟函子。
\end{lemma}

\begin{proof}
由 Derived Categories of Varieties，引理 \ref{equiv-lemma-duality-at-point}，
对每个支集为闭点的凝聚 $\mathcal{O}_X$-模 $\mathcal{F}$，都有同构
$$
H^0(X, M \otimes^\mathbf{L}_{\mathcal{O}_X} \mathcal{F}) =
H^0(X, F(M) \otimes^\mathbf{L}_{\mathcal{O}_X} \mathcal{F})
$$
，且关于 $M$ 函子性。

\medskip\noindent
令 $x \in X$ 为闭点，并将上述结论应用于在 $x$ 处取值为 $\kappa(x)$ 的叠置层
$\mathcal{F} = \mathcal{O}_x$。得到
$$
\dim_{\kappa(x)} 
\text{Tor}^{\mathcal{O}_{X, x}}_p(M_x, \kappa(x)) =
\dim_{\kappa(x)} 
\text{Tor}^{\mathcal{O}_{X, x}}_p(F(M)_x, \kappa(x))
$$
对所有 $p \in \mathbf{Z}$ 成立。特别地，若 $H^i(M) = 0$ 对 $i > 0$ 成立，
则 $H^i(F(M)) = 0$ 对 $i > 0$，这由 Derived Categories of Varieties，引理
\ref{equiv-lemma-orthogonal-point-sheaf}.

\medskip\noindent
若 $\mathcal{E}$ 是秩为 $r$ 的局部自由模，则 $F(\mathcal{E})$ 也是秩为 $r$ 的局部自由模。
这是因为在 $\mathcal{O}_{X, x}$ 上的完美复形 $K$ 若满足
$$
\dim_{\kappa(x)} \text{Tor}^{\mathcal{O}_{X, x}}_i(K, \kappa(x)) =
\left\{
\begin{matrix}
r & \text{若} & i = 0 \\
0 & \text{若} & i \not = 0
\end{matrix}
\right.
$$
则等于置于次数 $0$ 的秩为 $r$ 的自由模。例见 More on Algebra，引理
\ref{more-algebra-lemma-lift-perfect-from-residue-field}.

\medskip\noindent
若 $M$ 的支集包含于闭子概形 $Z \subset X$，则 $F(M)$ 的支集也包含于 $Z$。
这是显然的，因为对 $x \not \in Z$ 有
$M \otimes_{\mathcal{O}_X}^\mathbf{L} \mathcal{O}_x = 0$，于是 $F(M)$ 也满足同样结论，
再由 Derived Categories of Varieties，引理
\ref{equiv-lemma-orthogonal-point-sheaf}.

\medskip\noindent
特别地，$F(\mathcal{O}_x)$ 的支集为 $\{x\}$。令 $i \in \mathbf{Z}$ 为满足
$H^i(\mathcal{O}_x) \not = 0$ 的最小整数。我们知道 $i \leq 0$。
若 $i < 0$，则存在态射
$\mathcal{O}_x[-i] \to F(\mathcal{O}_x)$
，这与所有态射 $\mathcal{O}_x[-i] \to \mathcal{O}_x$ 均为零相矛盾。
因此 $F(\mathcal{O}_x) = \mathcal{H}[0]$，其中 $\mathcal{H}$ 是集中在 $x$ 的叠置层。

\medskip\noindent
设 $\mathcal{G}$ 是满足 $\dim(\text{Supp}(\mathcal{G})) = 0$ 的凝聚 $\mathcal{O}_X$-模。
则存在滤链
$$
0 = \mathcal{G}_0 \subset \mathcal{G}_1 \subset \ldots \subset
\mathcal{G}_n = \mathcal{G}
$$
使得对 $n \geq i \geq 1$，商 $\mathcal{G}_i/\mathcal{G}_{i - 1}$ 同构于某个闭点
$x_i \in X$ 处的 $\mathcal{O}_{x_i}$。于是得到三角形
$$
F(\mathcal{G}_{i - 1}) \to F(\mathcal{G}_i) \to F(\mathcal{O}_{x_i})
$$
利用归纳法可知 $F(\mathcal{G}_i)$ 是置于次数 $0$ 的凝聚层。

\medskip\noindent
设 $\mathcal{G}$ 是凝聚 $\mathcal{O}_X$-模。我们知道对 $i > 0$ 有
$H^i(F(\mathcal{G})) = 0$。为导出矛盾，假设对某个 $i < 0$，
$H^i(F(\mathcal{G}))$ 非零。取具有此性质的最小 $i$，于是存在
$H^i(F(\mathcal{G}))[-i] \to F(\mathcal{G})$ 的 $D_{perf}(\mathcal{O}_X)$ 中态射。
取 $H^i(F(\mathcal{G}))$ 支集中的闭点 $x \in X$。由 More on Algebra，引理
\ref{more-algebra-lemma-kollar-kovacs-pseudo-coherent}
存在 $n > 0$，使得
$$
H^i(F(\mathcal{G}))_x \otimes_{\mathcal{O}_{X, x}}
\mathcal{O}_{X, x}/\mathfrak m_x^n
\longrightarrow
\text{Tor}^{\mathcal{O}_{X, x}}_{-i}(F(\mathcal{G})_x,
\mathcal{O}_{X, x}/\mathfrak m_x^n)
$$
非零。接着取 $m \geq 1$，考虑短正合列
$$
0 \to \mathfrak m_x^m \mathcal{G} \to \mathcal{G} \to
\mathcal{G}/\mathfrak m_x^m\mathcal{G} \to 0
$$
由上面的结论，$F(\mathcal{G}/\mathfrak m_x^m\mathcal{G})$ 是置于次数 $0$ 的层。
因此 $H^i(F(\mathfrak m_x^m \mathcal{G})) \to H^i(F(\mathcal{G}))$ 是同构。
考虑交换图表
$$
\xymatrix{
H^i(F(\mathfrak m_x^m\mathcal{G}))_x \otimes_{\mathcal{O}_{X, x}}
\mathcal{O}_{X, x}/\mathfrak m_x^n \ar[r] \ar[d] &
\text{Tor}^{\mathcal{O}_{X, x}}_{-i}(F(\mathfrak m_x^m\mathcal{G})_x,
\mathcal{O}_{X, x}/\mathfrak m_x^n) \ar[d] \\
H^i(F(\mathcal{G}))_x \otimes_{\mathcal{O}_{X, x}}
\mathcal{O}_{X, x}/\mathfrak m_x^n \ar[r] &
\text{Tor}^{\mathcal{O}_{X, x}}_{-i}(F(\mathcal{G})_x,
\mathcal{O}_{X, x}/\mathfrak m_x^n)
}
$$
由于左侧竖直箭头是同构且底部箭头非零，得到右侧竖直箭头对所有 $m \geq 1$ 均非零。
另一方面，由证明的第一段，该箭头同构于下列箭头
$$
\text{Tor}^{\mathcal{O}_{X, x}}_{-i}(\mathfrak m_x^m\mathcal{G}_x,
\mathcal{O}_{X, x}/\mathfrak m_x^n)
\longrightarrow
\text{Tor}^{\mathcal{O}_{X, x}}_{-i}(\mathcal{G}_x,
\mathcal{O}_{X, x}/\mathfrak m_x^n)
$$
然而由 More on Algebra，引理 \ref{more-algebra-lemma-tor-annihilated}，当 $m \gg n$ 时该箭头为零，
这就是所需的矛盾。

\medskip\noindent
因此 $F$ 保持凝聚模。由 Derived Categories of Varieties，引理
\ref{equiv-lemma-exact-functor-preserving-Coh}
，可知 $F$ 是 Fourier--Mukai 函子 $F'$ 的兄弟函子，其中该函子由凝聚
$\mathcal{O}_{X \times X}$-模 $\mathcal{K}$ 给出，该模经 $\text{pr}_1$ 在 $X$ 上平坦，
经 $\text{pr}_2$ 在 $X$ 上有限。由于 $F(\mathcal{O}_X)$ 是置于次数 $0$ 的可逆
$\mathcal{O}_X$-模 $\mathcal{L}$，我们得到
$$
\mathcal{L} \cong F(\mathcal{O}_X) \cong F'(\mathcal{O}_X) \cong
\text{pr}_{2, *}\mathcal{K}
$$
于是由 Functors and Morphisms，引理 \ref{functors-lemma-pushforward-invertible-pre}，
存在态射 $s : X \to X \times X$，满足 $\text{pr}_2 \circ s = \text{id}_X$ 且
$\mathcal{K} = s_*\mathcal{L}$。令 $f = \text{pr}_1 \circ s$。则有
\begin{align*}
F'(M)
& =
R\text{pr}_{2, *}(L\text{pr}_1^*K \otimes \mathcal{K}) \\
& =
R\text{pr}_{2, *}(L\text{pr}_1^*M \otimes s_*\mathcal{L}) \\
& =
R\text{pr}_{2, *}(Rs_*(Lf^*M \otimes \mathcal{L})) \\
& =
Lf^*M \otimes \mathcal{L}
\end{align*}
第三步使用了 Derived Categories of Schemes，引理
\ref{perfect-lemma-cohomology-base-change}。由于对所有闭点 $x \in X$，模
$F(\mathcal{O}_x)$ 的支集在 $x$，可知 $f$ 在 $X$ 的底层拓扑空间上诱导恒等映射。
仍需证明 $f$ 是同构，这将在下一段完成。

\medskip\noindent
令 $x \in X$ 为闭点。对 $n \geq 1$，记 $\mathcal{O}_{x, n}$ 为在 $x$ 处取值
$\mathcal{O}_{X, x}/\mathfrak m_x^n$ 的叠置层。我们有
$$
\Hom_X(\mathcal{O}_{x, m}, \mathcal{O}_{x, n}) \cong
\Hom_X(\mathcal{O}_{x, m}, F(\mathcal{O}_{x, n})) \cong
\Hom_X(\mathcal{O}_{x, m}, f^*\mathcal{O}_{x, n} \otimes \mathcal{L})
$$
这些同构关于各 $\mathcal{O}_{x, n}$ 之间的 $\mathcal{O}_X$-模同态函子性。
（第一个同构由假设给出，第二个同构因为 $F$ 与 $F'$ 是兄弟函子。）
当 $m \geq n$ 时，通过对 $\mathcal{O}_{x, n}$ 的作用，有
$\mathcal{O}_{X, x}/\mathfrak m^n = \Hom_X(\mathcal{O}_{x, m}, \mathcal{O}_{x, n})$。
因此 $f^\sharp : \mathcal{O}_{X, x}/\mathfrak m_x^n \to
\mathcal{O}_{X, x}/\mathfrak m_x^n$ 对所有 $n$ 均为双射。
所以 $f$ 在闭点处的完备局部环上诱导同构，因而是 \'etale
（\'Etale Morphisms，引理 \ref{etale-lemma-characterize-etale-completions}）。
考察闭点可知，$\Delta_f : X \to X \times_{f, X, f} X$（由于 $f$ 是 \'etale，这是开浸入）
是双射，因而是同构。因此 $f$ 是单态射。最后由 Descent，引理
\ref{descent-lemma-flat-surjective-quasi-compact-monomorphism-isomorphism}
知它是开浸入，从而 $f$ 是同构。
\end{proof}










\section{正则真情形下的可表性}
\label{obsolete-section-regular-proper}

\noindent
本节已废弃，因为我们改进了 Derived Categories of Varieties，
定理 \ref{equiv-theorem-bondal-van-den-bergh}，使其适用于域上的所有真概形
（此前我们只对域上的射影概形证明了该结论）。

\begin{lemma}
\label{obsolete-lemma-trace-map}
\begin{reference}
此处的证明沿用
\cite[Remark 3.4]{MS}
\end{reference}
设 $f : X' \to X$ 是积分诺特概形之间的真双有理态射，且 $X$ 正则。
映射 $\mathcal{O}_X \to Rf_*\mathcal{O}_{X'}$ 在 $D(\mathcal{O}_X)$ 中典范分裂。
\end{lemma}

\begin{proof}
在 $D(\mathcal{O}_X)$ 中置 $E = Rf_*\mathcal{O}_{X'}$。由 Derived Categories of Schemes，引理
\ref{perfect-lemma-direct-image-coherent}，有 $E$ 属于 $D^b_{\textit{Coh}}(\mathcal{O}_X)$。
由 Derived Categories of Schemes，引理 \ref{perfect-lemma-perfect-on-regular}，可知 $E$ 是
$D(\mathcal{O}_X)$ 中的完美对象。
由于 $\mathcal{O}_{X'}$ 是代数层，由 Cohomology，注记
\ref{cohomology-remark-cup-product} 有相对杯积
$\mu : E \otimes_{\mathcal{O}_X}^\mathbf{L} E \to E$。
令 $\sigma : E \otimes E^\vee \to E^\vee \otimes E$ 为对称幺半范畴
$D(\mathcal{O}_X)$ 中的交换约束（Cohomology，引理
\ref{cohomology-lemma-symmetric-monoidal-derived}）。记
$\eta : \mathcal{O}_X \to E \otimes E^\vee$ 与
$\epsilon : E^\vee \otimes E \to \mathcal{O}_X$ 为 Cohomology，例
\ref{cohomology-example-dual-derived} 中构造的映射。于是可考虑映射
$$
E \xrightarrow{\eta \otimes 1} E \otimes E^\vee \otimes E
\xrightarrow{\sigma \otimes 1} E^\vee \otimes E \otimes E
\xrightarrow{1 \otimes \mu} E^\vee \otimes E
\xrightarrow{\epsilon} \mathcal{O}_X
$$
我们断言该映射是引理陈述中映射的单侧逆。只需证明复合
$\mathcal{O}_X \to \mathcal{O}_X$ 是恒等映射即可。可以在 $X$ 的一般点处
（或在 $f$ 为同构的 $X$ 的某个开子概形上）进行验证。此时
$E = \mathcal{O}_X$，而 $\mu$ 是通常的乘法映射，结论显然。
\end{proof}

\begin{lemma}
\label{obsolete-lemma-characterize-dbcoh-proper-regular}
设 $X$ 是域 $k$ 上的正则真概形。设
$K \in \Ob(D_\QCoh(\mathcal{O}_X))$。以下条件等价：
\begin{enumerate}
\item $K \in D^b_{\textit{Coh}}(\mathcal{O}_X) = D_{perf}(\mathcal{O}_X)$；
\item 对 $D(\mathcal{O}_X)$ 中所有完美对象 $E$，有
$\sum_{i \in \mathbf{Z}} \dim_k \Ext^i_X(E, K) < \infty$。
\end{enumerate}
\end{lemma}

\begin{proof}
条件 (1) 中的等式由 Derived Categories of Schemes，引理
\ref{perfect-lemma-perfect-on-regular} 给出。(1) $\Rightarrow$ (2) 由
Derived Categories of Varieties，引理 \ref{equiv-lemma-finiteness} 得出。
(2) $\Rightarrow$ (1) 由 More on Morphisms，引理
\ref{more-morphisms-lemma-characterize-relatively-perfect}.
\end{proof}

\begin{lemma}
\label{obsolete-lemma-bondal-van-den-bergh}
设 $X$ 是域 $k$ 上的正则真概形。
\begin{enumerate}
\item 设 $F : D_{perf}(\mathcal{O}_X)^{opp} \to \text{Vect}_k$ 是 $k$-线性上同调函子，满足
$$
\sum\nolimits_{n \in \mathbf{Z}} \dim_k F(E[n]) < \infty
$$
对所有 $E \in D_{perf}(\mathcal{O}_X)$ 成立。则 $F$ 同构于形如
$E \mapsto \Hom_X(E, K)$ 的函子，其中 $K \in D_{perf}(\mathcal{O}_X)$。
\item 设 $G : D_{perf}(\mathcal{O}_X) \to \text{Vect}_k$ 是 $k$-线性同调函子，满足
$$
\sum\nolimits_{n \in \mathbf{Z}} \dim_k G(E[n]) < \infty
$$
对所有 $E \in D_{perf}(\mathcal{O}_X)$ 成立。则 $G$ 同构于形如
$E \mapsto \Hom_X(K, E)$ 的函子，其中 $K \in D_{perf}(\mathcal{O}_X)$。
\end{enumerate}
\end{lemma}

\begin{proof}
这由 Derived Categories of Varieties，定理
\ref{equiv-theorem-bondal-van-den-bergh} 与
引理 \ref{equiv-lemma-homological-representable} 得出。
下面也给出另一种证明。

\medskip\noindent
证明 (1)。导出范畴 $D_\QCoh(\mathcal{O}_X)$ 有直和，是紧生成的，且
$D_{perf}(\mathcal{O}_X)$ 是紧对象的全子范畴；见
Derived Categories of Schemes，引理
\ref{perfect-lemma-quasi-coherence-direct-sums},
定理 \ref{perfect-theorem-bondal-van-den-Bergh} 与
命题 \ref{perfect-proposition-compact-is-perfect}。
由 Derived Categories of Varieties，引理 \ref{equiv-lemma-van-den-bergh}，可假设
$F(E) = \Hom_X(E, K)$，其中 $K \in \Ob(D_\QCoh(\mathcal{O}_X))$。
于是由引理 \ref{obsolete-lemma-characterize-dbcoh-proper-regular}，$K$ 属于
$D^b_{\textit{Coh}}(\mathcal{O}_X)$。

\medskip\noindent
证明 (2)。考虑 $D_{perf}(\mathcal{O}_X)$ 上的逆变函子 $E \mapsto E^\vee$；见
Cohomology，引理 \ref{cohomology-lemma-dual-perfect-complex}。
该函子是 $D_{perf}(\mathcal{O}_X)$ 的正合反自等价。
因此可将 (1) 应用于函子 $F(E) = G(E^\vee)$，得到
$K \in D_{perf}(\mathcal{O}_X)$，满足 $G(E^\vee) = \Hom_X(E, K)$。
于是 $G(E) = \Hom_X(E^\vee, K) = \Hom_X(K^\vee, E)$，
故取 $K^\vee$ 即可。
\end{proof}









\section{商函子}
\label{obsolete-section-quotients}

\begin{lemma}
\label{obsolete-lemma-factors-through-quotient}
设 $S = \Spec(R)$ 是仿射概形。设 $X$ 是 $S$ 上的代数空间。
设 $q_i : \mathcal{F} \to \mathcal{Q}_i$（$i = 1, 2$）是拟凝聚
$\mathcal{O}_X$-模的满射。假设 $\mathcal{Q}_1$ 在 $S$ 上平坦。
设 $T \to S$ 是概形的拟紧态射，并且存在如下分解
$$
\xymatrix{
& \mathcal{F}_T \ar[rd]^{q_{2, T}} \ar[ld]_{q_{1, T}} \\
\mathcal{Q}_{1, T} & & \mathcal{Q}_{2, T} \ar@{..>}[ll]
}
$$
则存在闭子概形 $Z \subset S$，使得 (a) $T \to S$ 经由 $Z$ 分解，且 (b)
$q_{1, Z}$ 经由 $q_{2, Z}$ 分解。若 $\Ker(q_2)$ 是有限型
$\mathcal{O}_X$-模且 $X$ 拟紧，则可取 $Z \to S$ 为有限表示。
\end{lemma}

\begin{proof}
将 Flatness on Spaces，引理 \ref{spaces-flat-lemma-F-zero-somewhat-closed}
应用于映射 $\Ker(q_2) \to \mathcal{Q}_1$。
\end{proof}







\section{空间与 fpqc 覆盖}
\label{obsolete-section-fpqc}

\noindent
本节材料因 Gabber 的论证而废弃；该论证表明代数空间相对于 fpqc 覆盖满足层条件。
请参阅
Properties of Spaces，小节 \ref{spaces-properties-section-fpqc}。

\begin{lemma}
\label{obsolete-lemma-separated-fpqc}
设 $S$ 是概形。设 $X$ 是 $S$ 上的代数空间。设
$\{f_i : T_i \to T\}_{i \in I}$ 是 $S$ 上概形的 fpqc 覆盖。
则映射
$$
\Mor_S(T, X)
\longrightarrow
\prod\nolimits_{i \in I} \Mor_S(T_i, X)
$$
是单射。
\end{lemma}

\begin{proof}
立即由 Properties of Spaces，命题
\ref{spaces-properties-proposition-sheaf-fpqc}.
\end{proof}

\begin{lemma}
\label{obsolete-lemma-sheaf-fpqc-open-covering}
设 $S$ 是概形。设 $X$ 是 $S$ 上的代数空间。设
$X = \bigcup_{j \in J} X_j$ 是 Zariski 覆盖，见
Spaces，定义 \ref{spaces-definition-Zariski-open-covering}。
若每个 $X_j$ 对 fpqc 拓扑满足层性质，则 $X$ 对 fpqc 拓扑也满足层性质。
\end{lemma}

\begin{proof}
这是因为所有代数空间对 fpqc 拓扑都满足层性质；见
Properties of Spaces，命题
\ref{spaces-properties-proposition-sheaf-fpqc}.
\end{proof}

\begin{lemma}
\label{obsolete-lemma-sheaf-fpqc-quasi-separated}
设 $S$ 是概形。设 $X$ 是 $S$ 上的代数空间。若 $X$ 在 Zariski 局部上相对于 $S$
拟分离，则 $X$ 对 fpqc 拓扑满足层条件。
\end{lemma}

\begin{proof}
立即由一般的 Properties of Spaces，命题
\ref{spaces-properties-proposition-sheaf-fpqc}.
\end{proof}

\begin{remark}
\label{obsolete-remark-proof-works-when}
本注记原用于讨论引理 \ref{obsolete-lemma-sheaf-fpqc-quasi-separated}（2009 年 12 月 18 日）
的原证明在多大程度上可以推广。
\end{remark}






\section{非常良好的代数空间}
\label{obsolete-section-very-reasonable}

\noindent
略已废弃的材料。

\begin{lemma}
\label{obsolete-lemma-reasonable-kolmogorov}
设 $S$ 是概形。设 $X$ 是 $S$ 上的良好代数空间。
则 $|X|$ 是 Kolmogorov 空间（见 Topology，定义
\ref{topology-definition-generic-point}）。
\end{lemma}

\begin{proof}
由定义与 Decent Spaces，引理 \ref{decent-spaces-lemma-kolmogorov} 得出。
\end{proof}

\noindent
本节其余部分讨论非常良好的代数空间。若存在概形 $U$ 及满射、\'etale、拟紧
态射 $U \to X$，则 $X$ 非常良好；见 Decent Spaces，引理
\ref{decent-spaces-lemma-characterize-very-reasonable}。

\begin{lemma}
\label{obsolete-lemma-scheme-very-reasonable}
概形是非常良好的。
\end{lemma}

\begin{proof}
这是因为恒等映射是拟紧满射 \'etale 态射。
\end{proof}

\begin{lemma}
\label{obsolete-lemma-very-reasonable-Zariski-local}
设 $S$ 是概形。设 $X$ 是 $S$ 上的代数空间。
若存在 Zariski 开覆盖 $X = \bigcup X_i$，使每个 $X_i$ 都非常良好，
则 $X$ 非常良好。
\end{lemma}

\begin{proof}
这是 Decent Spaces，引理 \ref{decent-spaces-lemma-properties-local} 的情形 $(\epsilon)$。
\end{proof}

\begin{lemma}
\label{obsolete-lemma-quasi-separated-very-reasonable}
Zariski 局部拟分离的代数空间非常良好。特别地，任何拟分离代数空间都非常良好。
\end{lemma}

\begin{proof}
这是 Decent Spaces，引理 \ref{decent-spaces-lemma-bounded-fibres} 的一个推论。
\end{proof}

\begin{lemma}
\label{obsolete-lemma-representable-very-reasonable}
设 $S$ 是概形。设 $X$、$Y$ 是 $S$ 上的代数空间。
设 $Y \to X$ 是可表态射。若 $X$ 非常良好，则 $Y$ 也非常良好。
\end{lemma}

\begin{proof}
这是 Decent Spaces，引理 \ref{decent-spaces-lemma-representable-properties} 的情形 $(\epsilon)$。
\end{proof}

\begin{remark}
\label{obsolete-remark-very-reasonable-Zariski-locally-quasi-separated}
非常良好的代数空间构成一个严格大于 Zariski 局部拟分离代数空间的集合。
考虑形如 $X = [U/G]$ 的代数空间（见 Spaces，定义
\ref{spaces-definition-quotient}），其中 $G$ 是有限群，在非拟分离概形 $U$ 上无不动点地作用。
在此情形，$U \times_X U = U \times G$，且显然到 $U$ 的两个投影都是拟紧的，
故 $X$ 非常良好。另一方面，对角映射
$U \times_X U \to U \times U$ 不是拟紧的，因此该代数空间不是拟分离的。
现在取特征为 $\not = 2$ 的域 $k$ 上的无限仿射空间 $U$，并将零点加倍；见
Schemes，例 \ref{schemes-example-not-quasi-separated}。
令 $0_1, 0_2$ 为 $U$ 的两个零点。令 $G = \{+1, -1\}$，令 $-1$ 在所有坐标上按
$-1$ 作用，并交换 $0_1$ 与 $0_2$。则 $[U/G]$ 非常良好，但不是 Zariski 局部拟分离的
（细节略）。
\end{remark}

\noindent
警告：使用以下引理时须谨慎，因为其中的概形 $U_i$ 不一定分离，甚至不一定拟分离。

\begin{lemma}
\label{obsolete-lemma-very-reasonable-quasi-compact-pieces}
设 $S$ 是概形。设 $X$ 是 $S$ 上的非常良好代数空间。
存在一族概形 $U_i$ 及态射 $U_i \to X$，使得
\begin{enumerate}
\item 每个 $U_i$ 是拟紧概形；
\item 每个 $U_i \to X$ 是 \'etale；
\item 两个投影 $U_i \times_X U_i \to U_i$ 都是拟紧的；
\item 态射 $\coprod U_i \to X$ 是满射（且为 \'etale）。
\end{enumerate}
\end{lemma}

\begin{proof}
Decent Spaces，定义 \ref{decent-spaces-definition-very-reasonable} 表明存在满足 (2)、(3)、(4) 的
$U_i \to X$。固定 $i$，置 $R_i = U_i \times_X U_i$，记投影为
$s, t : R_i \to U_i$。对任意仿射开集 $W \subset U_i$，开集
$W' = t(s^{-1}(W)) \subset U_i$ 是拟紧的 $R_i$-不变开集（见
Groupoids，引理 \ref{groupoids-lemma-constructing-invariant-opens}）。
因此 $W'$ 是拟紧概形，$W' \to X$ 是 \'etale，且
$W' \times_X W' = s^{-1}(W') = t^{-1}(W')$，所以两个投影
$W' \times_X W' \to W'$ 都是拟紧的。令 $W \subset U_i$ 遍历 $U_i$ 的仿射开覆盖中的成员，
所得族 $W' \to X$ 即满足要求。
\end{proof}




\section{关于代数空间的已废弃引理}
\label{obsolete-section-obsolete-on-spaces}

\noindent
看似多余或不再用于正文的引理。

\begin{lemma}
\label{obsolete-lemma-vanishing-surjective}
在 Cohomology of Spaces，情形 \ref{spaces-cohomology-situation-vanishing} 中，
态射 $p : X \to \Spec(A)$ 是满射。
\end{lemma}

\begin{proof}
本引理原用于 Cohomology of Spaces，命题
\ref{spaces-cohomology-proposition-vanishing-affine} 的证明，
但现在它是该命题的推论。
\end{proof}

\begin{lemma}
\label{obsolete-lemma-vanishing-universally-closed}
在 Cohomology of Spaces，情形 \ref{spaces-cohomology-situation-vanishing} 中，
态射 $p : X \to \Spec(A)$ 是泛闭的。
\end{lemma}

\begin{proof}
本引理原用于 Cohomology of Spaces，命题
\ref{spaces-cohomology-proposition-vanishing-affine} 的证明，
但现在它是该命题的推论。
\end{proof}

\begin{remark}
\label{obsolete-remark-equation-first}
此标签原指向 Formal Spaces，引理 \ref{formal-spaces-lemma-iff-adic-star} 证明中的一个方程。
\end{remark}

\begin{remark}
\label{obsolete-remark-equation-second}
此标签原指向 Formal Spaces，引理 \ref{formal-spaces-lemma-iff-adic-star} 证明中的一个方程。
\end{remark}





\section{关于代数叠的已废弃引理}
\label{obsolete-section-obsolete-on-stacks}

\noindent
看似多余或不再用于正文的引理。

\begin{lemma}
\label{obsolete-lemma-infinite-sequence}
设 $S$ 是局部诺特概形。设 $\mathcal{X}$ 是在 $(\Sch/S)_{fppf}$ 上纤维化于群胚的范畴，
具有 (RS*)。设 $x$ 是 $\mathcal{X}$ 中位于 $S$ 上有限型仿射概形 $U$ 上的对象。
设 $u_n \in U$（$n \geq 1$）是两两不同的有限型点，并且对所有 $n$，$x$ 在 $u_n$ 处都不是 versal 的。
将 $u_n$ 换为子列后，存在态射
$$
x \to x_1 \to x_2 \to \ldots
\quad\text{于 }\mathcal{X}\text{ 位于下列对象之上： }\quad
U \to U_1 \to U_2 \to \ldots
$$
它们定义在 $S$ 上，并满足
\begin{enumerate}
\item 对每个 $n$，态射 $U \to U_n$ 是一阶加厚；
\item 对每个 $n$，有短正合列
$$
0 \to \kappa(u_n) \to \mathcal{O}_{U_n} \to \mathcal{O}_{U_{n - 1}} \to 0
$$
其中 $n = 1$ 时 $U_0 = U$；
\item 对每个 $n$，{\bf 不}存在由 $u_n$ 在其中的开邻域 $W \subset U_n$ 与态射
$\alpha : x_n|_W \to x$ 构成的对 $(W, \alpha)$，使得复合
$$
x|_{U \cap W} \xrightarrow{\text{下列对象的限制： }x \to x_n}
x_n|_W \xrightarrow{\alpha} x
$$
是典范态射 $x|_{U \cap W} \to x$。
\end{enumerate}
\end{lemma}

\begin{proof}
本引理原用于 versal 性开放性判据的证明
（Artin's Axioms，引理 \ref{artin-lemma-SGE-implies-openness-versality}），但已被
Artin's Axioms，引理 \ref{artin-lemma-infinite-sequence-pre} 取代，后者可直接推出本引理。
具体地，将 $u_n$（$n \geq 1$）换为子列后，可假设这些点之间没有专化关系；见
Properties，引理 \ref{properties-lemma-thin-infinite-sequence}。然后应用
Artin's Axioms，引理 \ref{artin-lemma-infinite-sequence-pre}，即可完成证明。
\end{proof}







\section{概形余切复形的变体}
\label{obsolete-section-cotangent-schemes-variant}

\noindent
本节给出概形态射余切复形的另一种构造。目前本节位于已废弃章节中，
因为在应用中可以使用 Cotangent，小节
\ref{cotangent-section-cotangent-schemes-variant} 所讨论的较容易版本。

\medskip\noindent
设 $f : X \to Y$ 是概形态射。令 $\mathcal{C}_{X/Y}$ 为如下交换图表为对象的范畴
\begin{equation}
\label{obsolete-equation-object}
\vcenter{
\xymatrix{
X \ar[d] & U \ar[l] \ar[d] \ar[r]_i & A \ar[ld] \\
Y & V \ar[l]
}
}
\end{equation}
，其中各项为概形且
\begin{enumerate}
\item $U$ 是 $X$ 的开子概形；
\item $V$ 是 $Y$ 的开子概形；
\item 在 $V$ 上存在同构 $A = V \times \Spec(P)$，其中 $P$ 是
  $\mathbf{Z}$ 上的多项式代数（变量取自某个集合）。
\end{enumerate}
换言之，$A$ 是 $V$ 上的（无限维）仿射空间。态射由交换图表给出。

\medskip\noindent
{\bf 记号。} $\mathcal{C}_{X/Y}$ 的对象，即图表
(\ref{obsolete-equation-object})，常记为 $U \to A$，其中约定：(a) $U$ 是 $X$ 的开子概形；(b)
$U \to A$ 是在 $Y$ 上的态射；(c) 结构态射 $A \to Y$ 的像是开集 $V \subset Y$；(d)
$A \to V$ 是仿射空间。我们写 $U \to A/V$ 表示 $V \subset Y$ 是 $A \to Y$ 的像。
记 $X_{Zar}$ 为 $X$ 的小 Zariski site。存在遗忘函子
$$
\mathcal{C}_{X/Y} \to X_{Zar},\ (U \to A) \mapsto U
\quad\text{且}\quad
\mathcal{C}_{X/Y} \mapsto Y_{Zar},\ (U \to A/V) \mapsto V.
$$

\begin{lemma}
\label{obsolete-lemma-category-fibred}
设 $X \to Y$ 是概形态射。
\begin{enumerate}
\item 范畴 $\mathcal{C}_{X/Y}$ 在 $X_{Zar}$ 上纤维化；
\item 范畴 $\mathcal{C}_{X/Y}$ 在 $Y_{Zar}$ 上纤维化；
\item 范畴 $\mathcal{C}_{X/Y}$ 在如下二元组 $(U, V)$ 的范畴上纤维化：
  $U \subset X$、$V \subset Y$ 为开集且 $f(U) \subset V$。
\end{enumerate}
\end{lemma}

\begin{proof}
对 (1)。给定 $\mathcal{C}_{X/Y}$ 的对象 $U \to A$ 及 $X_{Zar}$ 中的态射
$U' \to U$，考虑 $\mathcal{C}_{X/Y}$ 的对象 $i' : U' \to A$，其中 $i'$ 是
$i$ 与 $U' \to U$ 的复合。$\mathcal{C}_{X/Y}$ 中的态射
$(U' \to A) \to (U \to A)$ 在 $X_{Zar}$ 上强笛卡尔。

\medskip\noindent
对 (2)。给定对象 $U \to A/V$ 与 $V' \to V$，置
$U' = U \cap f^{-1}(V')$、$A' = V' \times_V A$，即可得到在 $V' \to V$ 上强笛卡尔的
态射 $(U' \to A') \to (U \to A)$。

\medskip\noindent
对 (3)。记 (3) 中的范畴为 $(X/Y)_{Zar}$。给定 $U \to A/V$ 及
$(X/Y)_{Zar}$ 中的态射 $(U', V') \to (U, V)$，考虑
$A' = V' \times_V A$。则态射 $(U' \to A'/V') \to (U \to A/V)$
在 $\mathcal{C}_{X/Y}$ 中关于 $(X/Y)_{Zar}$ 强笛卡尔。
\end{proof}

\noindent
使用从 $X_{Zar}$ 继承的拓扑（见 Stacks，小节 \ref{stacks-section-topology}），
在 $\mathcal{C}_{X/Y}$ 上得到拓扑 $\tau_X$。除非另有说明，我们考虑的都是
$\mathcal{C}_{X/Y}$ 上的这个拓扑。准确地说，一族态射
$\{(U_i \to A_i) \to (U \to A)\}$ 是 $\mathcal{C}_{X/Y}$ 的覆盖，当且仅当
\begin{enumerate}
\item $U = \bigcup U_i$；
\item 对所有 $i$，$A_i = A$。
\end{enumerate}
若在 (2) 中允许 $A_i \cong A$，得到相同的层集合。函子 $u$ 定义拓扑斯态射
$\pi : \Sh(\mathcal{C}_{X/Y}) \to \Sh(X_{Zar})$.

\medskip\noindent
site $\mathcal{C}_{X/Y}$ 带有若干环层。
\begin{enumerate}
\item 由规则
$(U \to A) \mapsto \mathcal{O}(A)$.
给出的层 $\mathcal{O}$；
\item 由规则 $(U \to A) \mapsto \mathcal{O}(U)$ 给出的层
$\underline{\mathcal{O}}_X = \pi^{-1}\mathcal{O}_X$；
\item 由规则
$(U \to A/V) \mapsto \mathcal{O}(V)$.
给出的层 $\underline{\mathcal{O}}_Y$。
\end{enumerate}
得到环化拓扑斯的态射
\begin{equation}
\label{obsolete-equation-pi-schemes}
\vcenter{
\xymatrix{
(\Sh(\mathcal{C}_{X/Y}), \underline{\mathcal{O}}_X) \ar[r]_i \ar[d]_\pi &
(\Sh(\mathcal{C}_{X/Y}), \mathcal{O}) \\
(\Sh(X_{Zar}), \mathcal{O}_X)
}
}
\end{equation}
态射 $i$ 在底层拓扑斯上是恒等态射，且
$i^\sharp : \mathcal{O} \to \underline{\mathcal{O}}_X$ 是显然的映射。
映射 $\pi$ 是 Cohomology on Sites，情形
\ref{sites-cohomology-situation-fibred-category} 的特殊情形。
下列导出函子将在后文发挥重要作用：
$
Li^* : D(\mathcal{O}) \longrightarrow D(\underline{\mathcal{O}}_X)
$
是 $Ri_* = i_* : D(\underline{\mathcal{O}}_X) \to D(\mathcal{O})$ 的左伴随，以及
$
L\pi_! : D(\underline{\mathcal{O}}_X) \longrightarrow D(\mathcal{O}_X)
$
是下列函子的左伴随：
$\pi^* = \pi^{-1} : D(\mathcal{O}_X) \to D(\underline{\mathcal{O}}_X)$.

\begin{remark}
\label{obsolete-remark-different-topologies}
取从 $Y_{Zar}$ 继承的拓扑，可在 $\mathcal{C}_{X/Y}$ 上得到第二个拓扑 $\tau_Y$。
还有第三个拓扑 $\tau_{X \to Y}$：一族态射
$\{(U_i \to A_i) \to (U \to A)\}$ 是覆盖，当且仅当
$U = \bigcup U_i$、$V = \bigcup V_i$ 且 $A_i \cong V_i \times_V A$。
这是从 site $(X/Y)_{Zar}$ 上的拓扑继承而来的；其底层范畴是引理
\ref{obsolete-lemma-category-fibred} 第 (3) 部分中的二元组 $(U, V)$ 的范畴。
$(X/Y)_{Zar}$ 的覆盖是满足 $U = \bigcup U_i$ 与 $V = \bigcup V_i$ 的族
$\{(U_i, V_i) \to (U, V)\}$。存在拓扑斯态射
$$
\xymatrix{
\Sh(\mathcal{C}_{X/Y})
= \Sh(\mathcal{C}_{X/Y}, \tau_X) &
\Sh(\mathcal{C}_{X/Y}, \tau_{X \to Y}) \ar[l] \ar[r] &
\Sh(\mathcal{C}_{X/Y}, \tau_Y)
}
$$
（回忆 $\tau_X$ 是我们的“默认”拓扑）。这些箭头的拉回函子是层化，而推出在底层预层上是恒等函子。
拓扑斯图表
$$
\xymatrix{
\Sh(X_{Zar}) \ar[d]^f & \Sh(\mathcal{C}_{X/Y}) \ar[l]^\pi &
\Sh(\mathcal{C}_{X/Y}, \tau_{X \to Y}) \ar[l] \ar[d] \\
\Sh(Y_{Zar}) & & \Sh(\mathcal{C}_{X/Y}, \tau_Y) \ar[ll]
}
$$
{\bf 不}交换。事实上，$Y$ 上非零阿贝尔层的拉回是
$(\mathcal{C}_{X/Y}, \tau_{X \to Y})$ 上的非零阿贝尔层，
但确实可以找到该层在 $X$ 上拉回为零的例子。注意，$Y_{Zar}$ 上任意预层
$\mathcal{F}$ 都通过把集合 $\mathcal{F}(V)$ 赋给 $(U \to A/V)$ 的规则，
给出 $\mathcal{C}_{Y/X}$ 上的层 $\underline{\mathcal{F}}$。
即使 $\mathcal{F}$ 恰为层，一般也不成立
$\underline{\mathcal{F}} = \pi^{-1}f^{-1}\mathcal{F}$。这与上图不交换有关：
经下方水平箭头与右侧竖直箭头，可把 $\underline{\mathcal{F}}$ 描述为
$\mathcal{F}$ 到 $\Sh(\mathcal{C}_{X/Y}, \tau_{X \to Y})$ 的拉回之推出。
层 $\underline{\mathcal{O}}_Y$ 即为一例。
但确实存在映射 $\underline{\mathcal{F}} \to \pi^{-1}f^{-1}\mathcal{F}$，
它在应用 $\pi_!$ 后（如后文所见）变成同构。
\end{remark}










\section{平坦模的形变与障碍}
\label{obsolete-section-modules}

\noindent
本节概述环 $\Lambda$ 上任意代数空间 $X$ 的凝聚层叠之形变理论构造。
由于 Quot，小节 \ref{quot-section-not-flat} 中已有改进的讨论，本材料已废弃。

\medskip\noindent
设定如下。假设给定
\begin{enumerate}
\item 环 $\Lambda$；
\item $\Lambda$ 上的代数空间 $X$；
\item $\Lambda$-代数 $A$，置 $X_A = X \times_{\Spec(\Lambda)} \Spec(A)$；
\item 在 $A$ 上平坦的有限表示 $\mathcal{O}_{X_A}$-模 $\mathcal{F}$。
\end{enumerate}
在这种情形下，考虑所有可能的满射
$$
0 \to I \to A' \to A \to 0
$$
其中 $A'$ 是 $\Lambda$-代数，其核 $I$ 是 $A'$ 中平方为零的理想。
给定 $A'$，得到 $\Spec(\Lambda)$ 上代数空间的一阶加厚 $X_A \to X_{A'}$。
对每个这样的加厚，考虑把 $\mathcal{F}$ 提升为 $X_{A'}$ 上在 $A'$ 上平坦的有限表示模
$\mathcal{F}'$ 的问题。我们希望在此设定下复现 Deformation Theory，引理
\ref{defos-lemma-flat-ringed-topoi} 的结果。

\medskip\noindent
更准确地，令 $\textit{Lift}(\mathcal{F}, A')$ 表示如下二元组
$(\mathcal{F}', \alpha)$ 的范畴：$\mathcal{F}'$ 是 $X_{A'}$ 上在 $A'$ 上平坦的有限表示模，
而 $\alpha : \mathcal{F}'|_{X_A} \to \mathcal{F}$ 是同构。
态射 $(\mathcal{F}'_1, \alpha_1) \to (\mathcal{F}'_2, \alpha_2)$ 是与
$\alpha_1$、$\alpha_2$ 相容的同构 $\mathcal{F}'_1 \to \mathcal{F}'_2$。
$\textit{Lift}(\mathcal{F}, A')$ 的同构类集合记为 $\text{Lift}(\mathcal{F}, A')$。

\medskip\noindent
设 $\mathcal{G}$ 是 $X_\etale$ 上在 $A$ 上平坦的
$\mathcal{O}_X \otimes_\Lambda A$-模层。定义二元组
$(\mathcal{G}', \beta)$ 的范畴 $\textit{Lift}(\mathcal{G}, A')$，其中 $\mathcal{G}'$ 是
在 $A'$ 上平坦的 $\mathcal{O}_X \otimes_\Lambda A'$-模层，而 $\beta$ 是同构
$\mathcal{G}' \otimes_{A'} A \to \mathcal{G}$。

\begin{lemma}
\label{obsolete-lemma-equivalence}
记号与假设如上。记投影为 $p : X_A \to X$。给定 $A'$，记投影为
$p' : X_{A'} \to X$。函子 $p'_*$ 在下列范畴之间诱导等价：
\begin{enumerate}
\item 范畴 $\textit{Lift}(\mathcal{F}, A')$；
\item 范畴 $\textit{Lift}(p_*\mathcal{F}, A')$。
\end{enumerate}
\end{lemma}

\begin{proof}
FIXME.
\end{proof}

\noindent
设 $\mathcal{H}$ 是 $\mathcal{C}_{X/\Lambda}$ 上在 $A$ 上平坦的
$\mathcal{O} \otimes_\Lambda A$-模层。定义范畴
$\textit{Lift}_\mathcal{O}(\mathcal{H}, A')$，其对象是二元组
$(\mathcal{H}', \gamma)$，其中 $\mathcal{H}'$ 是在 $A'$ 上平坦的
$\mathcal{O} \otimes_\Lambda A'$-模层，而
$\gamma : \mathcal{H}' \otimes_A A' \to \mathcal{H}$ 是
$\mathcal{O} \otimes_\Lambda A$-模的同构。

\medskip\noindent
设 $\mathcal{G}$ 是 $X_\etale$ 上在 $A$ 上平坦的
$\mathcal{O}_X \otimes_\Lambda A$-模层。考虑 Cotangent，方程
(\ref{cotangent-equation-pi-spaces}) 中的态射 $i$ 与 $\pi$。
记 $\underline{\mathcal{G}} = \pi^{-1}(\mathcal{G})$。它由规则
$(U \to \mathbf{A}) \mapsto \mathcal{G}(U)$ 直接给出，因而是
$\underline{\mathcal{O}}_X \otimes_\Lambda A$-模层。
以 $i_*\underline{\mathcal{G}}$ 表示同一个层，但将其视为
$\mathcal{O} \otimes_\Lambda A$-模层。

\begin{lemma}
\label{obsolete-lemma-second-equivalence}
记号与假设同上。函子 $\pi_!$ 在下列两个范畴之间诱导范畴等价：
\begin{enumerate}
\item 范畴 $\textit{Lift}_\mathcal{O}(i_*\underline{\mathcal{G}}, A')$；
\item 范畴 $\textit{Lift}(\mathcal{G}, A')$。
\end{enumerate}
\end{lemma}

\begin{proof}
FIXME.
\end{proof}

\begin{lemma}
\label{obsolete-lemma-second-equivalence-obs}
记号与假设同引理 \ref{obsolete-lemma-second-equivalence}。考虑
$$
L = L(\Lambda, X, A, \mathcal{G}) = L\pi_!(Li^*(i_*(\underline{\mathcal{G}})))
$$
这一 $D(\mathcal{O}_X \otimes_\Lambda A)$ 中的对象。给定核为平方零理想
$I$ 的 $\Lambda$-代数满射 $A' \to A$，则有：
\begin{enumerate}
\item 范畴 $\textit{Lift}(\mathcal{G}, A')$ 非空，当且仅当某个类
$\xi \in \Ext^2_{\mathcal{O}_X \otimes A}(L, \mathcal{G} \otimes_A I)$
为零。
\item 若 $\textit{Lift}(\mathcal{G}, A')$ 非空，则
$\text{Lift}(\mathcal{G}, A')$ 是
$\Ext^1_{\mathcal{O}_X \otimes A}(L, \mathcal{G} \otimes_A I)$
作用下的主齐性集。
\item 给定一个提升 $\mathcal{G}'$，其拉回为
$\text{id}_\mathcal{G}$ 的 $\mathcal{G}'$ 自同构所成集合典范同构于
$\Ext^0_{\mathcal{O}_X \otimes A}(L, \mathcal{G} \otimes_A I)$.
\end{enumerate}
\end{lemma}

\begin{proof}
FIXME.
\end{proof}

\noindent
最后，将以上各项汇总如下。

\begin{proposition}
\label{obsolete-proposition-conclusion}
设 $\Lambda$、$X$、$A$、$\mathcal{F}$ 如上。存在 $D(X_A)$ 中的典范对象
$L = L(\Lambda, X, A, \mathcal{F})$，使得对任意核为平方零理想 $I$ 的
$\Lambda$-代数满射 $A' \to A$，均有：
\begin{enumerate}
\item 范畴 $\textit{Lift}(\mathcal{F}, A')$ 非空，当且仅当某个类
$\xi \in \Ext^2_{X_A}(L, \mathcal{F} \otimes_A I)$ 为零。
\item 若 $\textit{Lift}(\mathcal{F}, A')$ 非空，则
$\text{Lift}(\mathcal{F}, A')$ 是
$\Ext^1_{X_A}(L, \mathcal{F} \otimes_A I)$ 作用下的主齐性集。
\item 给定一个提升 $\mathcal{F}'$，其拉回为
$\text{id}_\mathcal{F}$ 的 $\mathcal{F}'$ 自同构所成集合典范同构于
$\Ext^0_{X_A}(L, \mathcal{F} \otimes_A I)$.
\end{enumerate}
\end{proposition}

\begin{proof}
FIXME.
\end{proof}

\begin{lemma}
\label{obsolete-lemma-pseudo-coherent}
在命题 \ref{obsolete-proposition-conclusion} 的情形中，若
$X \to \Spec(\Lambda)$ 局部有限型且 $\Lambda$ 为诺特环，则 $L$ 伪凝聚。
\end{lemma}

\begin{proof}
FIXME.
\end{proof}






\section{非平坦情形下的凝聚层叠}
\label{obsolete-section-not-flat}

\noindent
Quot 的定理 \ref{quot-theorem-coherent-algebraic} 实际上不需要
$f : X \to B$ 平坦这一假设。本节依据导出代数几何中的思想修改证明方法，
以绕开平坦性假设。Quot 的小节 \ref{quot-section-not-flat}
采用一种完全不同的方法得到完全相同的结果；因此本节的方法已经废弃。

\medskip\noindent
Quot 的定理 \ref{quot-theorem-coherent-algebraic} 的证明中，
唯一使用平坦性的步骤是应用 Quot 的引理
\ref{quot-lemma-coherent-defo-thy}。该引理用于构造 Artin's Axioms
小节 \ref{artin-section-dual} 中的那种障碍理论。其证明依赖
Deformation Theory 的引理 \ref{defos-lemma-flat-ringed-topoi} 与
\ref{defos-lemma-verify-iv-ringed-topoi}；这两个引理见 Deformation Theory
小节 \ref{defos-section-flat-ringed-topoi}。$f$ 的平坦性假设正由此产生。
继续之前请注意，即使不假设 $f$ 平坦，Deformation Theory 的引理
\ref{defos-lemma-flat-ringed-topoi} 中结论 (2) 与 (3) 仍然成立，
因为它们依赖的 Deformation Theory 引理
\ref{defos-lemma-inf-ext-rel-ringed-topoi} 与
\ref{defos-lemma-inf-map-rel-ringed-topoi} 均不含任何平坦性假设。

\medskip\noindent
在给出细节之前，我们先从导出代数几何的角度说明这一构造的动机，
因为这有助于澄清下文。设 $A$ 是局部诺特基底 $S$ 上的有限型代数。
以 $X \otimes^\mathbf{L} A$ 表示 $X$ 到 $A$ 的“导出基变换”，并以
$i : X_A \to X \otimes^\mathbf{L} A$ 表示典范嵌入态射。
对象 $X \otimes^\mathbf{L} A$ 在 Stacks 项目中（尚）无定义；
可以把它想成配备单纯环层
$\mathcal{O}_{X \otimes^\mathbf{L} A}$ 的代数空间 $X_A$，其同调层为
$$
H_i(\mathcal{O}_{X \otimes^\mathbf{L} A}) =
\text{Tor}^{\mathcal{O}_S}_i(\mathcal{O}_X, A).
$$
态射 $X \otimes^\mathbf{L} A \to \Spec(A)$ 是平坦的
（该单纯环层的各项均在 $A$ 上平坦），故平坦模形变的通常理论适用于它。
由此可见，我们用下列群得到一个障碍理论：
$$
\Ext^i_{X \otimes^\mathbf{L} A}(i_*\mathcal{F},
i_*\mathcal{F} \otimes_A M)
$$
其中 $i = 0, 1, 2$ 分别对应无穷小自同构、无穷小形变和障碍。
注意，$i_*\mathcal{F}$ 到 $X \otimes^\mathbf{L} A'$ 的平坦形变自动具有
$i'_*\mathcal{F}'$ 的形式，其中 $\mathcal{F}'$ 是 $\mathcal{F}$ 的平坦形变。
由函子 $Li^*$ 与 $i_* = Ri_*$ 的伴随性，这些 Ext 群等于
$$
\Ext^i_{X_A}(Li^*(i_*\mathcal{F}), \mathcal{F} \otimes_A M)
$$
因此所得障碍群与 Quot 的引理 \ref{quot-lemma-coherent-defo-thy}
证明中的形式完全相同；唯一的变化，是将第一次出现的 $\mathcal{F}$
替换为复形 $Li^*(i_*\mathcal{F})$。

\medskip\noindent
下文将“直接”构造 $E(\mathcal{F}) = Li^*(i_*\mathcal{F})$，并直接证明它与
$\mathcal{F}$ 的形变理论之间的关系，由此证明该引理的非平坦版本。
事实上，构造 $\tau_{\geq -2}E(\mathcal{F})$ 已经足够，因为我们只关心
$\Ext^i_{X_A}(Li^*(i_*\mathcal{F}), \mathcal{F} \otimes_A M)$
在 $i = 0, 1, 2$ 时的 Ext 群。我们甚至可以确定其上同调层：
$$
H^i(E(\mathcal{F})) =
\left\{
\begin{matrix}
0 & \text{若 }i > 0 \\
\mathcal{F} & \text{若 } i = 0 \\
0 & \text{若 } i = -1 \\
\text{Tor}_1^{\mathcal{O}_S}(\mathcal{O}_X, A)
\otimes_{\mathcal{O}_X} \mathcal{F} &
\text{若 } i = -2
\end{matrix}
\right.
$$
这一观察将指导下面各条注记中对 $E(\mathcal{F})$ 的构造。

\begin{remark}[直接构造]
\label{obsolete-remark-construction-E}
设 $S$ 为概形，$f : X \to B$ 为 $S$ 上代数空间的态射，
$U$ 为另一个 $B$ 上代数空间。以 $q : X \times_B U \to U$
表示第二投影。考虑优三角
$$
Lq^*L_{U/B} \to L_{X \times_B U/B} \to E \to Lq^*L_{U/B}[1]
$$
，见 Cotangent 小节 \ref{cotangent-section-fibre-product}。
对任意 $\mathcal{O}_{X \times_B U}$-模层 $\mathcal{F}$，都有阿蒂亚类
$$
\mathcal{F} \to
L_{X \times_B U/B}
\otimes_{\mathcal{O}_{X \times_B U}}^\mathbf{L} \mathcal{F}[1]
$$
（见 Cotangent 小节 \ref{cotangent-section-atiyah-general}）。
将其与到 $E$ 的态射复合，并在 $D(\mathcal{O}_{X \times_B U})$ 中选取优三角
$$
E(\mathcal{F}) \to \mathcal{F} \to
\mathcal{F} \otimes_{\mathcal{O}_{X \times_B U}}^\mathbf{L} E[1] \to
E(\mathcal{F})[1]
$$
。由构造，阿蒂亚类可提升为态射
$$
e_\mathcal{F} :
E(\mathcal{F})
\longrightarrow
Lq^*L_{U/B} \otimes_{\mathcal{O}_{X \times_B U}}^\mathbf{L} \mathcal{F}[1]
$$
，并嵌入下列优三角态射：
$$
\xymatrix{
\mathcal{F} \otimes^\mathbf{L} Lq^*L_{U/B}[1] \ar[r] &
\mathcal{F} \otimes^\mathbf{L} L_{X \times_B U/B}[1] \ar[r] &
\mathcal{F} \otimes^\mathbf{L} E[1] \\
E(\mathcal{F}) \ar[r] \ar[u]^{e_\mathcal{F}} &
\mathcal{F} \ar[r] \ar[u]^{Atiyah} &
\mathcal{F} \otimes^\mathbf{L} E[1] \ar[u]^{=}
}
$$
给定 $S, B, X, f, U, \mathcal{F}$ 后，我们固定 $E(\mathcal{F})$
与 $e_\mathcal{F}$ 的一种选取。
\end{remark}

\begin{remark}[障碍类的构造]
\label{obsolete-remark-construction-ob}
采用注记 \ref{obsolete-remark-construction-E} 的记号，设 $i : U \to U'$ 是
$U$ 在 $B$ 上的一阶增厚。设
% P12-ERR-0508
$\mathcal{I} \subset \mathcal{O}_{U'}$ 是定义 $B$ 在 $B'$ 中嵌入的
拟凝聚理想层。基本三角
$$
Li^*L_{U'/B} \to L_{U/B} \to L_{U/U'} \to Li^*L_{U'/B}[1]
$$
与态射 $L_{U/U'} \to \mathcal{I}[1]$ 共同确定态射
$e_{U'} : L_{U/B} \to \mathcal{I}[1]$。将它同上一注记中的态射
$e_\mathcal{F}$ 结合，得到
$$
(\text{id}_\mathcal{F} \otimes Lq^*e_{U'}) \cup e_\mathcal{F} :
E(\mathcal{F})
\longrightarrow
\mathcal{F} \otimes_{\mathcal{O}_{X \times_B U}} q^*\mathcal{I}[2]
$$
（这里还同从导出张量积到通常张量积的态射作了复合）。换言之，我们得到元素
$$
\xi_{U'} \in
\Ext^2_{\mathcal{O}_{X \times_B U}}(
E(\mathcal{F}),
\mathcal{F} \otimes_{\mathcal{O}_{X \times_B U}} q^*\mathcal{I})
$$
\end{remark}

\begin{lemma}
\label{obsolete-lemma-ob-is-obstruction}
在注记 \ref{obsolete-remark-construction-ob} 的情形中，假设 $\mathcal{F}$ 在 $U$ 上平坦。
则类 $\xi_{U'}$ 为零，是存在满足 $i^*\mathcal{F}' \cong \mathcal{F}$
且在 $U'$ 上平坦的 $\mathcal{O}_{X \times_B U'}$-模 $\mathcal{F}'$
的充分必要条件。
\end{lemma}

\begin{proof}[证明（概要）]
我们将使用 Deformation Theory 的引理
\ref{defos-lemma-inf-obs-ext-rel-ringed-topoi} 中的判据。
简记 $\mathcal{O} = \mathcal{O}_{X \times_B U}$ 以及
$\mathcal{O}' = \mathcal{O}_{X \times_B U'}$。考虑短正合列
$$
0 \to \mathcal{I} \to \mathcal{O}_{U'} \to \mathcal{O}_U \to 0.
$$
设 $\mathcal{J} \subset \mathcal{O}'$ 是定义 $X \times_B U$ 的拟凝聚理想层。
由上可得正合列
$$
\text{Tor}_1^{\mathcal{O}_B}(\mathcal{O}_X, \mathcal{O}_U) \to
q^*\mathcal{I} \to \mathcal{J} \to 0
$$
其中 $\text{Tor}_1^{\mathcal{O}_B}(\mathcal{O}_X, \mathcal{O}_U)$ 是下式的简写：
$$
\text{Tor}_1^{h^{-1}\mathcal{O}_B}(p^{-1}\mathcal{O}_X, q^{-1}\mathcal{O}_U)
\otimes_{(p^{-1}\mathcal{O}_X\otimes_{h^{-1}\mathcal{O}_B}q^{-1}\mathcal{O}_U)}
\mathcal{O}.
$$
与 $\mathcal{F}$ 作张量积，得到正合列
$$
\mathcal{F} \otimes_\mathcal{O}
\text{Tor}_1^{\mathcal{O}_B}(\mathcal{O}_X, \mathcal{O}_U) \to
\mathcal{F} \otimes_\mathcal{O}
q^*\mathcal{I} \to
\mathcal{F} \otimes_\mathcal{O} \mathcal{J} \to 0
$$
（注意，相对于 Deformation Theory 的引理
\ref{defos-lemma-inf-obs-ext-rel-ringed-topoi} 中的记号，字母
$\mathcal{I}$ 与 $\mathcal{J}$ 的作用恰好对调。）该引理的条件 (1)
要求上面的最后一个态射为同构，亦即第一个态射为零。
这一态射是否为零，可在几何点
$\overline{z} = (\overline{x}, \overline{u}) : \Spec(k) \to X \times_B U$
处的茎上检验。置 $R = \mathcal{O}_{B, \overline{b}}$、
$A = \mathcal{O}_{X, \overline{x}}$、$B = \mathcal{O}_{U, \overline{u}}$，以及
$C = \mathcal{O}_{\overline{z}}$。由 Cotangent 的引理
\ref{cotangent-lemma-fibre-product} 和 $E(\mathcal{F})$ 的定义三角可见
$$
H^{-2}(E(\mathcal{F}))_{\overline{z}} =
\mathcal{F}_{\overline{z}} \otimes \text{Tor}_1^R(A, B)
$$
因此态射 $\xi_{U'}$ 诱导态射
$$
\mathcal{F}_{\overline{z}} \otimes \text{Tor}_1^R(A, B)
\longrightarrow
\mathcal{F}_{\overline{z}} \otimes_B \mathcal{I}_{\overline{u}}
$$
我们断言，它与上述态射的茎相同（证明从略；这是一个纯粹的环论陈述）。
由此可见，Deformation Theory 的引理
\ref{defos-lemma-inf-obs-ext-rel-ringed-topoi} 的条件 (1) 等价于下列态射为零：
$H^{-2}(\xi_{U'}) :
H^{-2}(E(\mathcal{F})) \to \mathcal{F} \otimes \mathcal{I}$.

\medskip\noindent
为完成证明，我们说明：在条件 (1) 成立的假设下，条件 (2) 等价于
$\xi_{U'}$ 为零。在证明的余下部分，以 $\mathcal{F} \otimes \mathcal{I}$ 表示
$\mathcal{F} \otimes_\mathcal{O} q^*\mathcal{I} =
\mathcal{F} \otimes_\mathcal{O} \mathcal{J}$。考察谱序列
$$
\Ext^i(H^{-j}(E(\mathcal{F})), \mathcal{F} \otimes \mathcal{I})
\Rightarrow
\Ext^{i + j}(E(\mathcal{F}), \mathcal{F} \otimes \mathcal{I})
$$
并利用 $H^0(E(\mathcal{F})) = \mathcal{F}$ 与
$H^{-1}(E(\mathcal{F})) = 0$，可知存在正合列
$$
0 \to
\Ext^2(\mathcal{F}, \mathcal{F} \otimes \mathcal{I}) \to
\Ext^2(E(\mathcal{F}), \mathcal{F} \otimes \mathcal{I}) \to
\Hom(H^{-2}(E(\mathcal{F})), \mathcal{F} \otimes \mathcal{I})
$$
因此元素 $\xi_{U'}$ 属于
$\Ext^2(\mathcal{F}, \mathcal{F} \otimes \mathcal{I})$。
只需说明它与 Deformation Theory 的引理
\ref{defos-lemma-inf-obs-ext-rel-ringed-topoi} 中的元素相同即可完成证明；
这一验证从略。
\end{proof}

\begin{lemma}
\label{obsolete-lemma-coherent-defo-thy-general}
在 Quot 的情形 \ref{quot-situation-coherent} 中，假设 $S$ 是局部诺特概形且
$S = B$。令 $\mathcal{X} = \textit{Coh}_{X/B}$。则 $\mathcal{X}$
具有 versal 性的开性（见 Artin's Axioms 的定义
\ref{artin-definition-openness-versality}）。
\end{lemma}

\begin{proof}[证明（概要）]
设 $U \to S$ 为有限型概形态射，$x$ 为 $U$ 上的 $\mathcal{X}$-对象，
$u_0 \in U$ 为有限型点，且 $x$ 在 $u_0$ 处为 versal。缩小 $U$ 后，
可假设 $u_0$ 是闭点（Morphisms 的引理
\ref{morphisms-lemma-point-finite-type}），并且 $U = \Spec(A)$，
而 $U \to S$ 的像落在 $S$ 的仿射开集 $\Spec(\Lambda)$ 中。
我们将使用 Artin's Axioms 的引理 \ref{artin-lemma-dual-openness}
证明本引理。设 $\mathcal{F}$ 是 $X_A = \Spec(A) \times_S X$ 上
对应于给定对象 $x$、且在 $A$ 上平坦的凝聚模。

\medskip\noindent
按注记 \ref{obsolete-remark-construction-E} 选取 $E(\mathcal{F})$ 与
$e_\mathcal{F}$。由 $E(\mathcal{F})$ 的上同调层描述可知，对任意
$A$-模 $M$ 都有
$$
\Ext^1(E(\mathcal{F}), \mathcal{F} \otimes_A M) =
\Ext^1(\mathcal{F}, \mathcal{F} \otimes_A M)
$$
。结合这一点与 Deformation Theory 的引理
\ref{defos-lemma-inf-ext-rel-ringed-topoi}，得到函子同构
$$
T_x(M) = \Ext^1_{X_A}(E(\mathcal{F}), \mathcal{F} \otimes_A M)
$$
。由引理 \ref{obsolete-lemma-ob-is-obstruction}，对任意核为平方零理想 $I$ 的
$\Lambda$-代数满射 $A' \to A$，都有障碍类
$$
\xi_{A'} \in \Ext^2_{X_A}(E(\mathcal{F}), \mathcal{F} \otimes_A I)
$$
。将 Derived Categories of Spaces 的引理
\ref{spaces-perfect-lemma-compute-ext} 应用于 Ext 群
$\Ext^i_{X_A}(E(\mathcal{F}), \mathcal{F} \otimes_A M)$
在 $i \leq m$、$m = 2$ 时的计算。我们略去
$E(\mathcal{F})$ 属于 $D^-_{\textit{Coh}}$ 的验证；提示：使用
Cotangent 的引理 \ref{cotangent-lemma-cotangent-finite}。
于是得到完美对象 $K \in D(A)$ 以及函子性同构
$$
H^i(K \otimes_A^\mathbf{L} M)
\longrightarrow
\Ext^i_{X_A}(E(\mathcal{F}), \mathcal{F} \otimes_A M)
$$
，其中 $i \leq m$，且这些同构与边界态射相容。对象 $K$ 连同上述同一化，
给出了 Artin's Axioms 的情形 \ref{artin-situation-dual} 中的一组数据。
最后，Artin's Axioms 的引理 \ref{artin-lemma-dual-obstruction}
之条件 (iv)，可由 Deformation Theory 的引理
\ref{defos-lemma-verify-iv-ringed-topoi} 的一个变体推出；
我们略去该变体的陈述与证明。因此 Artin's Axioms 的引理
\ref{artin-lemma-dual-openness} 适用，引理得证。
\end{proof}

\begin{theorem}
\label{obsolete-theorem-coherent-algebraic-general}
设 $S$ 为概形，$f : X \to B$ 为 $S$ 上代数空间的态射。
假设 $f$ 有限表示且分离。则 $\textit{Coh}_{X/B}$ 是 $S$ 上的代数叠。
\end{theorem}

\begin{proof}
本定理是 Quot 的定理 \ref{quot-theorem-coherent-algebraic-general}
的一份副本。之所以在此保留它，是因为本节材料给出了第二种证明
（如本节开头所述）。具体而言，论证与 Quot 的定理
\ref{quot-theorem-coherent-algebraic} 的证明完全相同，
但以引理 \ref{obsolete-lemma-coherent-defo-thy-general} 取代 Quot 的引理
\ref{quot-lemma-coherent-defo-thy}。
\end{proof}







\section{修改}
\label{obsolete-section-modifications}

\noindent
下面是关于 Algebraization of Formal Spaces 中式
(\ref{restricted-equation-modification}) 所给范畴的一项已废弃结果。
现行材料请参见 Algebraization of Formal Spaces 的小节
\ref{restricted-section-modifications}。

\begin{lemma}
\label{obsolete-lemma-henselian}
设 $(A, \mathfrak m, \kappa)$ 为诺特局部环。Algebraization of Formal Spaces
的式 (\ref{restricted-equation-modification}) 给出的 $A$ 的范畴，
等价于同一式 (\ref{restricted-equation-modification}) 给出的
$A$ 的亨泽尔化 $A^h$ 的范畴。
\end{lemma}

\begin{proof}
这是 Algebraization of Formal Spaces 的引理
\ref{restricted-lemma-equivalence-to-completion} 的特殊情形。
\end{proof}

\noindent
下面关于有理奇点的引理，在曲面奇点消解一章中已不再需要。

\begin{lemma}
\label{obsolete-lemma-double-dual-rational}
在 Resolution of Surfaces 的情形 \ref{resolve-situation-rational} 中，
设 $M$ 为有限反身 $A$-模。以 $M \otimes_A \mathcal{O}_X$
表示相应 $\mathcal{O}_S$-模的拉回。则
$M \otimes_A \mathcal{O}_X$ 满射到其双对偶。
\end{lemma}

\begin{proof}
令 $\mathcal{F} = (M \otimes_A \mathcal{O}_X)^{**}$ 为双对偶，
并令 $\mathcal{F}' \subset \mathcal{F}$ 为求值态射
$M \otimes_A \mathcal{O}_X \to \mathcal{F}$ 的像。于是有短正合列
$$
0 \to \mathcal{F}' \to \mathcal{F} \to \mathcal{Q} \to 0
$$
。由于 $X$ 正规，对于余维 $1$ 的点，局部环 $\mathcal{O}_{X, x}$
均为离散赋值环（见 Properties 的引理
\ref{properties-lemma-criterion-normal}）。因此由 More on Algebra 的引理
\ref{more-algebra-lemma-cokernel-map-double-dual-dvr}，
在这类点上有 $\mathcal{Q}_x = 0$。故 $\mathcal{Q}$ 支撑在有限多个闭点上，
并由 Cohomology of Schemes 的引理
\ref{coherent-lemma-coherent-support-dimension-0} 得到它由整体截面生成。
我们得到正合列
$$
0 \to H^0(X, \mathcal{F}') \to H^0(X, \mathcal{F}) \to H^0(X, \mathcal{Q}) \to 0
$$
，因为 $\mathcal{F}'$ 由整体截面生成（Resolution of Surfaces 的引理
\ref{resolve-lemma-globally-generated}）。由于 $X \to \Spec(A)$
在闭点的补集上为同构，且 $M$ 反身，可见态射
$$
M \to H^0(X, \mathcal{F}') \to H^0(X, \mathcal{F})
$$
在 $A$ 的任意非极大素理想处局部化后诱导同构。由 More on Algebra 的引理
\ref{more-algebra-lemma-check-isomorphism-via-depth-and-ass}，以及正规环上的
反身模具有性质 $(S_2)$ 这一事实（More on Algebra 的引理
\ref{more-algebra-lemma-reflexive-over-normal}），这些态射因而都是同构。
故 $\mathcal{Q} = 0$，正如所需。
\end{proof}



\section{相交理论}
\label{obsolete-section-intersection-theory}

\begin{lemma}
\label{obsolete-lemma-good-blowing-up}
设 $b : X' \to X$ 是域 $k$ 上光滑射影概形沿其光滑闭子概形
$Z \subset X$ 的吹起。图示为
$$
\xymatrix{
E \ar[r]_j \ar[d]_\pi & X' \ar[d]^b \\
Z \ar[r]^i & X
}
$$
假设存在 $K_0(X)$ 中的元素，其在 $Z$ 上的限制等于
$K_0(Z)$ 中 $\mathcal{C}_{Z/X}$ 的类。则对某个
$\alpha'' \in K_0(X')$，在 $K_0(X')$ 中有
$[Lb^*\mathcal{O}_Z] = [\mathcal{O}_E] \cdot \alpha''$。
\end{lemma}

\begin{proof}
概形 $X$、$X'$、$E$、$Z$ 均在 $k$ 上光滑且射影，因而
$K'_0(X) = K_0(X) = K_0(\textit{Vect}(X)) =
K_0(D^b_{\textit{Coh}}(X)))$，其余 $3$ 个对象亦然。参见
Derived Categories of Schemes 的引理
\ref{perfect-lemma-Noetherian-Kprime}、\ref{perfect-lemma-Kprime-K}
与 \ref{perfect-lemma-K-is-old-K}。本证明中将任意换用这些版本。
考虑定义 $\mathcal{F}$ 的有限局部自由 $\mathcal{O}_E$-模短正合列
$$
0 \to \mathcal{F} \to \pi^*\mathcal{C}_{Z/X} \to \mathcal{C}_{E/X'} \to 0
$$
% P12-ERR-0509
。注意，$\mathcal{C}_{E/X'} = \mathcal{O}_{X'}(-E)|_E$
是可逆 $\mathcal{O}_X$-模 $\mathcal{O}_{X'}(-E)$ 的限制。
取 $\alpha \in K_0(X)$，使得在 $K_0(Z)$ 中
$i^*\alpha = [\mathcal{C}_{Z/X}]$。令
$\alpha' = b^*\alpha - [\mathcal{O}_{X'}(-E)]$。
于是 $j^*\alpha' = [\mathcal{F}]$。由 Weil Cohomology Theories 的引理
\ref{weil-lemma-lambda-operations}，推出
$j^*\lambda^i(\alpha') = [\wedge^i(\mathcal{F})]$。
% P12-ERR-0510 P12-ERR-0511
这意味着在 $K_0(X)$ 中
$[\mathcal{O}_E] \cdot \alpha' = [\wedge^i\mathcal{F}]$；参见
Derived Categories of Schemes 的引理
\ref{perfect-lemma-projection-formula}。令 $r$ 为 $Z$ 的不可约分支在 $X$
中的最大余维。一个从略的计算表明
% P12-ERR-0512
$H^{-i}(Lb^*\mathcal{O}_Z) = \wedge^i\mathcal{F}$
当 $i \geq 0, 1, \ldots, r - 1$ 时成立，而在其余次数为零。
因此在 $K_0(X)$ 中有
\begin{align*}
[Lb^*\mathcal{O}_Z] & =
\sum\nolimits_{i = 0, \ldots, r - 1} (-1)^i[\wedge^i\mathcal{F}] \\
& =
\sum\nolimits_{i = 0, \ldots, r - 1} (-1)^i[\mathcal{O}_E] \lambda^i(\alpha') \\
& =
[\mathcal{O}_E] \left(\sum\nolimits_{i = 0, \ldots, r - 1}
(-1)^i \lambda^i(\alpha')\right)
\end{align*}
取 $\alpha'' = \sum_{i = 0, \ldots, r - 1} (-1)^i \lambda^i(\alpha')$
即得本引理。
\end{proof}



\begin{lemma}
\label{obsolete-lemma-gysin-factors-principal}
设 $(S, \delta)$ 如 Chow Homology 的情形 \ref{chow-situation-setup}。
设 $X$ 在 $S$ 上局部有限型。假设 $X$ 整，且 $n = \dim_\delta(X)$。
设 $a \in \Gamma(X, \mathcal{O}_X)$ 为非零函数，并设
$i : D = Z(a) \to X$ 为 $a$ 的零概形之闭浸入。令 $f \in R(X)^*$。
此时在 $A_{n - 2}(D)$ 中有 $i^*\text{div}_X(f) = 0$。
\end{lemma}

\begin{proof}
这是 Chow Homology 的引理 \ref{chow-lemma-gysin-factors-general} 的特殊情形。
\end{proof}

\begin{remark}
\label{obsolete-remark-not-true-not-quasi-compact}
本注记过去曾说，尚不清楚 Chow Homology 的引理
\ref{chow-lemma-cycles-k-group} 中的箭头在一般情况下是否为同构。
不过，现在我们已经找到了这一事实的证明。
\end{remark}







\subsection{吹起引理}
\label{obsolete-section-blowing-up-lemmas}

\noindent
本节证明若干引理：在吹起上用适当的有效 Cartier 因子表示 Cartier 因子。
这些引理见 \cite[Section 2.4]{F}。我们调整了表述，使其在非有限型情形下
同样适用。引理 \ref{obsolete-lemma-blowing-up-intersections} 中的态射 $b$
可能是无穷多个吹起的复合，但在任意给定的拟紧开集 $W \subset X$ 上，
只需有限多个吹起（这正是同上引文中的结果；loc.\ cit.）。

\begin{lemma}
\label{obsolete-lemma-push-pull-effective-Cartier}
设 $(S, \delta)$ 如 Chow Homology 的情形 \ref{chow-situation-setup}。
设 $X$、$Y$ 在 $S$ 上局部有限型，$f : X \to Y$ 为固有态射，
$D \subset Y$ 为有效 Cartier 因子。假设 $X$、$Y$ 整，
$n = \dim_\delta(X) = \dim_\delta(Y)$，且 $f$ 支配。则
$$
f_*[f^{-1}(D)]_{n - 1} = [R(X) : R(Y)] [D]_{n - 1}.
$$
特别地，若 $f$ 双有理，则 $f_*[f^{-1}(D)]_{n - 1} = [D]_{n - 1}$。
\end{lemma}

\begin{proof}
这立即来自 Chow Homology 的引理 \ref{chow-lemma-equal-c1-as-cycles}，
以及 $D$ 是 $\mathcal{O}_X(D)$ 的典范截面 $1_D$ 的零概形这一事实。
\end{proof}

\begin{lemma}
\label{obsolete-lemma-blowing-up-denominators}
设 $(S, \delta)$ 如 Chow Homology 的情形 \ref{chow-situation-setup}。
设 $X$ 在 $S$ 上局部有限型，并假设 $X$ 整且 $\dim_\delta(X) = n$。
设 $\mathcal{L}$ 为可逆 $\mathcal{O}_X$-模，$s$ 为 $\mathcal{L}$ 的非零亚纯截面。
令 $U \subset X$ 为使 $s$ 对应于 $\mathcal{L}$ 在 $U$ 上一个截面的最大开子概形。
则存在射影态射
$$
\pi : X' \longrightarrow X
$$
，满足
\begin{enumerate}
\item $X'$ 整；
\item $\pi|_{\pi^{-1}(U)} : \pi^{-1}(U) \to U$ 是同构；
\item 存在有效 Cartier 因子 $D, E \subset X'$，使得
$$
\pi^*\mathcal{L} = \mathcal{O}_{X'}(D - E),
$$
\item 通过上述同构，亚纯截面 $s$ 对应于亚纯截面
$1_D \otimes (1_E)^{-1}$（见 Divisors 的定义
\ref{divisors-definition-invertible-sheaf-effective-Cartier-divisor}）；
\item 有
$$
\pi_*([D]_{n - 1} - [E]_{n - 1}) = \text{div}_\mathcal{L}(s)
$$
在 $Z_{n - 1}(X)$ 中成立。
\end{enumerate}
\end{lemma}

\begin{proof}
设 $\mathcal{I} \subset \mathcal{O}_X$ 为 $s$ 的拟凝聚分母理想层，
见 Divisors 的定义
\ref{divisors-definition-regular-meromorphic-ideal-denominators}。
由 Divisors 的引理 \ref{divisors-lemma-blowing-up-denominators}
得到 (2)、(3)、(4)。由 Divisors 的引理
\ref{divisors-lemma-blow-up-integral-scheme} 得到 (1)。
由 Divisors 的引理 \ref{divisors-lemma-blowing-up-projective}，
态射 $\pi$ 为射影态射。尚需证明 (5)。由 Chow Homology 的引理
\ref{chow-lemma-equal-c1-as-cycles}，有
$$
\pi_*(\text{div}_{\mathcal{L}'}(s')) = \text{div}_\mathcal{L}(s).
$$
因此只需证明
$\text{div}_{\mathcal{L}'}(s') = [D]_{n - 1} - [E]_{n - 1}$。
这由等式 $s' = 1_D \otimes 1_E^{-1}$ 和可加性推出；
见 Divisors 的引理 \ref{divisors-lemma-c1-additive}。
\end{proof}

\begin{definition}
\label{obsolete-definition-epsilon}
设 $(S, \delta)$ 如 Chow Homology 的情形 \ref{chow-situation-setup}。
设 $X$ 在 $S$ 上局部有限型，并假设 $X$ 整且 $\dim_\delta(X) = n$。
设 $D_1, D_2$ 为 $X$ 中的两个有效 Cartier 因子，
$Z \subset X$ 为满足 $\dim_\delta(Z) = n - 1$ 的整闭子概形。
这一情形的{\it $\epsilon$-不变量}为
$$
\epsilon_Z(D_1, D_2) = n_Z \cdot m_Z
$$
，其中 $n_Z$（resp.\ $m_Z$）是 $Z$ 在 $(n - 1)$-循环
$[D_1]_{n - 1}$（resp.\ $[D_2]_{n - 1}$）中的系数。
\end{definition}

\begin{lemma}
\label{obsolete-lemma-two-divisors}
设 $(S, \delta)$ 如 Chow Homology 的情形 \ref{chow-situation-setup}。
设 $X$ 在 $S$ 上局部有限型，并假设 $X$ 整且 $\dim_\delta(X) = n$。
设 $D_1, D_2$ 为 $X$ 中的两个有效 Cartier 因子，
$Z$ 为概形 $D_1 \cap D_2$ 的开闭子概形。假设
$\dim_\delta(D_1 \cap D_2 \setminus Z) \leq n - 2$。
则存在态射 $b : X' \to X$ 以及 $X'$ 上的 Cartier 因子
$D_1', D_2', E$，满足下列性质：
\begin{enumerate}
\item $X'$ 整；
\item $b$ 为射影态射；
\item $b$ 是 $X$ 沿闭子概形 $Z$ 的吹起；
\item $E = b^{-1}(Z)$；
\item $b^{-1}(D_1) = D'_1 + E$，且 $b^{-1}D_2 = D_2' + E$；
\item $\dim_\delta(D'_1 \cap D'_2) \leq n - 2$；并且若
$Z = D_1 \cap D_2$，则 $D'_1 \cap D'_2 = \emptyset$；
\item 对每个满足 $\dim_\delta(W') = n - 1$ 的整闭子概形 $W'$，均有：
\begin{enumerate}
\item 若 $\epsilon_{W'}(D'_1, E) > 0$，则置 $W = b(W')$ 后有
$\dim_\delta(W) = n - 1$，且
$$
\epsilon_{W'}(D'_1, E) < \epsilon_W(D_1, D_2),
$$
\item 若 $\epsilon_{W'}(D'_2, E) > 0$，则置 $W = b(W')$ 后有
$\dim_\delta(W) = n - 1$，且
$$
\epsilon_{W'}(D'_2, E) < \epsilon_W(D_1, D_2),
$$
\end{enumerate}
\end{enumerate}
\end{lemma}

\begin{proof}
注意，拟凝聚理想层
$\mathcal{I} = \mathcal{I}_{D_1} + \mathcal{I}_{D_2}$
定义了概形论交 $D_1 \cap D_2 \subset X$。由于 $Z$ 是
$D_1 \cap D_2$ 的若干连通分支之并，对每个 $z \in Z$，态射
$\mathcal{O}_{X, z} \to \mathcal{O}_{Z, z}$ 的核等于 $\mathcal{I}_z$。
令 $b : X' \to X$ 为 $X$ 沿 $Z$ 的吹起（所以在 $Z$ 附近 Zariski 局部，
它是 $X$ 沿 $\mathcal{I}$ 的吹起）。以 $E = b^{-1}(Z)$ 表示相应的
有效 Cartier 因子；见 Divisors 的引理
\ref{divisors-lemma-blowing-up-gives-effective-Cartier-divisor}。
% P12-ERR-0513
由于 $Z \subset D_1$，有 $E \subset f^{-1}(D_1)$，故对某个有效
Cartier 因子 $D'_1 \subset X'$，有 $D_1 = D_1' + E$；见 Divisors 的引理
\ref{divisors-lemma-difference-effective-Cartier-divisors}。
同理，$D_2 = D_2' + E$。这就处理了断言 (1) -- (5)。

\medskip\noindent
注意，若 $W'$ 如 (7) (a) 或 (7) (b) 所述，则 $W'$ 的像 $W$
包含于 $D_1 \cap D_2$。若 $W$ 不包含于 $Z$，则 $b$ 在 $W$ 的一般点处
为同构，从而 $\dim_\delta(W) = \dim_\delta(W') = n - 1$，
这与假设 $\dim_\delta(D_1 \cap D_2 \setminus Z) \leq n - 2$ 矛盾。
故 $W \subset Z$。这意味着为证明 (6) 与 (7)，可以在 $X$ 上 $Z$ 的附近局部工作。

\medskip\noindent
因此可假设 $X = \Spec(A)$，其中 $A$ 为诺特整环，并且
$D_1 = \Spec(A/a)$、$D_2 = \Spec(A/b)$、$Z = D_1 \cap D_2$。
置 $I = (a, b)$。由于 $A$ 为整环且 $a, b \not = 0$，可用两个片覆盖吹起，即
$U = \Spec(A[s]/(as - b))$ 与 $V = \Spec(A[t]/(bt -a))$。
这两个片通过同构
$A[s, s^{-1}]/(as - b) \cong A[t, t^{-1}]/(bt - a)$
黏合，该同构将 $s$ 映到 $t^{-1}$。有效 Cartier 因子 $E$ 由
$\Spec(A[s]/(as - b, a)) \subset U$ 与
$\Spec(A[t]/(bt - a, b)) \subset V$ 描述。闭子概形 $D'_1$ 对应于
% P12-ERR-0514
$\Spec(A[t]/(bt - a, t)) \subset U$。闭子概形 $D'_2$ 对应于
$\Spec(A[s]/(as -b, s)) \subset V$。由于“$ts = 1$”，可见
$D'_1 \cap D'_2 = \emptyset$。

\medskip\noindent
假设有高度为一的素理想 $\mathfrak q \subset A[s]/(as - b)$，且
$s, a \in \mathfrak q$。令 $\mathfrak p \subset A$ 为 $A$ 中相应的素理想。
注意 $a, b \in \mathfrak p$。由维数公式亦有
$\dim(A_{\mathfrak p}) = 1$。最后需要证明的断言是
$$
\text{ord}_{A_{\mathfrak p}}(a)
\text{ord}_{A_{\mathfrak p}}(b)
>
\text{ord}_{B_{\mathfrak q}}(a)
\text{ord}_{B_{\mathfrak q}}(s)
$$
，其中 $B = A[s]/(as - b)$。由 Algebra 的引理
\ref{algebra-lemma-quasi-finite-extension-dim-1}，对 $x = a, b$ 有
$\text{ord}_{A_{\mathfrak p}}(x) \geq \text{ord}_{B_{\mathfrak q}}(x)$。
又因 $\text{ord}_{B_{\mathfrak q}}(s) > 0$，由 $\text{ord}$ 函数的可加性
以及 $as = b$，所需不等式成立。
\end{proof}

\begin{definition}
\label{obsolete-definition-locally-finite-sum-effective-Cartier-divisors}
设 $X$ 为概形，$\{D_i\}_{i \in I}$ 为 $X$ 上一族局部有限的
有效 Cartier 因子。给定函数
$I \to \mathbf{Z}_{\geq 0}$，$i \mapsto n_i$。这些有效 Cartier 因子的
{\it 和} $D = \sum n_i D_i$，是唯一满足下列性质的有效 Cartier 因子
$D \subset X$：对任意拟紧开集 $U \subset X$，
$D|_U = \sum_{D_i \cap U \not = \emptyset} n_iD_i|_U$ 是 Divisors 的定义
\ref{divisors-definition-sum-effective-Cartier-divisors} 中所定义的和。
\end{definition}

\begin{lemma}
\label{obsolete-lemma-sum-divisors-associated-Weil}
设 $(S, \delta)$ 如 Chow Homology 的情形 \ref{chow-situation-setup}。
设 $X$ 在 $S$ 上局部有限型，并假设 $X$ 整且 $\dim_\delta(X) = n$。
设 $\{D_i\}_{i \in I}$ 为 $X$ 上一族局部有限的有效 Cartier 因子，
并对 $i \in I$ 给定 $n_i \geq 0$。则
$$
[D]_{n - 1} = \sum\nolimits_i n_i[D_i]_{n - 1}
$$
在 $Z_{n - 1}(X)$ 中成立。
\end{lemma}

\begin{proof}
由于要证明的是循环的等式，可以在 $X$ 上局部工作。因此问题化为有限和，
再由归纳化为两个有效 Cartier 因子的和 $D = D_1 + D_2$。
由 Chow Homology 的引理 \ref{chow-lemma-compute-c1} 可见
$D_1 = \text{div}_{\mathcal{O}_X(D_1)}(1_{D_1})$，其中
$1_{D_1}$ 表示 $\mathcal{O}_X(D_1)$ 的典范截面。$D_2$ 与 $D$ 当然也有同样的陈述。
在同一化 $\mathcal{O}_X(D) = \mathcal{O}_X(D_1) \otimes \mathcal{O}_X(D_2)$
下，$1_D = 1_{D_1} \otimes 1_{D_2}$，故结论由 Divisors 的引理
\ref{divisors-lemma-c1-additive} 得出。
\end{proof}

\begin{lemma}
\label{obsolete-lemma-blowing-up-intersections}
设 $(S, \delta)$ 如 Chow Homology 的情形 \ref{chow-situation-setup}。
设 $X$ 在 $S$ 上局部有限型，并假设 $X$ 整且 $\dim_\delta(X) = d$。
设 $\{D_i\}_{i \in I}$ 为 $X$ 上一族局部有限的有效 Cartier 因子。
假设对所有满足 $\#\{i, j, k\} = 3$ 的 $\{i, j, k\} \subset I$，都有
$D_i \cap D_j \cap D_k = \emptyset$。则存在
\begin{enumerate}
\item 开子概形 $U \subset X$，满足
$\dim_\delta(X \setminus U) \leq d - 3$；
\item 态射 $b : U' \to U$；
\item $U'$ 上的有效 Cartier 因子 $\{D'_j\}_{j \in J}$；
\end{enumerate}
它们具有下列性质：
\begin{enumerate}
\item $b$ 是固有态射 $b : U' \to U$；
\item $U'$ 整；
\item 在各 $D_i|_U$ 两两交之并的补集上，$b$ 为同构；
\item $\{D'_j\}_{j \in J}$ 是 $U'$ 上一族局部有限的有效 Cartier 因子；
\item 若 $j \not = j'$，则 $\dim_\delta(D'_j \cap D'_{j'}) \leq d - 2$；
\item 对某些 $n_{ij} \geq 0$，有
$b^{-1}(D_i|_U) = \sum n_{ij} D'_j$。
\end{enumerate}
此外，若 $X$ 拟紧，则上面可以假设 $U = X$。
\end{lemma}

\begin{proof}
先证明拟紧情形，因为这也许是最有意思的情形。在此情形中，归纳构造一列吹起
$$
X = X_0 \xleftarrow{b_0} X_1 \xleftarrow{b_1} X_2 \leftarrow \ldots
$$
以及有效 Cartier 因子的有限集 $\{D_{n, i}\}_{i \in I_n}$。
在每一步，任意三重交 $D_{n, i} \cap D_{n, j} \cap D_{n, k}$ 都为空。
此外，对每个 $n \geq 0$，都有
$I_{n + 1} = I_n \amalg P(I_n)$，其中 $P(I_n)$ 表示 $I_n$ 中元素对的集合。
最后，还有
$$
b_n^{-1}(D_{n, i}) = D_{n + 1, i} +
\sum\nolimits_{i' \in I_n, i' \not = i} D_{n + 1, \{i, i'\}}
$$
从而对每个 $n \geq 0$，$(b_0 \circ \ldots \circ b_n)^{-1}(D_i)$
都是各因子 $D_{n + 1, j}$（$j \in I_{n + 1}$）的非负整系数组合。

\medskip\noindent
为开始归纳，置 $X_0 = X$、$I_0 = I$、$D_{0, i} = D_i$。

\medskip\noindent
给定 $(X_n, \{D_{n, i}\}_{i \in I_n})$，令 $X_{n + 1}$ 为 $X_n$
沿闭子概形
$Z_n = \bigcup_{\{i, i'\} \in P(I_n)} D_{n, i} \cap D_{n, i'}$
的吹起。由关于三重交的假设，各闭子概形
$D_{n, i} \cap D_{n, i'}$ 两两不交。换言之，可以写成
$Z_n = \coprod_{\{i, i'\} \in P(I_n)} D_{n, i} \cap D_{n, i'}$.
此外，在 $D_{n, i} \cap D_{n, i'}$ 的 Zariski 邻域中，态射 $b_n$
等于概形 $X_n$ 沿闭子概形 $D_{n, i} \cap D_{n, i'}$ 的吹起，
故引理 \ref{obsolete-lemma-two-divisors} 的结论适用。因此置
% P12-ERR-0515
$D_{n + 1, \{i, i'\}} = b_n^{-1}(D_i \cap D_{i'})$，
便得到有效 Cartier 因子。各 Cartier 因子
$D_{n + 1, \{i, i'\}}$ 两两不交。显然，对每个
$i' \in I_n$、$i' \not = i$，都有
$b_n^{-1}(D_{n, i}) \supset D_{n + 1, \{i, i'\}}$。
因此应用 Divisors 的引理
\ref{divisors-lemma-difference-effective-Cartier-divisors}，可见确有
% P12-ERR-0516
$b^{-1}(D_{n, i}) = D_{n + 1, i} +
\sum\nolimits_{i' \in I_n, i' \not = i} D_{n + 1, \{i, i'\}}$
，其中 $D_{n + 1, i}$ 是 $X_{n + 1}$ 上的某个有效 Cartier 因子。
在 $D_{n + 1, \{i, i'\}}$ 的邻域中，这些因子 $D_{n + 1, i}$
扮演引理 \ref{obsolete-lemma-two-divisors} 中带撇因子的角色。特别地，
由引理 \ref{obsolete-lemma-two-divisors} 第 (6) 部分可知，若
$i \not = i'$ 且 $i, i' \in I_n$，则
$D_{n + 1, i} \cap D_{n + 1, i'} = \emptyset$。
这已经蕴含各因子 $D_{n + 1, i}$ 的三重交为零。

\medskip\noindent
至此，可以利用 $X$ 的拟紧性断定不变量
\begin{equation}
\label{obsolete-equation-invariant}
\epsilon(X, \{D_i\}_{i \in I}) =
\max\{\epsilon_Z(D_i, D_{i'}) \mid
Z \subset X,
\dim_\delta(Z) = d - 1,
\{i, i'\} \in P(I)\}
\end{equation}
是有限的；毕竟每个 $D_i$ 至多只有有限多个不可约分支。我们断言，对某个
$n$，不变量 $\epsilon(X_n, \{D_{n, i}\}_{i \in I_n})$ 为零。
否则，由引理 \ref{obsolete-lemma-two-divisors} 将得到严格递减序列
$$
\epsilon(X, \{D_i\}_{i \in I})
=
\epsilon(X_0, \{D_{0, i}\}_{i \in I_0})
>
\epsilon(X_1, \{D_{1, i}\}_{i \in I_1})
>
\ldots
$$
，其各项均为正整数，这就产生矛盾。取 $n$ 使得不变量
$\epsilon(X_n, \{D_{n, i}\}_{i \in I_n})$ 为零。这意味着不存在整闭子概形
$Z \subset X_n$ 和指标对 $i, i' \in I_n$，使得
$\epsilon_Z(D_{n, i}, D_{n, i'}) > 0$。换言之，对所有
$\{i, i'\} \in P(I_n)$，都有
% P12-ERR-0517
$\dim_\delta(D_{n, i}, D_{n, i'}) \leq d - 2$，正如所需。

\medskip\noindent
下面转向不再假设概形 $X$ 拟紧的一般情形。证明第一部分的思路存在一个问题：
可能得到一列无穷吹起，其中心支配 $X$ 的某个固定点。为避免这一点，
我们在每一步删去适当的余维 $\geq 3$ 的闭子集。具体而言，
将归纳构造如下形状的一列态射：
$$
\xymatrix{
X = X_0 \\
U_0 \ar[u]^{j_0} & X_1 \ar[l]_{b_0} \\
 & U_1 \ar[u]^{j_1} & X_2 \ar[l]_{b_1} \\
 & & U_2 \ar[u]^{j_2} & X_3 \ar[l]_{b_2}
}
$$
每个态射 $j_n : U_n \to X_n$ 都是开浸入。每个态射
$b_n : X_{n + 1} \to U_n$ 都是整概形之间的固有双有理态射。
同拟紧情形一样，在 $X_n$ 上有有效 Cartier 因子
$\{D_{n, i}\}_{i \in I_n}$。在每一步，任意三重交
$D_{n, i} \cap D_{n, j} \cap D_{n, k}$ 都为空。此外，对每个
$n \geq 0$，有 $I_{n + 1} = I_n \amalg P(I_n)$，其中 $P(I_n)$
表示 $I_n$ 中元素对的集合。最后，将使得
$$
b_n^{-1}(D_{n, i}|_{U_n}) = D_{n + 1, i} +
\sum\nolimits_{i' \in I_n, i' \not = i} D_{n + 1, \{i, i'\}}
$$

\medskip\noindent
归纳从置 $X_0 = X$、$I_0 = I$、$D_{0, i} = D_i$ 开始。

\medskip\noindent
给定 $(X_n, \{D_{n, i}\})$ 后，按如下方式构造开子概形 $U_n$。
对每个指标对 $\{i, i'\} \in P(I_n)$，考虑闭子概形
$D_{n, i} \cap D_{n, i'}$。它有 $\delta$-维数为 $d - 2$ 的“良”不可约分支，
以及 $\delta$-维数为 $d - 1$ 的“坏”不可约分支。置
$$
\text{Bad}(i, i')
=
\bigcup\nolimits_{W \subset D_{n, i} \cap D_{n, i'}
\text{ irred.\ comp. with }\dim_\delta(W) = d - 1} W
$$
，类似地置
$$
\text{Good}(i, i')
=
\bigcup\nolimits_{W \subset D_{n, i} \cap D_{n, i'}
\text{ irred.\ comp. with }\dim_\delta(W) = d - 2} W.
$$
于是 $D_{n, i} \cap D_{n, i'} = \text{Bad}(i, i') \cup \text{Good}(i, i')$，
并且
$\dim_\delta(\text{Bad}(i, i') \cap \text{Good}(i, i')) \leq d - 3$。
我们选取 $U_n$ 如下：
$$
U_n
=
X_n
\setminus
\bigcup\nolimits_{\{i, i'\} \in P(I_n)}
\text{Bad}(i, i') \cap \text{Good}(i, i').
$$
由各因子 $D_{n, i}$ 的三重交条件可见，该并实际上是不交并。
此外，作为概形有
$$
D_{n, i}|_{U_n} \cap D_{n, i'}|_{U_n}
=
Z_{n, i, i'} \amalg G_{n, i, i'}
$$
，其中 $Z_{n, i, i'}$ 为 $\delta$-等维的，维数为 $d - 1$；
$G_{n, i, i'}$ 为 $\delta$-等维的，维数为 $d - 2$。
（因此从拓扑上看，$Z_{n, i, i'}$ 是坏分支之并去掉它们同良分支的交。）
最后置
$$
Z_n =
\bigcup\nolimits_{\{i, i'\} \in P(I_n)} Z_{n, i, i'} =
\coprod\nolimits_{\{i, i'\} \in P(I_n)} Z_{n, i, i'},
$$
% P12-ERR-0518
，并令 $b_n : X_{n + 1} \to X_n$ 为沿 $Z_n$ 的吹起。注意，在每个集合
$D_{n, i}|_{U_n} \cap D_{n, i'}|_{U_n}$ 的邻域中，引理
\ref{obsolete-lemma-two-divisors} 都适用于态射 $b_n : X_{n + 1} \to X_n$。
因此，完全如证明第一部分那样，对 $\{i, i'\} \in P(I_n)$ 得到有效
Cartier 因子 $D_{n + 1, \{i, i'\}}$，并对 $i \in I_n$
得到有效 Cartier 因子 $D_{n + 1, i}$，使得
$b_n^{-1}(D_{n, i}|_{U_n}) = D_{n + 1, i} +
\sum\nolimits_{i' \in I_n, i' \not = i} D_{n + 1, \{i, i'\}}$.
对每个 $n$，以 $\pi_n : X_n \to X$ 表示复合
$j_0 \circ \ldots \circ j_{n - 1} \circ b_{n - 1}$ 所得的态射。

\medskip\noindent
{\bf 断言：}给定任意拟紧开集 $V \subset X$，则对所有充分大的 $n$，态射
$$
\pi_n^{-1}(V) \leftarrow \pi_{n + 1}^{-1}(V) \leftarrow \ldots
$$
全为同构。事实上，若态射
$\pi_n^{-1}(V) \leftarrow \pi_{n + 1}^{-1}(V)$ 不是同构，则对某个
$\{i, i'\} \in P(I_n)$，有
$Z_{n, i, i'} \cap \pi_n^{-1}(V) \not = \emptyset$。因此存在不可约分支
$W \subset D_{n, i} \cap D_{n, i'}$，满足 $\dim_\delta(W) = d - 1$。
特别地，$\epsilon_W(D_{n, i}, D_{n, i'}) > 0$。反复应用引理
\ref{obsolete-lemma-two-divisors} 可见
$$
\epsilon_W(D_{n, i}, D_{n, i'})
<
\epsilon(V, \{D_i|_V\}) - n
$$
，其中 $\epsilon(V, \{D_i|_V\})$ 如 (\ref{obsolete-equation-invariant}) 所定义。
由于 $V$ 拟紧，有 $\epsilon(V, \{D_i|_V\}) < \infty$；取
$n > \epsilon(V, \{D_i|_V\})$ 即得结论。

\medskip\noindent
注意，差集 $X_n \setminus U_n$ 由构造满足
$\dim_\delta(X_n \setminus U_n) \leq d - 3$。
令 $T_n = \pi_n(X_n \setminus U_n)$ 为其在 $X$ 中的像。
沿上述态射图逐步考察，并利用每个 $b_n$ 都是闭态射，可知对每个 $n$，
$T_0 \cup \ldots \cup T_n$ 都是 $X$ 的闭子集。由构造，任意
$t \in T_n$ 均满足 $\delta(t) \leq d - 3$。因此
$\overline{T_n} \subset X$ 是满足 $\dim_\delta(T_n) \leq d - 3$ 的闭子集。
由上述断言，对任意拟紧开集 $V \subset X$，仅有有限多个 $n$ 满足
$T_n \cap V \not = \emptyset$。故 $\{\overline{T_n}\}_{n \geq 0}$
是一族局部有限的闭子集，可以置
$U = X \setminus \bigcup \overline{T_n}$。这就是引理中的 $U$。

\medskip\noindent
由 $U$ 的构造，$U_n \cap \pi_n^{-1}(U) = \pi_n^{-1}(U)$。因此所有态射
$$
b_n : \pi_{n + 1}^{-1}(U) \longrightarrow \pi_n^{-1}(U)
$$
均为固有态射。此外，由上述断言，它们在 $X$ 的每个拟紧开集上最终成为同构。
因此可以定义
$$
U' = \lim_n \pi_n^{-1}(U).
$$
诱导态射 $b : U' \to U$ 是固有的，因为此性质在 $U$ 上是局部的，
而在每个拟紧开集上该极限都稳定。类似地，利用构造中的嵌入
$I_n \to I_{n + 1}$，置 $J = \bigcup_{n \geq 0} I_n$。
对 $j \in J$，选取 $n_0$ 使 $j$ 对应于某个 $i \in I_{n_0}$，并定义
$D'_j = \lim_{n \geq n_0} D_{n, i}$。这同样有意义，因为态射在 $X$ 上
局部稳定。引理的其余断言与 $X$ 拟紧的情形同样验证。
\end{proof}









\section{相交因子的交换性}
\label{obsolete-section-commutativity}

\noindent
本节结果原本用于给出 Chow Homology 小节
\ref{chow-section-commutativity} 中各引理的另一种证明，
以及 Chow Homology 的引理 \ref{chow-lemma-gysin-commutes-gysin}
的一个弱版本。

\begin{lemma}
\label{obsolete-lemma-improved-additivity}
设 $(S, \delta)$ 如 Chow Homology 的情形 \ref{chow-situation-setup}。
设 $X$ 在 $S$ 上局部有限型，且
$\{i_j : D_j \to X \}_{j \in J}$ 是 $X$ 上一族局部有限的有效
Cartier 因子。对 $j\in J$，设 $n_j > 0$。置
$D = \sum_{j \in J} n_j D_j$，并以 $i : D \to X$ 表示包含态射。
设 $\alpha \in Z_{k + 1}(X)$。则
$$
p : \coprod\nolimits_{j \in J} D_j \longrightarrow D
$$
为固有态射，且
$$
i^*\alpha = p_*\left(\sum n_j i_j^*\alpha\right)
$$
在 $\CH_k(D)$ 中成立。
\end{lemma}

\begin{proof}
由于无穷个有理等价之和涉及一个细微问题，本引理的证明比预期稍长。
在拟紧情形中，族 $D_j$ 是有限的，结论十分容易，直接由 Chow Homology
的引理 \ref{chow-lemma-compute-c1}、Divisors 的引理
\ref{divisors-lemma-c1-additive} 以及各项定义推出。

\medskip\noindent
由于族 $\{D_j\}_{j \in J}$ 局部有限，态射 $p$ 为固有态射。写成
$\alpha = \sum_{a \in A} m_a [W_a]$，其中 $W_a \subset X$
为 $\delta$-维数 $k + 1$ 的整闭子概形。以 $i_a : W_a \to X$
表示闭浸入。我们假设对所有 $a \in A$，有 $m_a \not = 0$，
且 $\{W_a\}_{a \in A}$ 在 $X$ 上局部有限。

\medskip\noindent
注意，由 Chow Homology 的定义 \ref{chow-definition-gysin-homomorphism}，
类 $i^*\alpha$ 是某个循环 $\sum m_a\beta_a$ 的类，
其中 $\beta_a \in Z_k(W_a \cap D)$。具体而言，若
$W_a \not \subset D$，则 $\beta_a = [D \cap W_a]_k$；若
$W_a \subset D$，则 $\beta_a$ 是表示
$c_1(\mathcal{O}_X(D)) \cap [W_a]$ 的循环。

\medskip\noindent
对每个 $a \in A$，写成 $J = J_{a, 1} \amalg J_{a, 2} \amalg J_{a, 3}$，其中
\begin{enumerate}
\item $j \in J_{a, 1}$ 当且仅当 $W_a \cap D_j = \emptyset$；
% P12-ERR-0519
\item $j \in J_{a, 2}$ 当且仅当
$W_a \not = W_a \cap D_1 \not = \emptyset$；
\item $j \in J_{a, 3}$ 当且仅当 $W_a \subset D_j$。
\end{enumerate}
由于族 $\{D_j\}$ 局部有限，$J_{a, 3}$ 是有限集。
对每个 $a \in A$ 与 $j \in J$，按如下方式选取循环
$\beta_{a, j} \in Z_k(W_a \cap D_j)$：
\begin{enumerate}
\item 若 $j \in J_{a, 1}$，置 $\beta_{a, j} = 0$；
\item 若 $j \in J_{a, 2}$，置 $\beta_{a, j} = [D_j \cap W_a]_k$；
\item 若 $j \in J_{a, 3}$，选取 $\beta_{a, j} \in Z_k(W_a)$，
% P12-ERR-0520
使其表示 $c_1(i_a^*\mathcal{O}_X(D_j)) \cap [W_j]$。
\end{enumerate}
我们断言
$$
\beta_a \sim_{rat}
\sum\nolimits_{j \in J} n_j \beta_{a, j}
$$
在 $\CH_k(W_a \cap D)$ 中成立。

\medskip\noindent
情形 I：$W_a \not \subset D$。此时 $J_{a, 3} = \emptyset$。
所以只需证明循环等式
$[D \cap W_a]_k = \sum n_j [D_j \cap W_a]_k$。
这正是引理 \ref{obsolete-lemma-sum-divisors-associated-Weil}。

\medskip\noindent
情形 II：$W_a \subset D$。此时 $\beta_a$ 是表示
$c_1(i_a^*\mathcal{O}_X(D)) \cap [W_a]$ 的循环。写成
$D = D_{a, 1} + D_{a, 2} + D_{a, 3}$，其中
$D_{a, s} = \sum_{j \in J_{a, s}} n_jD_j$。
由 Divisors 的引理 \ref{divisors-lemma-c1-additive}，有
\begin{eqnarray*}
c_1(i_a^*\mathcal{O}_X(D)) \cap [W_a] & = &
c_1(i_a^*\mathcal{O}_X(D_{a, 1})) \cap [W_a] +
c_1(i_a^*\mathcal{O}_X(D_{a, 2})) \cap [W_a] \\
& &
 + c_1(i_a^*\mathcal{O}_X(D_{a, 3})) \cap [W_a].
\end{eqnarray*}
显然，和式第一项为零。由于 $J_{a, 3}$ 有限，最后一项等于
% P12-ERR-0521
$\sum\nolimits_{j \in J_{a, 3}} n_jc_1(i_a^*\mathcal{L}_j) \cap [W_a]$；
见 Divisors 的引理 \ref{divisors-lemma-c1-additive}。
它由 $\sum_{j \in J_{a, 3}} n_j \beta_{a, j}$ 表示。
最后，由情形 I，中间一项由循环
$\sum\nolimits_{j \in J_{a, 2}} n_j[D_j \cap W_a]_k =
\sum_{j \in J_{a, 2}} n_j\beta_{a, j}$ 表示。
因此此情形下断言成立。

\medskip\noindent
现在可以完成引理的证明。由 $\beta_a$ 的选取，有
% P12-ERR-0522
$i^*D \sim_{rat} \sum m_a\beta_a$。对每个 $a$，有
% P12-ERR-0523
$\beta_a \sim_{rat} \sum_j \beta_{a, j}$，该有理等价发生在
$D \cap W_a$ 上。由于各闭子概形 $D \cap W_a$ 在 $D$ 上局部有限，
还可在 $D$ 上得到
$\sum m_a \beta_a \sim_{rat} \sum_{a, j} m_a\beta_{a, j}$！
（见 Chow Homology 的注记
\ref{chow-remark-infinite-sums-rational-equivalences}。）
最后，$\sum_a m_a\beta_{a, j}$（视为 $D_j$ 上的循环）显然表示
$i_j^*\alpha$，故 $\sum_{a, j} m_a\beta_{a, j}$ 表示
$p_* \sum_j i_j^*\alpha$，结论得证。
\end{proof}

\begin{lemma}
\label{obsolete-lemma-commutativity-effective-Cartier-proper-intersection}
设 $(S, \delta)$ 如 Chow Homology 的情形 \ref{chow-situation-setup}。
设 $X$ 在 $S$ 上局部有限型，并假设 $X$ 整且 $\dim_\delta(X) = n$。
设 $D$、$D'$ 为 $X$ 上的有效 Cartier 因子，并假设
$\dim_\delta(D \cap D') = n - 2$。以 $i : D \to X$
（resp.\ $i' : D' \to X$）表示相应的闭浸入。则
\begin{enumerate}
\item 存在循环 $\alpha \in Z_{n - 2}(D \cap D')$，其到 $D$ 的前推表示
$i^*[D']_{n - 1} \in \CH_{n - 2}(D)$，而其到 $D'$ 的前推表示
$(i')^*[D]_{n - 1} \in \CH_{n - 2}(D')$；
\item 有
$$
D \cdot [D']_{n - 1}
=
D' \cdot [D]_{n - 1}
$$
在 $\CH_{n - 2}(X)$ 中成立。
\end{enumerate}
\end{lemma}

\begin{proof}
第 (2) 部分是第 (1) 部分的直接推论。写成
$[D]_{n - 1} = \sum n_a[Z_a]$ 与
$[D']_{n - 1} = \sum m_b[Z_b]$，其中 $Z_a$ 是 $D$ 的不可约分支，
$[Z_b]$ 是 $D'$ 的不可约分支。由 Chow Homology 的定义
\ref{chow-definition-gysin-homomorphism}，有
$i^*D' = \sum m_b i^*[Z_b]$ 以及
$(i')^*D = \sum n_a(i')^*[Z_a]$。由假设，没有不可约分支 $Z_b$
包含于 $D$，故由定义 $i^*[Z_b] = [Z_b\cap D]_{n - 2}$。
同理，$(i')^*[Z_a] = [Z_a \cap D']_{n - 2}$。因此要证明循环等式
$$
\sum n_a[Z_a \cap D']_{n - 2} = \sum m_b[Z_b \cap D]_{n - 2}
$$
；二者确实都支撑在 $D \cap D'$ 上。设 $W \subset X$ 为满足
$\dim_\delta(W) = n - 2$ 的整闭子概形，并设 $\xi \in W$ 为其一般点。
置 $R = \mathcal{O}_{X, \xi}$；这是诺特局部整环。注意
$\dim(R) = 2$。设 $f \in R$（resp.\ $f' \in R$）为定义 $D$
（resp.\ $D'$）之理想的元素。由假设，$\dim(R/(f, f')) = 0$。
设 $\mathfrak q'_1, \ldots, \mathfrak q'_t \subset R$ 为 $(f')$ 上方的
极小素理想，$\mathfrak q_1, \ldots, \mathfrak q_s \subset R$ 为
$(f)$ 上方的极小素理想。上述等式归结为
$$
\sum_{i = 1, \ldots, s}
\text{length}_{R_{\mathfrak q_i}}(R_{\mathfrak q_i}/(f))
\text{ord}_{R/\mathfrak q_i}(f')
=
\sum_{j = 1, \ldots, t}
\text{length}_{R_{\mathfrak q'_j}}(R_{\mathfrak q'_j}/(f'))
\text{ord}_{R/\mathfrak q'_j}(f).
$$
将 Chow Homology 的引理 \ref{chow-lemma-length-multiplication}
应用于 $M = R/(f)$，该等式左边等于
$$
\text{length}_R(R/(f, f'))
-
\text{length}_R(\Ker(f' : R/(f) \to R/(f)))
$$
。注意，$\Ker(f' : R/(f) \to R/(f))$ 通过态射
$x \bmod (f) \mapsto f'x \bmod (ff')$ 典范同构于
$((f) \cap (f'))/(ff')$。所以左边为
$$
\text{length}_R(R/(f, f'))
-
\text{length}_R((f) \cap (f')/(ff'))
$$
。由于该表达式关于 $f$ 与 $f'$ 对称，结论得证。
\end{proof}

\begin{lemma}
\label{obsolete-lemma-commutativity-effective-Cartier-proper-intersection-infinite}
设 $(S, \delta)$ 如 Chow Homology 的情形 \ref{chow-situation-setup}。
设 $X$ 在 $S$ 上局部有限型，并假设 $X$ 整且 $\dim_\delta(X) = n$。
设 $\{D_j\}_{j \in J}$ 为 $X$ 上一族局部有限的有效 Cartier 因子，
$n_j, m_j \geq 0$ 为两族非负整数。置
$D = \sum n_j D_j$ 与 $D' = \sum m_j D_j$。假设对每个
$j \not = j'$ 都有 $\dim_\delta(D_j \cap D_{j'}) = n - 2$。
则在 $\CH_{n - 2}(X)$ 中，
$D \cdot [D']_{n - 1} = D' \cdot [D]_{n - 1}$。
\end{lemma}

\begin{proof}
当各和式有限时，例如 $X$ 拟紧时，本引理是引理
\ref{obsolete-lemma-sum-divisors-associated-Weil} 与
\ref{obsolete-lemma-commutativity-effective-Cartier-proper-intersection}
的直接推论。因此建议读者略过本证明。

\medskip\noindent
下面给出一般情形的证明。设 $i_j : D_j \to X$ 为闭浸入，
以 $p : \coprod D_j \to X$ 表示各态射 $i_j$ 的余积。
设 $\{Z_a\}_{a \in A}$ 为 $\bigcup D_j$ 的不可约分支所成的族。
对每个 $j$，写成
$$
[D_j]_{n - 1} = \sum d_{j, a}[Z_a].
$$
由引理 \ref{obsolete-lemma-sum-divisors-associated-Weil}，有
$$
[D]_{n - 1} = \sum n_j d_{j, a} [Z_a],
\quad
[D']_{n - 1} = \sum m_j d_{j, a} [Z_a].
$$
由引理 \ref{obsolete-lemma-improved-additivity}，有
$$
D \cdot [D']_{n - 1} = p_*\left(\sum n_j i_j^*[D']_{n - 1} \right),
\quad
D' \cdot [D]_{n - 1} = p_*\left(\sum m_{j'} i_{j'}^*[D]_{n - 1} \right).
$$
如 Gysin 同态的定义（见 Chow Homology 的定义
\ref{chow-definition-gysin-homomorphism}）那样，选取
$D_j \cap Z_a$ 上表示 $i_j^*[Z_a]$ 的循环 $\beta_{a, j}$。
（注意，若 $Z_a$ 不包含于 $D_j$，则事实上
$\beta_{a, j} = [D_j \cap Z_a]_{n - 2}$，即此时并无选择。）
由于 $p$ 限制在每个 $D_j$ 上都是闭浸入，可以而且将把
$\beta_{a, j}$ 视为 $X$ 上的循环。代入上面所得
$[D]_{n - 1}$ 与 $[D']_{n - 1}$ 的公式，得到
$$
D \cdot [D']_{n - 1} =
\sum\nolimits_{j, j', a} n_j m_{j'} d_{j', a} \beta_{a, j},
\quad
D' \cdot [D]_{n - 1} =
\sum\nolimits_{j, j', a} m_{j'} n_j d_{j, a} \beta_{a, j'}.
$$
此外，在同样约定下，还有
$$
D_j \cdot [D_{j'}]_{n - 1} = \sum d_{j', a} \beta_{a, j}.
$$
用这些记号，引理
\ref{obsolete-lemma-commutativity-effective-Cartier-proper-intersection}
（亦见其证明）表明：当 $j \not = j'$ 时，循环
$\sum d_{j', a} \beta_{a, j}$ 与 $\sum d_{j, a} \beta_{a, j'}$
作为循环相等！因此
\begin{eqnarray*}
D \cdot [D']_{n - 1}
& = &
\sum\nolimits_{j, j', a} n_j m_{j'} d_{j', a} \beta_{a, j} \\
& = &
\sum\nolimits_{j \not = j'} n_j m_{j'}
\left(\sum\nolimits_a d_{j', a} \beta_{a, j}\right) +
\sum\nolimits_{j, a} n_j m_j d_{j, a} \beta_{a, j} \\
& = &
\sum\nolimits_{j \not = j'} n_j m_{j'}
\left(\sum\nolimits_a d_{j, a} \beta_{a, j'}\right) +
\sum\nolimits_{j, a} n_j m_j d_{j, a} \beta_{a, j} \\
& = &
\sum\nolimits_{j, j', a} m_{j'} n_j d_{j, a} \beta_{a, j'} \\
& = &
D' \cdot [D]_{n - 1}
\end{eqnarray*}
，结论得证。
\end{proof}

\begin{lemma}
\label{obsolete-lemma-commutativity-effective-Cartier}
设 $(S, \delta)$ 如 Chow Homology 的情形 \ref{chow-situation-setup}。
设 $X$ 在 $S$ 上局部有限型，并假设 $X$ 整且 $\dim_\delta(X) = n$。
设 $D$、$D'$ 为 $X$ 上的有效 Cartier 因子。则
$$
D \cdot [D']_{n - 1} = D' \cdot [D]_{n - 1}
$$
在 $\CH_{n - 2}(X)$ 中成立。
\end{lemma}

\begin{proof}[引理 \ref{obsolete-lemma-commutativity-effective-Cartier} 的第一种证明]
先证明 $X$ 拟紧的情形。此时，对 $X$ 和由有效 Cartier 因子构成的
二元集合 $\{D, D'\}$ 应用引理 \ref{obsolete-lemma-blowing-up-intersections}。
于是得到固有态射 $b : X' \to X$，以及一族有限的有效 Cartier 因子
$D'_j \subset X'$；它们两两以余维 $\geq 2$ 相交，并且
$b^{-1}(D) = \sum n_j D'_j$、$b^{-1}(D') = \sum m_j D'_j$。
注意，在 $Z_{n - 1}(X)$ 中
$b_*[b^{-1}(D)]_{n - 1} = [D]_{n - 1}$，对 $D'$ 亦然；
见引理 \ref{obsolete-lemma-push-pull-effective-Cartier}。因此由 Chow Homology
的引理 \ref{chow-lemma-pushforward-cap-c1}，有
$$
D \cdot [D']_{n - 1} = b_*\left(b^{-1}(D) \cdot [b^{-1}(D')]_{n - 1}\right)
$$
在 $\CH_{n - 2}(X)$ 中成立；另一项亦然。因此本引理由等式
$b^{-1}(D) \cdot [b^{-1}(D')]_{n - 1} = b^{-1}(D') \cdot [b^{-1}(D)]_{n - 1}$
推出；该等式由引理
\ref{obsolete-lemma-commutativity-effective-Cartier-proper-intersection-infinite}
在 $\CH_{n - 2}(X')$ 中给出。

\medskip\noindent
注意，上述证明中引用的每个引理在一般情形下也成立（此时不假设
$X$ 拟紧）。一般情形唯一的细小变化是：应用引理
\ref{obsolete-lemma-blowing-up-intersections} 得到的态射
$b : U' \to U$，其目标是开集 $U \subset X$，且该开集的补集余维
$\geq 3$。因此由 Chow Homology 的引理
\ref{chow-lemma-restrict-to-open}，有
$\CH_{n - 2}(U) = \CH_{n - 2}(X)$；以 $U$ 替换 $X$ 后，
证明余下部分原封不动地成立。
\end{proof}

\begin{proof}[引理 \ref{obsolete-lemma-commutativity-effective-Cartier} 的第二种证明]
设 $\mathcal{I} = \mathcal{O}_X(-D)$ 与
$\mathcal{I}' = \mathcal{O}_X(-D')$ 分别为 $D$ 与 $D'$ 的可逆理想层。
记
$\mathcal{I}_{D'} = \mathcal{I} \otimes_{\mathcal{O}_X} \mathcal{O}_{D'}$
以及
$\mathcal{I}'_D = \mathcal{I}' \otimes_{\mathcal{O}_X} \mathcal{O}_D$。
将包含态射 $\mathcal{I} \to \mathcal{O}_X$ 限制到 $D'$，得到态射
$$
\varphi : \mathcal{I}_{D'} \longrightarrow \mathcal{O}_{D'}
$$
；类似地得到
$$
\psi : \mathcal{I}'_D \longrightarrow \mathcal{O}_D
$$
。显然
$$
\Coker(\varphi)
\cong
\mathcal{O}_{D \cap D'}
\cong
\Coker(\psi)
$$
，并且
$$
\Ker(\varphi)
\cong
\frac{\mathcal{I} \cap \mathcal{I}'}{\mathcal{I}\mathcal{I}'}
\cong
\Ker(\psi).
$$
因此可见
$$
\gamma =
[\mathcal{I}_{D'}] - [\mathcal{O}_{D'}]
=
[\mathcal{I}'_D] - [\mathcal{O}_D]
$$
在 $K_0(\textit{Coh}_{\leq n - 1}(X))$ 中成立。另一方面，显然
$$
[\mathcal{I}'_D]_{n - 1} = [D]_{n - 1}, \quad
[\mathcal{I}_{D'}]_{n - 1} = [D']_{n - 1}.
$$
并且
$$
\mathcal{O}_X(D') \otimes \mathcal{I}'_D = \mathcal{O}_D, \quad
\mathcal{O}_X(D) \otimes \mathcal{I}_{D'} = \mathcal{O}_{D'}.
$$
将 Chow Homology 的引理 \ref{chow-lemma-coherent-sheaf-cap-c1}
应用两次，这意味着元素 $\gamma$ 属于 $B_{n - 2}(X)$，并同时映到
$c_1(\mathcal{O}_X(D')) \cap [D]_{n - 1}$ 与
$c_1(\mathcal{O}_X(D)) \cap [D']_{n - 1}$，结论得证
（因为态射 $B_{n - 2}(X) \to \CH_{n - 2}(X)$ 是良定义的——
这正是本证明的关键）。
\end{proof}











\section{正则固有模型上的对偶化模}
\label{obsolete-section-dualizing}

\noindent
在 Semistable Reduction 的情形 \ref{models-situation-regular-model} 中，
令 $\omega_{X/R}^\bullet = f^!\mathcal{O}_{\Spec(R)}$ 为态射
$f : X \to \Spec(R)$ 的相对对偶化复形；见 Duality for Schemes 的注记
\ref{duality-remark-relative-dualizing-complex}。由 Semistable Reduction
的引理 \ref{models-lemma-gorenstein}，$f$ 是相对维数 $1$ 的戈伦斯坦态射，
故可以使用 Duality for Schemes 的引理
\ref{duality-lemma-affine-flat-Noetherian-gorenstein},
\ref{duality-lemma-CM-shriek} 与
\ref{duality-lemma-gorenstein-CM-morphism}
得到
$$
\omega_{X/R}^\bullet = \omega_X[1]
$$
，其中 $\omega_X$ 是某个可逆 $\mathcal{O}_X$-模。这个可逆模通常称为
$X$ 在 $R$ 上的{\it 相对对偶化模}。
由于 $R$ 正则（因而戈伦斯坦）且维数为 $1$，
$\omega_R^\bullet = R[1]$ 是 $R$ 的归一化对偶化复形。所以
$\omega_X = H^{-2}(f^!\omega_R^\bullet)$；由此可见，$\omega_X$
不仅是相对对偶化模，也是对偶化模，见 Duality for Schemes 的例
\ref{duality-example-proper-over-local}。因此 $\omega_X$ 表示函子
$$
\textit{Coh}(\mathcal{O}_X) \to \textit{Sets},\quad
\mathcal{F} \mapsto \Hom_R(H^1(X, \mathcal{F}), R)
$$
，这是 Duality for Schemes 的引理
\ref{duality-lemma-dualizing-module-proper-over-A}。这给出了
Semistable Reduction 的情形 \ref{models-situation-regular-model}
中相对对偶化模的另一种定义。$\omega_X$ 的形成与任意基变换可换
（对任意给定相对维数的固有戈伦斯坦态射均如此）；这由相对对偶化复形
的相应事实推出，该事实见 Duality for Schemes 的注记
\ref{duality-remark-relative-dualizing-complex}，其源自 Duality for Schemes
的引理 \ref{duality-lemma-proper-flat-base-change}。因此 $\omega_X$
拉回为 Algebraic Curves 的引理 \ref{curves-lemma-duality-dim-1-CM}
中讨论的 $C$ 在 $K$ 上的对偶化模 $\omega_C$。注意，由 Algebraic Curves
的引理 \ref{curves-lemma-duality-dim-1}，$\omega_C$ 同构于
$\Omega_{C/K}$。类似地，$\omega_X|_{X_k}$ 是 $X_k$ 在 $k$ 上的
对偶化模 $\omega_{X_k}$。

\begin{lemma}
\label{obsolete-lemma-dualizing-components}
在 Semistable Reduction 的情形 \ref{models-situation-regular-model} 中，
$C_i$ 在 $k$ 上的对偶化模为
$$
\omega_{C_i} = \omega_X(C_i)|_{C_i}
$$
，其中 $\omega_X$ 如上。
\end{lemma}

\begin{proof}
设 $t : C_i \to X$ 为闭浸入。由于 $t$ 是有效 Cartier 因子的包含，
由 Duality for Schemes 的引理
\ref{duality-lemma-twisted-inverse-image-closed} 与
\ref{duality-lemma-sheaf-with-exact-support-effective-Cartier} 可知，
对每个可逆 $\mathcal{O}_X$-模 $\mathcal{L}$，都有
$t^!(\mathcal{L}) = \mathcal{L}(C_i)|_{C_i}$。考虑交换图
$$
\xymatrix{
C_i \ar[r]_t \ar[d]_g & X \ar[d]^f \\
\Spec(k) \ar[r]^s & \Spec(R)
}
$$
。注意，$C_i$ 是戈伦斯坦曲线（Semistable Reduction 的引理
\ref{models-lemma-gorenstein}），其可逆对偶化模 $\omega_{C_i}$
由性质 $\omega_{C_i}[0] = g^!\mathcal{O}_{\Spec(k)}$ 刻画。
参见 Algebraic Curves 的引理 \ref{curves-lemma-duality-dim-1} 及其证明，
以及 Algebraic Curves 的引理 \ref{curves-lemma-duality-dim-1-CM}
与 \ref{curves-lemma-rr}。另一方面，$s^!(R[1]) = k$，故
$$
\omega_{C_i}[0] =
g^! s^!(R[1]) = t^!f^!(R[1]) = t^!\omega_X
$$
。结合以上各式即得引理结论。
\end{proof}






\section{重复项与拆分出的引用}
\label{obsolete-section-duplicates}

\noindent
本节收集重复项，以及那些过去同时陈述若干件事、如今已拆分为各组成部分的引用。

\begin{lemma}
\label{obsolete-lemma-characterize-epi-mono}
设 $X$ 为拓扑空间。
\begin{enumerate}
\item 预层范畴中的满态射（resp.\ 单态射）恰好是预层的满射
（resp.\ 单射）。
\item 层范畴中的满态射（resp.\ 单态射）恰好是层的满射
（resp.\ 单射），也恰好是在所有茎上均为满射（resp.\ 单射）的态射。
\item 集合预层的满射（resp.\ 单射）态射的层化仍为满射（resp.\ 单射）。
\end{enumerate}
\end{lemma}

\begin{proof}
合并 Sheaves 的引理 \ref{sheaves-lemma-characterize-epi-mono}
与 \ref{sheaves-lemma-sheafify-epi-mono} 即得。
\end{proof}

\begin{lemma}
\label{obsolete-lemma-directed-colimit-finite-type}
设 $X$ 为概形，并假设 $X$ 拟紧且拟分离。
设 $\mathcal{F}$ 为拟凝聚 $\mathcal{O}_X$-模。
则 $\mathcal{F}$ 是其有限型拟凝聚子模的滤过余极限。
\end{lemma}

\begin{proof}
这是 Properties 的引理
\ref{properties-lemma-quasi-coherent-colimit-finite-type} 的重复项。
\end{proof}

\begin{lemma}
\label{obsolete-lemma-points-monomorphism}
设 $S$ 为概形，$X$ 为 $S$ 上的代数空间。态射
$\{\Spec(k) \to X \text{ monomorphism}\} \to |X|$ 为单射。
\end{lemma}

\begin{proof}
这是 Properties of Spaces 的引理
\ref{spaces-properties-lemma-points-monomorphism} 的重复项。
\end{proof}

\begin{theorem}
\label{obsolete-theorem-equivalence-sheaves-point}
设 $S = \Spec(K)$，其中 $K$ 为域。设 $\overline{s}$ 为 $S$ 的几何点，
$G = \text{Gal}_{\kappa(s)}$ 表示绝对伽罗瓦群。则存在范畴等价
$\Sh(S_\etale) \to G\textit{-Sets}$,
$\mathcal{F} \mapsto \mathcal{F}_{\overline{s}}$.
\end{theorem}

\begin{proof}
这是 \'Etale Cohomology 的定理
\ref{etale-cohomology-theorem-equivalence-sheaves-point} 的重复项。
\end{proof}

\begin{remark}
\label{obsolete-remark-tangent-spaces}
你来到这里是因为标签重复。实际内容请见 Formal Deformation Theory
的小节 \ref{formal-defos-section-tangent-spaces}。
\end{remark}

\begin{lemma}
\label{obsolete-lemma-locally-ringed-space-direct-summand-free}
设 $X$ 为局部环化空间。有限自由 $\mathcal{O}_X$-模的直和因子有限局部自由。
\end{lemma}

\begin{proof}
这是 Modules 的引理
\ref{modules-lemma-direct-summand-of-locally-free-is-locally-free} 的重复项。
\end{proof}

\begin{lemma}
\label{obsolete-lemma-characterize-injective}
设 $R$ 为环，$E$ 为 $R$-模。下列条件等价：
\begin{enumerate}
\item $E$ 是内射 $R$-模；
\item 给定理想 $I \subset R$ 和模态射 $\varphi : I \to E$，
存在 $\varphi$ 到 $R$-模态射 $R \to E$ 的延拓。
\end{enumerate}
\end{lemma}

\begin{proof}
这是 Baer 判据；见 Injectives 的引理 \ref{injectives-lemma-criterion-baer}。
\end{proof}

\begin{lemma}
\label{obsolete-lemma-periodic-length}
设 $R$ 为局部环。
\begin{enumerate}
\item 若 $(M, N, \varphi, \psi)$ 是 $2$-周期复形，且 $M$、$N$ 长度有限，
则 $e_R(M, N, \varphi, \psi) = \text{length}_R(M) - \text{length}_R(N)$。
\item 若 $(M, \varphi, \psi)$ 是 $(2, 1)$-周期复形，且 $M$ 长度有限，
则 $e_R(M, \varphi, \psi) = 0$。
\item 假设有 $2$-周期复形的短正合列
$$
0 \to (M_1, N_1, \varphi_1, \psi_1)
\to (M_2, N_2, \varphi_2, \psi_2)
\to (M_3, N_3, \varphi_3, \psi_3)
\to 0
$$
。若三者中有两个的上同调模长度有限，则第三个也如此，并且有
$$
e_R(M_2, N_2, \varphi_2, \psi_2) =
e_R(M_1, N_1, \varphi_1, \psi_1) +
e_R(M_3, N_3, \varphi_3, \psi_3).
$$
\end{enumerate}
\end{lemma}

\begin{proof}
这由 Chow Homology 的引理 \ref{chow-lemma-additivity-periodic-length}
与 \ref{chow-lemma-finite-periodic-length} 推出。
\end{proof}


\begin{lemma}
\label{obsolete-lemma-extensions-of-rings}
设 $A$ 为环，$I$ 为 $A$-模。
\begin{enumerate}
\item 以 $I$ 为平方零理想的环扩张
$0 \to I \to A' \to A \to 0$ 所成集合，典范双射于
$\Ext^1_A(\NL_{A/\mathbf{Z}}, I)$。
\item 给定环同态 $A \to B$、$B$-模 $N$、$A$-模态射
$c : I \to N$，并给定下列核平方为零的环扩张：
\begin{enumerate}
\item[(a)] $0 \to I \to A' \to A \to 0$，对应于
$\alpha \in \Ext^1_A(\NL_{A/\mathbf{Z}}, I)$；
\item[(b)] $0 \to N \to B' \to B \to 0$，对应于
$\beta \in \Ext^1_B(\NL_{B/\mathbf{Z}}, N)$；
\end{enumerate}
则存在符合 Deformation Theory 的式 (\ref{defos-equation-to-solve})
的态射 $A' \to B'$，当且仅当 $\beta$ 与 $\alpha$ 映到
$\Ext^1_A(\NL_{A/\mathbf{Z}}, N)$ 中的同一元素。
\end{enumerate}
\end{lemma}

\begin{proof}
这由 Deformation Theory 的引理
\ref{defos-lemma-extensions-of-algebras} 与
\ref{defos-lemma-extensions-of-algebras-functorial} 推出。
\end{proof}

\begin{lemma}
\label{obsolete-lemma-extensions-of-ringed-spaces}
设 $(S, \mathcal{O}_S)$ 为环化空间，$\mathcal{J}$ 为 $\mathcal{O}_S$-模。
\begin{enumerate}
\item 以 $\mathcal{J}$ 为平方零理想的环层扩张
$0 \to \mathcal{J} \to \mathcal{O}_{S'} \to \mathcal{O}_S \to 0$
所成集合，典范双射于
$\Ext^1_{\mathcal{O}_S}(\NL_{S/\mathbf{Z}}, \mathcal{J})$。
\item 给定环化空间态射
$f : (X, \mathcal{O}_X) \to (S, \mathcal{O}_S)$、$\mathcal{O}_X$-模
$\mathcal{G}$、$f$-态射 $c : \mathcal{J} \to \mathcal{G}$，
并给定下列核平方为零的环层扩张：
\begin{enumerate}
\item[(a)] $0 \to \mathcal{J} \to \mathcal{O}_{S'} \to \mathcal{O}_S \to 0$
，对应于
$\alpha \in \Ext^1_{\mathcal{O}_S}(\NL_{S/\mathbf{Z}}, \mathcal{J})$；
\item[(b)] $0 \to \mathcal{G} \to \mathcal{O}_{X'} \to \mathcal{O}_X \to 0$
，对应于
$\beta \in \Ext^1_{\mathcal{O}_X}(\NL_{X/\mathbf{Z}}, \mathcal{G})$；
\end{enumerate}
则存在符合 Deformation Theory 的式
(\ref{defos-equation-to-solve-ringed-spaces}) 的态射 $X' \to S'$，
当且仅当 $\beta$ 与 $\alpha$ 映到
$\Ext^1_{\mathcal{O}_X}(Lf^*\NL_{S/\mathbf{Z}}, \mathcal{G})$
中的同一元素。
\end{enumerate}
\end{lemma}

\begin{proof}
这由 Deformation Theory 的引理
\ref{defos-lemma-extensions-of-relative-ringed-spaces} 与
\ref{defos-lemma-extensions-of-relative-ringed-spaces-functorial} 推出。
\end{proof}

\begin{lemma}
\label{obsolete-lemma-extensions-of-ringed-topoi}
设 $(\Sh(\mathcal{B}), \mathcal{O}_\mathcal{B})$ 为环化拓扑斯，
$\mathcal{J}$ 为 $\mathcal{O}_\mathcal{B}$-模。
\begin{enumerate}
\item 以 $\mathcal{J}$ 为平方零理想的环层扩张
$0 \to \mathcal{J} \to \mathcal{O}_{\mathcal{B}'} \to
\mathcal{O}_\mathcal{B} \to 0$
所成集合，典范双射于
$\Ext^1_{\mathcal{O}_\mathcal{B}}(
\NL_{\mathcal{O}_\mathcal{B}/\mathbf{Z}}, \mathcal{J})$.
\item 给定环化拓扑斯态射
$f : (\Sh(\mathcal{C}), \mathcal{O}) \to
(\Sh(\mathcal{B}), \mathcal{O}_\mathcal{B})$、$\mathcal{O}$-模
$\mathcal{G}$、$f^{-1}\mathcal{O}_\mathcal{B}$-模态射
$c : f^{-1}\mathcal{J} \to \mathcal{G}$，并给定下列核平方为零的环层扩张：
\begin{enumerate}
\item[(a)] $0 \to \mathcal{J} \to \mathcal{O}_{\mathcal{B}'} \to
\mathcal{O}_\mathcal{B} \to 0$，对应于
$\alpha \in \Ext^1_{\mathcal{O}_\mathcal{B}}(
\NL_{\mathcal{O}_\mathcal{B}/\mathbf{Z}}, \mathcal{J})$；
\item[(b)] $0 \to \mathcal{G} \to \mathcal{O}' \to \mathcal{O} \to 0$
，对应于
$\beta \in \Ext^1_\mathcal{O}(\NL_{\mathcal{O}/\mathbf{Z}}, \mathcal{G})$
\end{enumerate}
则存在符合 Deformation Theory 的式
(\ref{defos-equation-to-solve-ringed-topoi}) 的态射
$(\Sh(\mathcal{C}), \mathcal{O}') \to
(\Sh(\mathcal{B}, \mathcal{O}_{\mathcal{B}'})$，当且仅当
$\beta$ 与 $\alpha$ 映到下列 Ext 群中的同一元素：
$\Ext^1_\mathcal{O}(
Lf^*\NL_{\mathcal{O}_\mathcal{B}/\mathbf{Z}}, \mathcal{G})$.
\end{enumerate}
\end{lemma}

\begin{proof}
这由 Deformation Theory 的引理
\ref{defos-lemma-extensions-of-relative-ringed-topoi} 与
\ref{defos-lemma-extensions-of-relative-ringed-topoi-functorial} 推出。
\end{proof}

\begin{remark}
\label{obsolete-remark-examples-formal-defos}
这个标签过去指向一个描述若干形变问题例子的章节。现在每个例子都有自己的章节。
见 Deformation Problems 的小节
\ref{examples-defos-section-finite-projective-modules},
\ref{examples-defos-section-representations},
\ref{examples-defos-section-continuous-representations} 与
\ref{examples-defos-section-graded-algebras}.
\end{remark}

\begin{lemma}
\label{obsolete-lemma-examples-have-RS}
Deformation Problems 的例
\ref{examples-defos-example-finite-projective-modules},
\ref{examples-defos-example-representations},
\ref{examples-defos-example-continuous-representations} 与
\ref{examples-defos-example-graded-algebras}
满足 Rim--Schlessinger 条件 (RS)。
\end{lemma}

\begin{proof}
这由 Deformation Problems 的引理
\ref{examples-defos-lemma-finite-projective-modules-RS},
\ref{examples-defos-lemma-representations-RS},
\ref{examples-defos-lemma-continuous-representations-RS} 与
\ref{examples-defos-lemma-graded-algebras-RS} 推出。
\end{proof}

\begin{lemma}
\label{obsolete-lemma-tangent-and-inf}
有下列典范 $k$-向量空间同一化：
\begin{enumerate}
\item 在 Deformation Problems 的例
\ref{examples-defos-example-finite-projective-modules}
中，若 $x_0 = (k, V)$，则 $T_{x_0}\mathcal{F} = (0)$ 与
$\text{Inf}_{x_0}(\mathcal{F}) = \text{End}_k(V)$ 均为有限维。
\item 在 Deformation Problems 的例
\ref{examples-defos-example-representations}
中，若 $x_0 = (k, V, \rho_0)$，则
$T_{x_0}\mathcal{F} = \Ext^1_{k[\Gamma]}(V, V) = H^1(\Gamma, \text{End}_k(V))$
与 $\text{Inf}_{x_0}(\mathcal{F}) = H^0(\Gamma, \text{End}_k(V))$
在 $\Gamma$ 有限生成时均为有限维。
\item 在 Deformation Problems 的例
\ref{examples-defos-example-continuous-representations}
中，若 $x_0 = (k, V, \rho_0)$，则
$T_{x_0}\mathcal{F} = H^1_{cont}(\Gamma, \text{End}_k(V))$
且
$\text{Inf}_{x_0}(\mathcal{F}) = H^0_{cont}(\Gamma, \text{End}_k(V))$
在 $\Gamma$ 拓扑有限生成时均为有限维。
\item 在 Deformation Problems 的例
\ref{examples-defos-example-graded-algebras}
中，若 $x_0 = (k, P)$，则
$T_{x_0}\mathcal{F}$ 与 $\text{Inf}_{x_0}(\mathcal{F}) = \text{Der}_k(P, P)$
在 $P$ 为 $k$ 上有限生成时均为有限维。
\end{enumerate}
\end{lemma}

\begin{proof}
这由 Deformation Problems 的引理
\ref{examples-defos-lemma-finite-projective-modules-TI},
\ref{examples-defos-lemma-representations-TI},
\ref{examples-defos-lemma-continuous-representations-TI} 与
\ref{examples-defos-lemma-graded-algebras-TI} 推出。
\end{proof}
